% Les composés oxygénés : alcools, aldéhydes, cétones — Chimie 1ère D — Cours signé OnBuch+
% Programme officiel MINESEC, Première C, D, E, TI. Compiler avec : tectonic N03-les-composes-oxygenes.tex
\newcommand{\DOCMATIERE}{Chimie}
\newcommand{\DOCNIVEAU}{1ère D}
\newcommand{\DOCMODULE}{Module 1 : Chimie organique}
\newcommand{\DOCLECON}{03}
\newcommand{\DOCTITRE}{Les composés oxygénés : alcools, aldéhydes, cétones}
% OnBuch+ — préambule « cours signés » ENRICHI (compilé avec tectonic / XeLaTeX)
% Système visuel « Pop » d'OnBuch+ : fond crème (jamais blanc), bordures encre
% épaisses, ombres « dures » décalées, titres Archivo Black en capitales.
% Version enrichie pour les cours de chimie : chimie (mhchem, chemfig),
% graphes (pgfplots), boîtes pédagogiques supplémentaires (définition,
% propriété, méthode, attention, le savais-tu, exercice résolu, exercice,
% TP, bilan), page de garde refaite et signature OnBuch+ ancrée en bas.
%
% Macros attendues dans main.tex AVANT \input{preamble} :
%   \DOCMATIERE  \DOCNIVEAU  \DOCMODULE  \DOCLECON (numéro)  \DOCTITRE
\documentclass[11pt]{article}

\usepackage{fontspec}
\usepackage[french]{babel}
\usepackage[a4paper,margin=2cm,top=3.4cm,bottom=2.8cm,headheight=1.4cm,footskip=1.3cm]{geometry}
\usepackage{amsmath,amssymb,amsfonts}
\usepackage[table]{xcolor} % [table] : fournit aussi \rowcolors (alternance de lignes)
\usepackage{pagecolor}
\usepackage{tikz}
\usetikzlibrary{arrows.meta,positioning,calc,shapes.geometric,shapes.misc,shapes.arrows,decorations.pathmorphing,decorations.markings,patterns,shadows,mindmap,trees,fit,backgrounds,matrix,intersections}
\usepackage{pgfplots}
\pgfplotsset{compat=1.18}
\usepackage[version=4]{mhchem}
\usepackage{chemfig}
\usepackage{enumitem}
\usepackage{fancyhdr}
% [most] charge toutes les bibliothèques tcolorbox (listings, minted, raster,
% documentation...) et gonfle inutilement la mémoire TeX sur un document long
% (~300 boîtes) : on ne charge que ce qui est réellement utilisé.
\usepackage[skins,breakable]{tcolorbox}
\usepackage{graphicx}
\usepackage{colortbl}
\usepackage{booktabs}
\usepackage{tabularx}
\usepackage{array}
\usepackage{multicol}
\usepackage{siunitx}
% parse-numbers=false : accepte aussi \SI{2,5 \times 10^{-3}}{...}
\sisetup{locale=FR,output-decimal-marker={,},group-minimum-digits=4,parse-numbers=false}
\DeclareSIUnit\ohm{\ensuremath{\symup{\Omega}}} % U+2126 absent de la police
\usepackage{qrcode}
% \degree : commande fréquemment utilisée par les modèles pour un angle ou une
% température en dehors de \SI{}{} (non fournie par les paquets chargés).
\providecommand{\degree}{^{\circ}}
\usepackage{lastpage}
\usepackage{hyperref}
\hypersetup{hidelinks}

% ============================================================
% Palette « Pop » OnBuch+ — mêmes hex que la classe P (plus_ui.dart)
% ============================================================
\definecolor{poporange}{HTML}{F59321}
\definecolor{popdark}{HTML}{E07A0C}
\definecolor{popcream}{HTML}{FFF6E8}
\definecolor{popink}{HTML}{1C1714}
\definecolor{poppurple}{HTML}{7A5AE0}
\definecolor{popgreen}{HTML}{2E9E6A}
\definecolor{poppink}{HTML}{E0457B}
\definecolor{popblue}{HTML}{2F7DE1}
\definecolor{popgold}{HTML}{C98A12}
\definecolor{popmuted}{HTML}{8A7F76}
\definecolor{popline}{HTML}{EADFD2}
\definecolor{popgray}{HTML}{8A7F76}
\colorlet{popgrey}{popgray}
% Alias tolérants (le modèle nomme parfois les couleurs différemment)
\colorlet{popwhite}{white}
\colorlet{popblack}{popink}
\colorlet{popred}{poppink}
\colorlet{popyellow}{popgold}
% Teintes claires pour fonds de tableaux / zones
\colorlet{popblueL}{popblue!10!white}
\colorlet{poporangeL}{poporange!14!white}
\colorlet{popgreenL}{popgreen!12!white}
\colorlet{poppurpleL}{poppurple!12!white}
\colorlet{poppinkL}{poppink!12!white}
\colorlet{popgoldL}{popgold!14!white}

\pagecolor{popcream}

% ============================================================
% Typographie
% ============================================================
\defaultfontfeatures{Ligatures=TeX}
\setmainfont{PlusJakartaSans-Medium.ttf}[Path=../../../fonts/,BoldFont=PlusJakartaSans-Bold.ttf]
% Maths en sans-serif (Fira Math) avec chiffres et lettres droites de Plus
% Jakarta, pour que les formules se fondent dans le texte.
\usepackage[math-style=ISO,bold-style=ISO]{unicode-math}
\setmathfont{FiraMath-Regular.otf}[Path=../../../fonts/]
\setmathfont{PlusJakartaSans-Medium.ttf}[Path=../../../fonts/,range={up/num,up/{Latin,latin},\mathup}]
\usepackage{icomma} % virgule décimale sans espace en mode math
% Caractères absents de la police de texte (sinon carré vide) : rendus par
% leur équivalent mathématique, en texte comme en formule.
\usepackage{newunicodechar}
\newunicodechar{α}{\ensuremath{\alpha}}
\newunicodechar{Α}{\ensuremath{\alpha}}
\newunicodechar{β}{\ensuremath{\beta}}
\newunicodechar{δ}{\ensuremath{\delta}}
\newunicodechar{✓}{\popcheck{popgreen}}
\newfontfamily\popdisplay{ArchivoBlack-Regular.ttf}[Path=../../../fonts/]
\newfontfamily\poplogo{SpaceGrotesk-Bold.ttf}[Path=../../../fonts/]
\newfontfamily\popsemi{PlusJakartaSans-SemiBold.ttf}[Path=../../../fonts/]
\newfontfamily\popxbold{PlusJakartaSans-ExtraBold.ttf}[Path=../../../fonts/]
\newfontfamily\popxboldls{PlusJakartaSans-ExtraBold.ttf}[Path=../../../fonts/,LetterSpace=3.0]

\setlist[enumerate]{leftmargin=1.7em,itemsep=5pt,topsep=4pt}
\setlist[itemize]{leftmargin=1.7em,itemsep=4pt,topsep=4pt,label={\color{poporange}\textbullet}}
\setlength{\parindent}{0pt}
\setlength{\parskip}{5pt}
\renewcommand{\arraystretch}{1.25}

% chemfig : liaisons un peu plus épaisses, encre Pop
\setchemfig{atom sep=2em,bond style={line width=0.9pt,popink},cram width=3pt}

% pgfplots : style Pop par défaut
\pgfplotsset{
  every axis/.append style={axis line style={popink,line width=1pt},tick style={popink},
    grid style={popline,line width=0.5pt},label style={font=\small},tick label style={font=\footnotesize}},
  popcurve/.style={poporange,line width=1.8pt,no markers,smooth},
}

% ============================================================
% Pill / squiggle
% ============================================================
\newcommand{\poppill}[3]{%
  \tikz[baseline=-0.55ex]\node[draw=popink,line width=1.2pt,rounded corners=2.6pt,%
    fill=#2,inner xsep=6.5pt,inner ysep=3pt]{%
    \popxboldls\fontsize{7}{7}\selectfont\color{#3}\MakeUppercase{#1}};%
}
\newcommand{\popsquiggle}[1]{%
  \tikz[baseline=-0.4ex]{
    \draw[#1,line width=2.6pt,line cap=round]
      plot [smooth] coordinates {(0,0.09) (0.34,0.01) (0.68,0.21) (1.02,0.01) (1.36,0.10)};
  }%
}
% Mot-clé surligné (marqueur orange sous le texte)
\newcommand{\cle}[1]{{\popxbold\color{popdark}#1}}

% ============================================================
% En-tête / pied
% ============================================================
\pagestyle{fancy}
\fancyhf{}
\renewcommand{\headrulewidth}{0pt}
\renewcommand{\footrulewidth}{0pt}
\lhead{\raisebox{-0.15ex}{\poplogo\fontsize{13}{13}\selectfont\color{popink}OnBuch\color{poporange}+}}
\rhead{\poppill{\DOCMATIERE\ \textbullet\ \DOCNIVEAU\ \textbullet\ Leçon \DOCLECON}{white}{popink}}
\lfoot{\popxbold\fontsize{7.6}{7.6}\selectfont\color{popmuted}\MakeUppercase{Cours signé OnBuch+}\ \textbullet\ {\popsemi\DOCTITRE}}
\rfoot{\tikz[baseline=-0.6ex]\node[rounded corners=8.5pt,fill=popink,minimum height=17pt,minimum width=17pt,inner xsep=5pt,inner sep=0]{\popxbold\fontsize{8}{8}\selectfont\color{popcream}\thepage\,/\,\pageref*{LastPage}};}

% ============================================================
% Titres
% ============================================================
\newcommand{\chaptitre}[2]{%
  \begin{tcolorbox}[enhanced,colback=poporange,colframe=popink,boxrule=2.6pt,arc=11pt,
      drop shadow=popink,left=15pt,right=15pt,top=12pt,bottom=11pt,before skip=3pt,after skip=18pt]
    \poppill{#2}{popink}{popcream}\par\vspace{9pt}
    {\raggedright\hyphenpenalty=10000\exhyphenpenalty=10000\popdisplay\fontsize{22}{24}\selectfont\color{popcream}\MakeUppercase{#1}\par}
  \end{tcolorbox}
}
% Section numérotée : pastille orange avec le numéro + titre Archivo
\newcounter{popsec}
\newcommand{\coursec}[1]{%
  \stepcounter{popsec}\setcounter{popsub}{0}%
  \par\vspace{16pt}\noindent
  \begin{minipage}{\linewidth}
  \tikz[baseline=-0.7ex]\node[draw=popink,line width=1.6pt,fill=poporange,rounded corners=4pt,minimum size=22pt,inner sep=0]{\popdisplay\fontsize{12}{12}\selectfont\color{popcream}\thepopsec};%
  \hspace{8pt}{\popdisplay\fontsize{15}{17}\selectfont\color{popink}\MakeUppercase{#1}}\par
  \vspace{4pt}\noindent{\color{poporange}\rule{36pt}{3.6pt}}
  \end{minipage}\par\nopagebreak\vspace{8pt}
}
\newcounter{popsub}
\newcommand{\courssub}[1]{%
  \stepcounter{popsub}%
  \par\vspace{9pt}\noindent{\popxbold\fontsize{11.5}{13}\selectfont\color{popink}{\color{poporange}\thepopsec.\thepopsub}\hspace{6pt}#1}\par\nopagebreak\vspace{4pt}
}

% ============================================================
% Boîtes pédagogiques — même recette : fond blanc, bordure épaisse colorée,
% ombre dure encre, badge plein. Titre optionnel [#1].
% ============================================================
\tcbset{popbase/.style={breakable,enhanced,colback=white,boxrule=2.3pt,arc=9pt,
  shadow={3pt}{-3pt}{0pt}{popink},left=14pt,right=14pt,top=11pt,bottom=11pt,before skip=11pt,after skip=15pt,
  fontupper=\popsemi\fontsize{10.6}{14.5}\selectfont\color{popink},
  fontlower=\popsemi\fontsize{10.6}{14.5}\selectfont\color{popink}}}
% Coche dessinée (le glyphe ✓ n'existe pas dans les polices du document)
\newcommand{\popcheck}[1]{\tikz[baseline=-0.5ex]\draw[#1,line width=1.6pt,line cap=round,line join=round](-0.13,0.0)--(-0.04,-0.09)--(0.13,0.1);}
\newcommand{\popboxhead}[3]{\poppill{#1}{#2}{white}\ifx\relax#3\relax\else\hspace{6pt}{\popxbold\fontsize{10.6}{12}\selectfont\color{popink}#3}\fi\par\vspace{7pt}}

\newtcolorbox{aretenir}[1][]{popbase,colframe=popgreen,before upper={\popboxhead{À retenir}{popgreen}{#1}}}
\newtcolorbox{exemplebox}[1][]{popbase,colframe=poporange,before upper={\popboxhead{Exemple}{poporange}{#1}}}
\newtcolorbox{definition}[1][]{popbase,colframe=popblue,before upper={\popboxhead{Définition}{popblue}{#1}}}
\newtcolorbox{propriete}[1][]{popbase,colframe=poppurple,before upper={\popboxhead{Propriété}{poppurple}{#1}}}
\newtcolorbox{methode}[1][]{popbase,colframe=popgold,before upper={\popboxhead{Méthode}{popgold}{#1}}}
\newtcolorbox{attention}[1][]{popbase,colframe=poppink,before upper={\popboxhead{Attention}{poppink}{#1}}}
\newtcolorbox{experience}[1][]{popbase,colframe=popblue,colback=popblueL,before upper={\popboxhead{Expérience}{popblue}{#1}}}
% « Le savais-tu ? » : carte violette pleine (contexte, culture, Cameroun)
\newtcolorbox{savaistu}[1][]{popbase,colback=poppurple,colframe=popink,
  fontupper=\popsemi\fontsize{10.4}{14.2}\selectfont\color{white},
  before upper={\poppill{Le savais-tu\,?}{white}{poppurple}\ifx\relax#1\relax\else\hspace{6pt}{\popxbold\color{white}#1}\fi\par\vspace{7pt}}}
% Exercice résolu : énoncé en haut, \tcblower puis la solution sur fond crème
\newcounter{popexo}\newcounter{popexr}
\newtcolorbox{exoresolu}[1][]{popbase,colframe=poporange,colbacklower=popcream,
  segmentation style={popink,line width=1.2pt,dashed},
  before upper={\stepcounter{popexr}\popboxhead{Exercice résolu \thepopexr}{poporange}{#1}},
  before lower={\poppill{Solution}{popink}{popcream}\par\vspace{6pt}}}
% Exercice (non résolu en place) — la correction est dans la partie Corrigés
\newtcolorbox{exercice}[1][]{popbase,colframe=popink,
  before upper={\stepcounter{popexo}\popboxhead{Exercice \thepopexo}{popink}{#1}}}
% Corrigé
\newtcolorbox{corrige}[1][]{popbase,colframe=popgreen,colback=popgreenL,
  before upper={\popboxhead{Corrigé}{popgreen}{#1}}}
% Objectifs / prérequis / bilan
\newtcolorbox{objectifs}{popbase,colframe=popink,colback=poporangeL,
  before upper={\popboxhead{Objectifs de la leçon}{popink}{}}}
\newtcolorbox{prerequis}{popbase,colframe=popink,colback=popblueL,
  before upper={\popboxhead{Prérequis}{popblue}{}}}
\newtcolorbox{bilan}{popbase,colframe=popink,colback=white,boxrule=2.8pt,
  before upper={\popboxhead{Fiche bilan}{poporange}{L'essentiel à connaître pour l'examen}}}
% Figure encadrée avec légende
\newtcolorbox{popfigure}[1][]{enhanced,breakable=false,colback=white,colframe=popline,boxrule=1.4pt,arc=8pt,
  left=10pt,right=10pt,top=10pt,bottom=8pt,before skip=10pt,after skip=12pt,halign=center,
  lower separated=false,halign lower=center,
  fontlower=\popsemi\fontsize{9}{11.5}\selectfont\color{popmuted},
  #1}
\newcommand{\legende}[1]{\par\vspace{6pt}{\popsemi\fontsize{9}{11.5}\selectfont\color{popmuted}#1}}

% ============================================================
% Signature OnBuch+
% ============================================================
\newcommand{\poplogoinline}[1]{{\poplogo\fontsize{#1}{#1}\selectfont\color{popink}OnBuch\color{poporange}+}}
\newcommand{\signatureonbuch}{%
  \begin{tcolorbox}[enhanced,colback=popink,colframe=popink,boxrule=2.2pt,arc=10pt,
      drop shadow=poporange,left=14pt,right=14pt,top=10pt,bottom=10pt,before skip=0pt,after skip=0pt]
    \tikz[baseline=-0.7ex]\node[circle,fill=popgreen,draw=popcream,line width=1.3pt,minimum size=22pt,inner sep=0]{\popcheck{white}};%
    \hspace{9pt}\begin{minipage}[c]{0.62\linewidth}
      {\popxbold\fontsize{12}{13}\selectfont\color{popcream}Signé\ \textbullet\ {\poplogo OnBuch\color{poporange}+}}\par\vspace{2pt}
      {\raggedright\popsemi\fontsize{8}{10.5}\selectfont\color{popline}\MakeUppercase{Équipe pédagogique OnBuch+}\\\MakeUppercase{Conforme au programme officiel MINESEC}\par}
    \end{minipage}\hfill
    {\poplogo\fontsize{22}{22}\selectfont\color{popcream}OnBuch\color{poporange}+}
  \end{tcolorbox}%
}

% ============================================================
% Page de garde
% #1 titre · #2 matière · #3 niveau · #4 module · #5 n° leçon · #6 durée ·
% #7 description / objectifs (texte ou itemize)
% ============================================================
\newcommand{\pagegarde}[7]{%
  \thispagestyle{empty}
  \newgeometry{a4paper,margin=1.7cm,top=1.6cm,bottom=1.4cm}%
  \begin{tikzpicture}[remember picture,overlay]
    \fill[poporange!18] ([xshift=-3.2cm,yshift=-3.6cm]current page.north east) circle (4.6cm);
    \fill[poppurple!14] ([xshift=2.4cm,yshift=7.6cm]current page.south west) circle (3.8cm);
  \end{tikzpicture}%
  {\poplogo\fontsize{30}{30}\selectfont\color{popink}OnBuch\color{poporange}+}\hfill
  \raisebox{0.6ex}{\poppill{Cours signé \textbullet\ Premium}{popink}{popcream}}\par
  \vspace{1pt}\popsquiggle{poppurple}\par
  \vspace{8pt}
  \begin{tcolorbox}[enhanced,colback=poporange,colframe=popink,boxrule=2.8pt,arc=13pt,
      drop shadow=popink,left=18pt,right=18pt,top=12pt,bottom=13pt,after skip=10pt]
    \poppill{#2 \textbullet\ #3}{popink}{popcream}\hspace{5pt}\poppill{#4}{white}{popink}\par\vspace{8pt}
    {\popdisplay\fontsize{14}{14}\selectfont\color{popink}LEÇON #5}\par\vspace{4pt}
    {\raggedright\hyphenpenalty=10000\exhyphenpenalty=10000\popdisplay\fontsize{25}{27}\selectfont\color{popcream}\MakeUppercase{#1}\par}
    \vspace{8pt}
    {\popxbold\fontsize{9.5}{9.5}\selectfont\color{popink}\MakeUppercase{Durée conseillée : #6}}
  \end{tcolorbox}
  \begin{tcolorbox}[enhanced,colback=white,colframe=popink,boxrule=2.2pt,arc=10pt,
      drop shadow=popink,left=16pt,right=16pt,top=10pt,bottom=10pt,after skip=8pt]
    \poppill{Dans cette leçon}{poppurple}{white}\par\vspace{5pt}
    \popsemi\fontsize{9.3}{12.6}\selectfont\color{popink}#7
  \end{tcolorbox}
  \noindent\begin{minipage}[t]{0.487\linewidth}
    \begin{tcolorbox}[enhanced,colback=popgreen,colframe=popink,boxrule=2.2pt,arc=10pt,
        drop shadow=popink,left=11pt,right=11pt,top=10pt,bottom=10pt,height=3.1cm,valign=center]
      \tcbox[colback=white,colframe=popink,boxrule=1.3pt,arc=5pt,boxsep=0pt,nobeforeafter,
        left=3pt,right=3pt,top=3pt,bottom=3pt,valign=center]{\qrcode[height=1.9cm]{https://play.google.com/store/apps/details?id=com.luvvix.onbuch&hl=fr}}%
      \hspace{8pt}\begin{minipage}[c]{0.5\linewidth}
        {\popdisplay\fontsize{11}{12.5}\selectfont\color{white}\MakeUppercase{Télécharge OnBuch}}\par\vspace{4pt}
        {\raggedright\popsemi\fontsize{8.4}{11}\selectfont\color{white}Gratuit \textbullet\ Google Play\\Tuteur IA, annales, concours, résultats\par}
      \end{minipage}
    \end{tcolorbox}
  \end{minipage}\hfill
  \begin{minipage}[t]{0.487\linewidth}
    \begin{tcolorbox}[enhanced,colback=poppurple,colframe=popink,boxrule=2.2pt,arc=10pt,
        drop shadow=popink,left=11pt,right=11pt,top=10pt,bottom=10pt,height=3.1cm,valign=center]
      \tikz[baseline=-0.7ex]\node[circle,fill=white,draw=popink,line width=1.3pt,minimum size=1.6cm,inner sep=0]{\popdisplay\fontsize{24}{24}\selectfont\color{poppurple}+};%
      \hspace{8pt}\begin{minipage}[c]{0.6\linewidth}
        {\popdisplay\fontsize{11}{12.5}\selectfont\color{white}\MakeUppercase{Passe à OnBuch+}}\par\vspace{4pt}
        {\raggedright\popsemi\fontsize{8.4}{11}\selectfont\color{white}Tous les cours signés, Tuteur IA illimité, fiches PDF — onglet OnBuch+ de l'app\par}
      \end{minipage}
    \end{tcolorbox}
  \end{minipage}
  \vspace{8pt}
  \popcontenu
  \vfill
  \signatureonbuch
  \restoregeometry
  \newpage
}

% Rangée « Ce cours contient » (page de garde)
\newcommand{\poptile}[3]{% #1 couleur #2 gros texte #3 légende
  \begin{tcolorbox}[enhanced,colback=#1,colframe=popink,boxrule=2pt,arc=9pt,shadow={2.5pt}{-2.5pt}{0pt}{popink},
      left=6pt,right=6pt,top=9pt,bottom=9pt,halign=center,valign=center,height=1.75cm,width=\linewidth,nobeforeafter]
    {\popdisplay\fontsize{12.5}{13}\selectfont\color{white}\MakeUppercase{#2}}\par\vspace{3pt}
    {\centering\popsemi\fontsize{7.8}{9.5}\selectfont\color{white}#3\par}
  \end{tcolorbox}}
\newcommand{\popcontenu}{%
  \par\noindent{\popxboldls\fontsize{7.5}{7.5}\selectfont\color{popmuted}CE COURS SIGNÉ CONTIENT}\par\vspace{7pt}
  \noindent\begin{minipage}[t]{0.235\linewidth}\poptile{popblue}{Cours}{complet, illustré, pas à pas}\end{minipage}\hfill
  \begin{minipage}[t]{0.235\linewidth}\poptile{poporange}{Méthodes}{et exercices résolus}\end{minipage}\hfill
  \begin{minipage}[t]{0.235\linewidth}\poptile{poppink}{Exercices}{type examen + corrigés}\end{minipage}\hfill
  \begin{minipage}[t]{0.235\linewidth}\poptile{popgreen}{Fiche bilan}{l'essentiel à retenir}\end{minipage}\par}

% Sommaire
\newtcolorbox{sommaire}{enhanced,colback=white,colframe=popink,boxrule=2.2pt,arc=10pt,
  drop shadow=popink,left=18pt,right=18pt,top=13pt,bottom=13pt,before skip=6pt,after skip=20pt,
  before upper={\poppill{Sommaire}{poppurple}{white}\par\vspace{9pt}\popsemi\fontsize{10.6}{14.5}\selectfont\color{popink}}}

% Signature de fin de cours — poussée tout en bas de la dernière page
\newcommand{\findecours}{%
  \par\vfill\nopagebreak
  \begin{center}\popsquiggle{poporange}\end{center}\vspace{4pt}
  \signatureonbuch
}
% Glyphes absents de Fira Math : redéfinis à partir de glyphes disponibles
\AtBeginDocument{%
  \def\setminus{\mathbin{\backslash}}\let\smallsetminus\setminus
  \def\vdots{\mathord{\vbox{\baselineskip3.2pt\lineskiplimit0pt\kern4pt\hbox{.}\hbox{.}\hbox{.}}}}%
  \def\ddots{\mathinner{\mkern1mu\raise6.5pt\hbox{.}\mkern2mu\raise3.6pt\hbox{.}\mkern2mu\raise0.7pt\hbox{.}\mkern1mu}}%
  \def\triangle{\mathord{\scalebox{0.9}{$\symup{\Delta}$}}}%
  \def\nabla{\mathord{\raisebox{\depth}{\scalebox{1}[-1]{$\symup{\Delta}$}}}}%
  \def\checkmark{\popcheck{popdark}}%
  \def\star{\mathbin{\ast}}\def\bigstar{\mathord{\ast}}%
  \def\diamond{\mathbin{\scalebox{0.6}{$\Diamond$}}}%
  \def\bigcup{\mathop{\mathchoice{\raisebox{-0.3em}{\scalebox{1.7}{$\cup$}}}{\scalebox{1.25}{$\cup$}}{\cup}{\cup}}\displaylimits}%
  \def\bigcap{\mathop{\mathchoice{\raisebox{-0.3em}{\scalebox{1.7}{$\cap$}}}{\scalebox{1.25}{$\cap$}}{\cap}{\cap}}\displaylimits}%
  \def\lesseqqgtr{\mathrel{\lessgtr}}\def\bigoplus{\mathop{\oplus}\displaylimits}%
}
\providecommand{\ssi}{si et seulement si} % abréviation fréquente

%% FIGFIX-BEGIN v5 — correctifs globaux des figures (docs/onbuch-generation/tools/figfix)
\usepackage{adjustbox}\usepackage{etoolbox}\tcbuselibrary{raster}
% toute figure TikZ/pgfplots plus large que la ligne est réduite à la largeur disponible
\BeforeBeginEnvironment{tikzpicture}{\begin{adjustbox}{max width=\linewidth}}
\AfterEndEnvironment{tikzpicture}{\end{adjustbox}}
% légendes pgfplots lisibles par-dessus les courbes
\pgfplotsset{every axis/.append style={legend style={fill=white,fill opacity=0.92,text opacity=1,draw=black!15,inner sep=3pt}}}
% étiquettes d'arêtes (syntaxe edge["..."]) avec fond blanc
\tikzset{every edge quotes/.append style={fill=white,inner sep=1.2pt}}
%% FIGFIX-END
\begin{document}

\pagegarde{Les composés oxygénés : alcools, aldéhydes, cétones}{Chimie}{1ère D}{Module 1 : Chimie organique}{03}{7 h}{Cette leçon présente les principales familles de composés organiques oxygénés : alcools, aldéhydes et cétones. Elle développe leur nomenclature, leurs classes, leurs propriétés physiques et les tests caractéristiques permettant de les identifier au laboratoire.
\vspace{4pt}
\begin{multicols}{2}\small
\begin{itemize}
\item À la fin de cette leçon, je sais identifier le groupe fonctionnel d'un alcool, d'un aldéhyde, d'une cétone et d'un éther-oxyde dans une formule développée.
\item À la fin de cette leçon, je sais écrire la formule générale et la formule brute d'un alcool, d'un aldéhyde et d'une cétone saturés.
\item À la fin de cette leçon, je sais nommer un alcool, un éther-oxyde, un aldéhyde et une cétone selon les règles de nomenclature officielle.
\item À la fin de cette leçon, je sais distinguer les trois classes d'alcools (primaire, secondaire, tertiaire) à partir de leur formule développée.
\item À la fin de cette leçon, je sais expliquer les propriétés physiques des alcools (ébullition, solubilité) par la présence de liaisons hydrogène.
\item À la fin de cette leçon, je sais citer des exemples de polyalcools (glycol, glycérol) et leurs usages.
\end{itemize}\end{multicols}}

\begin{sommaire}
\begin{multicols}{2}
\begin{enumerate}
\item Généralités sur les composés oxygénés
\item Les alcools : formules, nomenclature et classes
\item Propriétés physiques des alcools et polyalcools
\item Aldéhydes et cétones : formules et nomenclature
\item Tests d'identification des aldéhydes et cétones
\item Synthèse et applications au quotidien
\item Travaux pratiques
\item Méthodes et exercices résolus
\item Exercices
\item Corrigés des exercices
\item Fiche bilan
\end{enumerate}
\end{multicols}
\end{sommaire}

\begin{objectifs}
\begin{itemize}
\item À la fin de cette leçon, je sais identifier le groupe fonctionnel d'un alcool, d'un aldéhyde, d'une cétone et d'un éther-oxyde dans une formule développée.
\item À la fin de cette leçon, je sais écrire la formule générale et la formule brute d'un alcool, d'un aldéhyde et d'une cétone saturés.
\item À la fin de cette leçon, je sais nommer un alcool, un éther-oxyde, un aldéhyde et une cétone selon les règles de nomenclature officielle.
\item À la fin de cette leçon, je sais distinguer les trois classes d'alcools (primaire, secondaire, tertiaire) à partir de leur formule développée.
\item À la fin de cette leçon, je sais expliquer les propriétés physiques des alcools (ébullition, solubilité) par la présence de liaisons hydrogène.
\item À la fin de cette leçon, je sais citer des exemples de polyalcools (glycol, glycérol) et leurs usages.
\item À la fin de cette leçon, je sais distinguer expérimentalement un aldéhyde d'une cétone à l'aide de la 2,4-DNPH, du réactif de Schiff, de la liqueur de Fehling et du réactif de Tollens.
\item À la fin de cette leçon, je sais écrire les équations-bilans des réactions de test des aldéhydes.
\end{itemize}
\end{objectifs}

\begin{prerequis}
\begin{itemize}
\item Formules brutes, développées et semi-développées des molécules organiques (Seconde)
\item Nomenclature des alcanes et règles de numérotation de la chaîne carbonée
\item Notion d'isomérie (isomères de constitution)
\item Écriture et équilibrage d'une équation-bilan chimique
\item Notion de réaction d'oxydoréduction (oxydant, réducteur)
\end{itemize}
\end{prerequis}

\newpage

\coursec{Généralités sur les composés oxygénés}

\courssub{Situation d'accroche et problématique}

À Douala comme à Bafoussam, l'éthanol est omniprésent : dans le vin de palme et le bili-bili qui animent les veillées, dans les flacons de désinfectant à \SI{70}{\%} des pharmacies, ou encore comme additif dans l'essence à la station SONARA. Pourtant, un simple déplacement d'un atome d'oxygène transforme cet alcool buvable en méthanol, responsable d'intoxications mortelles lors de la consommation d'alcools frelatés. Comment distinguer ces molécules si proches ? Comment les nommer, les classer et les identifier formellement au laboratoire ?

La réponse tient en une idée directrice de la chimie organique : ce ne sont pas seulement les atomes présents qui comptent, mais la \emph{manière dont ils sont assemblés}. Un atome d'oxygène portant un hydrogène donne un alcool ; le même oxygène en double liaison avec un carbone en bout de chaîne donne un aldéhyde ; au milieu de la chaîne, une cétone ; coincé entre deux carbones par des liaisons simples, un éther-oxyde. Cette leçon donne les clés pour comprendre, nommer et identifier la grande famille des \cle{composés oxygénés}.

\textbf{Problématique :} comment reconnaître, nommer et distinguer les différentes familles de composés oxygénés (alcools, aldéhydes, cétones, éthers-oxydes) à partir de la structure de leurs molécules ?

\courssub{Qu'est-ce qu'un composé oxygéné ?}

Jusqu'ici, tu as rencontré les hydrocarbures (alcanes, alcènes, alcynes, composés aromatiques) : leurs molécules ne contiennent que du carbone et de l'hydrogène. La vie quotidienne et l'industrie camerounaise utilisent pourtant d'immenses quantités de molécules organiques qui contiennent, en plus, de l'oxygène : l'éthanol du vin de palme, le glucose du jus de bissap, l'acétone du dissolvant à vernis, le formol des laboratoires de biologie. Toutes appartiennent à la grande famille des composés oxygénés.

\begin{definition}[Composé oxygéné]
Un \cle{composé oxygéné} est un composé organique dont la molécule contient \textbf{au moins un atome d'oxygène}, en plus du carbone et de l'hydrogène.

La \cle{famille} à laquelle il appartient est déterminée par son \cle{groupe fonctionnel}, c'est-à-dire le groupement d'atomes particulier, incluant l'oxygène, qui confère à la molécule ses propriétés chimiques caractéristiques.
\end{definition}

L'idée directrice est la suivante : une molécule organique se décompose conceptuellement en deux parties.
\begin{itemize}
\item Une partie \textbf{hydrocarbonée}, notée $R$ (ou $R'$, $R''$), appelée \cle{groupe alkyle} : c'est le « squelette » de la molécule, chimiquement assez peu réactif. Exemples : $R = \ce{CH3-CH2}-$ (groupe éthyle), $R = \ce{CH3}-$ (groupe méthyle), $R = \ce{CH3-CH2-CH2}-$ (groupe propyle).
\item Un \textbf{groupe fonctionnel}, petite zone de la molécule où se concentre la réactivité chimique. C'est lui qui « signe » la famille.
\end{itemize}

\begin{attention}[L'atome d'oxygène ne suffit pas à tout dire]
Deux molécules peuvent avoir la \emph{même formule brute} et appartenir à des familles différentes. Par exemple, \ce{C2H6O} correspond à l'éthanol \ce{CH3-CH2-OH} (un alcool) \emph{et} à la \textbf{méthoxyméthane} \ce{CH3-O-CH3} (un éther-oxyde). Seule la \textbf{formule développée} (ou semi-développée) permet de trancher. Retiens : en chimie organique, on raisonne toujours sur la formule développée, jamais sur la seule formule brute.
\end{attention}

\courssub{Les quatre familles de composés oxygénés du programme}

\textbf{Les alcools : le groupe hydroxyle $-\ce{OH}$}\par

Dans un alcool, l'atome d'oxygène est lié \emph{d'une part} à un atome de carbone (tétravalent, porteur d'une chaîne carbonée), \emph{d'autre part} à un atome d'hydrogène. Le groupe fonctionnel $-\ce{OH}$ s'appelle le groupe \cle{hydroxyle}. La formule générale d'un alcool saturé acyclique s'écrit :

\[ R-\ce{OH} \qquad \text{avec } R \text{ groupe alkyle (donc } R \neq \ce{H}\text{)} \]

\textbf{Précision importante :} dans $R-\ce{OH}$, $R$ est un groupe \emph{carboné} (alkyle). Si l'on remplaçait $R$ par un simple H, on obtiendrait \ce{H-OH}, c'est-à-dire l'eau, qui n'est évidemment pas un alcool. Le plus simple des alcools est donc le méthanol \ce{CH3-OH} ($R = \ce{CH3}-$).

L'exemple emblématique est l'\cle{éthanol} \ce{CH3-CH2-OH}, de formule brute \ce{C2H6O} : c'est l'alcool des boissons fermentées traditionnelles (vin de palme, bili-bili), du gel hydroalcoolique et du désinfectant à \SI{70}{\%} vendu en pharmacie, et un solvant industriel majeur.

\courssub{Aldéhydes et cétones : le groupe carbonyle $\ce{C=O}$}

Ici, l'atome d'oxygène n'est plus lié à un hydrogène : il forme une \textbf{double liaison} avec un atome de carbone. Le groupement $\ce{C=O}$ s'appelle le groupe \cle{carbonyle}. Mais la présence du carbonyle ne suffit pas à déterminer la famille. Tout dépend de la \emph{position} du carbone carbonyle dans la chaîne carbonée.

\begin{definition}[Aldéhyde et cétone]
\begin{itemize}
\item Un \cle{aldéhyde} est un composé oxygéné dont le groupe carbonyle est porté par un carbone \textbf{en bout de chaîne} : ce carbone est donc aussi lié à \emph{au moins un atome d'hydrogène}. Formule générale : $R-\ce{CHO}$ (souvent écrit $R-\ce{CH=O}$ en développé). Le cas particulier $R = \ce{H}$ donne le méthanal \ce{H-CHO}.
\item Une \cle{cétone} est un composé oxygéné dont le groupe carbonyle est porté par un carbone \textbf{à l'intérieur de la chaîne} : ce carbone est lié à \emph{deux autres atomes de carbone}. Formule générale : $R-\ce{CO}-R'$, avec $R$ et $R'$ groupes alkyles (donc $R \neq \ce{H}$ et $R' \neq \ce{H}$).
\end{itemize}
\end{definition}

Moyen mnémotechnique : dans un aldéhyde, le carbone du carbonyle garde « un H sous le coude » (en bout de chaîne) ; dans une cétone, il est « coincé entre deux carbones ».

Exemples de référence :
\begin{itemize}
\item Le \cle{méthanal} \ce{H-CHO} (formule brute \ce{CH2O}), l'aldéhyde le plus simple : sa solution aqueuse à \SI{37}{\%} massique, le \emph{formol}, sert à conserver les organes et spécimens au laboratoire.
\item L'\cle{éthanal} \ce{CH3-CHO} (\ce{C2H4O}), aldéhyde à deux carbones.
\item La \cle{propanone} \ce{CH3-CO-CH3} (formule brute \ce{C3H6O}), appelée couramment \emph{acétone} : c'est la cétone la plus simple, utilisée comme dissolvant pour vernis à ongles et solvant industriel.
\end{itemize}

\begin{attention}[Pas de cétone à 2 carbones]
Pour une cétone, le carbone carbonyle doit être lié à \textbf{deux} groupes alkyles. Il faut donc au minimum 3 atomes de carbone au total. Il n'existe \emph{pas} de cétone de formule brute \ce{C2H4O}.
\end{attention}

\courssub{Les éthers-oxydes : le pont oxygène $-\ce{O}-$}

Dans un éther-oxyde, l'atome d'oxygène est lié par des \emph{liaisons simples} à \textbf{deux atomes de carbone} : il forme un « pont » entre deux groupes alkyles. La formule générale est :

\[ R-\ce{O}-R' \]

L'exemple historique est l'\cle{éthoxyéthane} \ce{CH3-CH2-O-CH2-CH3} (formule brute \ce{C4H10O}), le fameux « éther » utilisé autrefois comme anesthésiant en chirurgie. Les éthers-oxydes sont chimiquement assez peu réactifs ; on les utilise surtout comme solvants (ex. : éther de pétrole, THF).

\begin{attention}[Éther ou alcool : même formule brute, familles différentes]
\ce{CH3-O-CH3} (méthoxyméthane, un éther) et \ce{CH3-CH2-OH} (éthanol, un alcool) ont tous deux pour formule brute \ce{C2H6O}. Ce sont des \cle{isomères de fonction}. Pour les distinguer, cherche l'hydrogène porté par l'oxygène : s'il existe, c'est un alcool ; sinon, c'est un éther.
\end{attention}

\courssub{Identifier visuellement les quatre groupes fonctionnels}

La figure ci-dessous rassemble les quatre groupes fonctionnels à reconnaître instantanément. Entraîne-toi à les repérer dans n'importe quelle formule développée : c'est le réflexe fondamental de tout le chapitre.

\begin{popfigure}
\begin{center}
\begin{tikzpicture}[node distance=1.5cm, every node/.style={align=center, font=\small}]
% Alcool
\node (alc) at (0,0) {\chemfig{[:30]R-[:-30]O-[:-90]H}};
\node[below=0.2cm of alc] (alcT) {\textcolor{popblue}{\textbf{Alcool}}\\groupe \cle{hydroxyle} $-\ce{OH}$};
% Aldéhyde
\node (ald) at (4,0) {\chemfig{[:30]R-[:-30]C(=[:90]O)-[:-90]H}};
\node[below=0.2cm of ald] (aldT) {\textcolor{poporange}{\textbf{Aldéhyde}}\\carbonyle en \textbf{bout de chaîne}\\$-\ce{CHO}$};
% Cétone
\node (ket) at (8,0) {\chemfig{[:30]R-[:-30]C(=[:90]O)-[:-30]R'}};
\node[below=0.2cm of ket] (ketT) {\textcolor{poppurple}{\textbf{Cétone}}\\carbonyle \textbf{dans la chaîne}\\$-\ce{CO}-$};
% Éther
\node (eth) at (12,0) {\chemfig{[:30]R-[:-30]O-[:-30]R'}};
\node[below=0.2cm of eth] (ethT) {\textcolor{popgreen}{\textbf{Éther-oxyde}}\\pont $-\ce{O}-$};
\end{tikzpicture}
\end{center}
\legende{Les quatre groupes fonctionnels oxygénés. L'oxygène est le point de départ du diagnostic : porte-t-il un H (alcool) ? Est-il en double liaison en bout de chaîne (aldéhyde) ou à l'intérieur (cétone) ? Ou entre deux carbones par liaisons simples (éther) ?}
\end{popfigure}

\begin{methode}[Identifier la famille d'un composé oxygéné]
Face à une formule développée ou semi-développée inconnue, procède en quatre étapes systématiques :
\begin{enumerate}
\item \textbf{Repérer l'atome (ou les atomes) d'oxygène} dans la molécule.
\item \textbf{Observer son environnement immédiat} :
      \begin{itemize}
      \item L'oxygène porte-t-il un atome d'hydrogène (liaison O--H) ? $\Rightarrow$ Groupe hydroxyle $\Rightarrow$ \textbf{Alcool}.
      \item L'oxygène est-il en \textbf{double liaison $\ce{C=O}$} (carbonyle) ?
            \begin{itemize}
            \item Le carbone du carbonyle porte-t-il un H (carbone terminal) ? $\Rightarrow$ \textbf{Aldéhyde}.
            \item Le carbone du carbonyle est-il lié à deux carbones (carbone interne) ? $\Rightarrow$ \textbf{Cétone}.
            \end{itemize}
      \item L'oxygène est-il lié par \textbf{deux liaisons simples à deux carbones} (sans H, sans double liaison) ? $\Rightarrow$ Pont $-\ce{O}-$ $\Rightarrow$ \textbf{Éther-oxyde}.
      \end{itemize}
\end{enumerate}
\end{methode}

\begin{exemplebox}[Application immédiate de la méthode]
Identifions la famille de trois molécules de formule brute \ce{C3H6O} ou \ce{C3H8O} :
\begin{itemize}
\item \chemfig{CH_3-CH_2-CH_2-OH} : l'oxygène porte un H $\Rightarrow$ groupe hydroxyle $\Rightarrow$ \textbf{alcool} (propan-1-ol, \ce{C3H8O}).
\item \chemfig{CH_3-CH_2-C(=[1]O)-[7]H} : double liaison $\ce{C=O}$, le carbone du carbonyle porte un H $\Rightarrow$ \textbf{aldéhyde} (propanal, \ce{C3H6O}).
\item \chemfig{CH_3-C(=[1]O)-[7]CH_3} : double liaison $\ce{C=O}$, le carbone du carbonyle est entre deux carbones $\Rightarrow$ \textbf{cétone} (propanone, \ce{C3H6O}).
\end{itemize}
Remarque cruciale : le propanal et la propanone ont la même formule brute \ce{C3H6O} : ce sont des \textbf{isomères de fonction}. Cela justifie l'existence de tests chimiques spécifiques pour les distinguer (section 5).
\end{exemplebox}

\courssub{Tableau récapitulatif des quatre familles}

\begin{center}
\begin{tabularx}{\linewidth}{|l|X|X|X|}
\hline
\rowcolor{popblueL}
\textbf{Famille} & \textbf{Formule générale} & \textbf{Groupe fonctionnel} & \textbf{Exemple courant (contexte camerounais)} \\
\hline
Alcool & $R-\ce{OH}$ & Hydroxyle $-\ce{OH}$ & Éthanol \ce{CH3-CH2-OH} (vin de palme, bili-bili, désinfectant \SI{70}{\%}) \\
\hline
Aldéhyde & $R-\ce{CHO}$ & Carbonyle $\ce{C=O}$ en bout de chaîne & Méthanal \ce{H-CHO} (formol, conservation spécimens) \\
\hline
Cétone & $R-\ce{CO}-R'$ & Carbonyle $\ce{C=O}$ à l'intérieur de la chaîne & Propanone \ce{CH3-CO-CH3} (acétone, dissolvant vernis) \\
\hline
Éther-oxyde & $R-\ce{O}-R'$ & Pont oxygène $-\ce{O}-$ & Éthoxyéthane \ce{C2H5-O-C2H5} (ancien anesthésique) \\
\hline
\end{tabularx}
\end{center}

\courssub{Le groupe fonctionnel commande les propriétés chimiques}

Pourquoi accorde-t-on tant d'importance au groupe fonctionnel ? Parce que c'est lui qui détermine le \emph{comportement chimique} de la molécule : toutes les molécules d'une même famille réagissent de manière analogue face à un même réactif.

\begin{propriete}[Famille et réactivité]
Tous les composés portant le même groupe fonctionnel possèdent des \textbf{propriétés chimiques semblables} : ils subissent les mêmes types de réactions et donnent les mêmes tests d'identification. Ainsi :
\begin{itemize}
\item tous les alcools réagissent avec le sodium métallique en dégageant du dihydrogène ;
\item tous les aldéhydes réduisent la liqueur de Fehling et le réactif de Tollens (ce sont des \emph{réducteurs}) ;
\item les cétones, elles, ne subissent pas ces réactions d'oxydation ménagée : c'est ce qui permet de les distinguer des aldéhydes.
\end{itemize}
\end{propriete}

Un premier exemple de réaction caractéristique : l'action du sodium sur un alcool. Le sodium métallique \ce{Na} arrache l'hydrogène du groupe hydroxyle ; il se forme un \emph{alcoolate} (sel de sodium de l'alcool) et un dégagement de dihydrogène :

\[ \ce{2 CH3CH2OH + 2 Na -> 2 CH3CH2O- Na+ + H2 ^} \]

Cette réaction, vive et exothermique, est caractéristique du groupe $-\ce{OH}$ : elle se produit avec \emph{tous} les alcools (et aussi avec l'eau, de façon encore plus violente, d'où l'interdiction formelle de jeter du sodium dans un évier !). Elle illustre parfaitement le principe : \textbf{même groupe fonctionnel, même réactivité}.

\begin{attention}[Le groupe $-\ce{OH}$ d'un alcool n'est \textbf{pas} l'ion hydroxyde \ce{HO-}]
Erreur très fréquente au Probatoire ! Le groupe hydroxyle $-\ce{OH}$ d'un alcool est \textbf{lié par covalence} à un atome de carbone : il ne se libère \emph{pas} en ions \ce{HO-} dans l'eau. Conséquences majeures :
\begin{itemize}
\item une solution d'éthanol dans l'eau est \textbf{neutre} : son pH vaut environ 7, elle ne rougit pas le papier tournesol ;
\item l'alcool \textbf{n'est pas une base}, contrairement à la soude \ce{NaOH} qui, elle, libère bien des ions hydroxyde \ce{HO-} en solution.
\end{itemize}
Retiens la formulation exacte : « le groupe hydroxyle est un groupe \emph{fonctionnel neutre}, l'ion hydroxyde est un \emph{ion basique} ». L'alcool peut certes céder l'hydrogène de son $-\ce{OH}$ à un métal très réducteur comme le sodium (réaction vue ci-dessus), mais jamais spontanément dans l'eau.
\end{attention}

\begin{savaistu}[Le formol, gardien des collections scientifiques... et danger public]
Le liquide dans lequel baignent les organes et les spécimens des laboratoires de biologie — grenouilles, serpents, fœtus — est le \emph{formol} : une solution aqueuse à environ \SI{37}{\%} en masse de \textbf{méthanal} \ce{H-CHO}, l'aldéhyde le plus simple. Le méthanal est un excellent agent de conservation car il « fige » les protéines en créant des ponts covalents entre elles (on dit qu'il les \emph{dénature}), ce qui empêche toute décomposition microbienne. C'est aussi ce qui le rend toxique, irritant pour les voies respiratoires et cancérigène avéré (CIRC groupe 1) : il faut manipuler le formol sous la hotte, avec des gants et des lunettes. Au Cameroun, des contrôles sanitaires saisissent régulièrement du poisson ou de la viande frauduleusement traités au formol par des commerçants peu scrupuleux pour masquer un début de décomposition : une pratique dangereuse, strictement interdite et punie par la loi.
\end{savaistu}

\begin{exoresolu}[Identification de familles à partir de formules semi-développées]
\textbf{Énoncé.} On donne les formules semi-développées de quatre composés organiques :
\[
\text{A : } \ce{CH3-CH2-CH2-OH} \qquad
\text{B : } \ce{CH3-CH2-CHO} \qquad
\text{C : } \ce{CH3-CO-CH3} \qquad
\text{D : } \ce{CH3-CH2-O-CH3}
\]
\begin{enumerate}
\item Pour chaque composé, repérer l'atome d'oxygène et décrire son environnement de liaison.
\item En déduire la famille de chaque composé et donner son nom usuel ou systématique.
\item Les composés B et C ont-ils la même formule brute ? Commenter le lien avec l'isomérie.
\end{enumerate}
\tcblower
\textbf{Solution détaillée.}
\begin{enumerate}
\item \textbf{Composé A :} l'oxygène est lié à un carbone (liaison simple) et porte un hydrogène (liaison O--H) : groupe $-\ce{OH}$ (hydroxyle).\\
      \textbf{Composé B :} l'oxygène est en double liaison avec un carbone terminal (bout de chaîne), ce carbone porte aussi un H : groupe $-\ce{CHO}$ (carbonyle aldéhydique).\\
      \textbf{Composé C :} l'oxygène est en double liaison avec un carbone situé entre deux groupes méthyles (carbone secondaire) : groupe $-\ce{CO}-$ (carbonyle cétonique).\\
      \textbf{Composé D :} l'oxygène est lié par deux liaisons simples à deux carbones (un éthyle, un méthyle), sans hydrogène, sans double liaison : pont $-\ce{O}-$.
\item A est un \textbf{alcool} : propan-1-ol (ou propanol-1).\\
      B est un \textbf{aldéhyde} : propanal.\\
      C est une \textbf{cétone} : propanone (acétone).\\
      D est un \textbf{éther-oxyde} : méthoxyéthane.
\item Comptons les atomes pour B : 3 C, 6 H (5 sur la chaîne + 1 sur CHO), 1 O $\Rightarrow$ \ce{C3H6O}.\\
      Pour C : 3 C, 6 H (3+3 sur les deux CH3), 1 O $\Rightarrow$ \ce{C3H6O}.\\
      B et C ont donc la \textbf{même formule brute} \ce{C3H6O} tout en appartenant à deux familles différentes (aldéhyde vs cétone) : ce sont des \textbf{isomères de fonction}. Cela confirme qu'on ne peut pas déterminer la famille d'un composé oxygéné à partir de sa seule formule brute : il faut sa formule développée ou un test chimique d'identification.
\end{enumerate}
\end{exoresolu}

\begin{aretenir}[L'essentiel de la section 1]
\begin{itemize}
\item Un \cle{composé oxygéné} contient au moins un atome d'oxygène ; sa famille est fixée par son \cle{groupe fonctionnel}.
\item Quatre familles au programme : \textbf{alcool} $R-\ce{OH}$ (hydroxyle), \textbf{aldéhyde} $R-\ce{CHO}$ (carbonyle terminal), \textbf{cétone} $R-\ce{CO}-R'$ (carbonyle interne), \textbf{éther-oxyde} $R-\ce{O}-R'$ (pont oxygène).
\item Méthode de diagnostic : repérer l'oxygène $\rightarrow$ observer liaisons (H ? double liaison ? position ?).
\item Le groupe fonctionnel commande la réactivité : même groupe $\Rightarrow$ mêmes réactions, mêmes tests.
\item Le $-\ce{OH}$ alcoolique est covalent, neutre, non basique (différent de l'ion \ce{HO-}).
\end{itemize}
\end{aretenir}

% ============================================================
% SECTION 2 : LES ALCOOLS : FORMULES, NOMENCLATURE ET CLASSES
% ============================================================
\coursec{Les alcools : formules, nomenclature et classes}

\courssub{Formule générale et formule brute des alcools saturés acycliques}

Un alcool saturé acyclique (sans cycle, sans double liaison C=C) dérive d'un alcane par substitution d'un atome d'hydrogène par un groupe hydroxyle $-\ce{OH}$. Sa formule brute suit la loi générale :

\[ \textbf{$C_nH_{2n+2}O$} \quad (n \ge 1) \]

\begin{exemplebox}[Vérification sur les premiers alcools]
\begin{center}
\begin{tabular}{c|c|c|c}
\textbf{Alcool} & \textbf{Formule semi-développée} & \textbf{n} & \textbf{Formule brute $C_nH_{2n+2}O$} \\
\hline
Méthanol & \ce{CH3-OH} & 1 & \ce{CH4O} \\
Éthanol & \ce{CH3-CH2-OH} & 2 & \ce{C2H6O} \\
Propan-1-ol & \ce{CH3-CH2-CH2-OH} & 3 & \ce{C3H8O} \\
Butan-1-ol & \ce{CH3-CH2-CH2-CH2-OH} & 4 & \ce{C4H10O} \\
\end{tabular}
\end{center}
Attention : cette formule $C_nH_{2n+2}O$ est aussi celle des \textbf{éthers-oxydes} saturés acycliques (isomérie de fonction). Elle ne caractérise pas à elle seule la fonction alcool.
\end{exemplebox}

\courssub{Nomenclature systématique (IUPAC) des alcools}

La nomenclature suit des règles précises, identiques à celles des alcanes avec une priorité pour le groupe $-\ce{OH}$.

\begin{methode}[Nommer un alcool saturé acyclique (IUPAC)]
\begin{enumerate}
\item \textbf{Choisir la chaîne principale :} la plus longue chaîne carbonée \textbf{portant le groupe $-\ce{OH}$}. Elle détermine le radical de base (méth-, éth-, prop-, but-, pent-, hex-, hept-, oct-, non-, déc-).
\item \textbf{Numéroter la chaîne :} le numérotage part de l'extrémité la plus proche du groupe $-\ce{OH}$ (règle du plus bas locant pour la fonction prioritaire).
\item \textbf{Indiquer la position du $-\ce{OH}$ :} par un chiffre (locant) placé \emph{avant} le suffixe \textbf{-ol} (ex. : propan-\textbf{2}-ol). Si un seul $-\ce{OH}$ et pas d'ambiguïté, le locant peut être omis pour le propan-1-ol (souvent dit propan-1-ol ou 1-propanol, mais IUPAC recommande le locant).
\item \textbf{Nommer et numéroter les substituants :} groupes alkyles (méthyle, éthyle...), halogènes, etc., par ordre alphabétique, avec leurs locants.
\item \textbf{Assembler :} \emph{locants-substituants-radical-\textbf{locant}-ol}.
\end{enumerate}
\end{methode}

\begin{exemplebox}[Exemples de nomenclature]
\begin{enumerate}
\item \chemfig{CH_3-CH_2-CH_2-OH} : chaîne 3 C (propane), OH sur C1 $\Rightarrow$ \textbf{propan-1-ol}.
\item \chemfig{CH_3-CH(-[2]OH)-CH_3} : chaîne 3 C, OH sur C2 $\Rightarrow$ \textbf{propan-2-ol} (ancien : isopropanol).
\item \chemfig{CH_3-CH_2-CH(-[2]CH_3)-CH_2-OH} : chaîne 4 C (butane), OH sur C1, méthyle sur C3 $\Rightarrow$ \textbf{3-méthylbutan-1-ol}.
\item \chemfig{HO-CH_2-CH_2-OH} : chaîne 2 C, deux OH $\Rightarrow$ \textbf{éthane-1,2-diol} (glycol). Le suffixe devient \textbf{-diol}, \textbf{-triol}, etc., et on conserve le 'e' de l'alcane (éthane, propane).
\end{enumerate}
\end{exemplebox}

\begin{attention}[Pièges fréquents en nomenclature]
\begin{itemize}
\item \textbf{Oublier le locant de l'OH :} écrire « butanol » au lieu de « butan-1-ol » ou « butan-2-ol » est une erreur (sauf pour le méthanol et l'éthanol où l'OH est nécessairement en 1).
\item \textbf{Mauvaise numérotation :} \chemfig{CH_3-CH_2-CH(OH)-CH_3} se numérote depuis la droite (OH sur C2) $\Rightarrow$ butan-2-ol, \emph{pas} butan-3-ol.
\item \textbf{Chaîne principale mal choisie :} pour \chemfig{CH_3-CH_2-CH(CH_3)-CH_2-OH}, la chaîne la plus longue portant l'OH a 4 C (butane), pas 3 C (propane). C'est le 3-méthylbutan-1-ol.
\item \textbf{Ordre alphabétique :} « 3-éthyl-2-méthyl... » (é avant m), pas l'inverse.
\end{itemize}
\end{attention}

\courssub{Les trois classes d'alcools : primaire, secondaire, tertiaire}

La classification dépend du nombre d'atomes de \textbf{carbone} liés au carbone porteur du groupe $-\ce{OH}$ (carbone \textbf{$\alpha$} ou carbone \emph{carbinolique}). C'est une notion cruciale pour prévoir la réactivité (oxydation, déshydratation, substitution).

\begin{definition}[Classes d'alcools]
Soit $C_\alpha$ le carbone portant le groupe $-\ce{OH}$.
\begin{itemize}
\item \textbf{Alcool primaire ($1^\circ$) :} $C_\alpha$ est lié à \textbf{un seul} atome de carbone (et deux H). Groupe $-\mathrm{CH_2OH}$.
\item \textbf{Alcool secondaire ($2^\circ$) :} $C_\alpha$ est lié à \textbf{deux} atomes de carbone (et un H). Groupe $-\mathrm{CH(OH)}-$.
\item \textbf{Alcool tertiaire ($3^\circ$) :} $C_\alpha$ est lié à \textbf{trois} atomes de carbone (pas de H). Groupe $-\mathrm{C(OH)(R)(R')}-$.
\end{itemize}
\end{definition}

\begin{popfigure}
\begin{center}
\begin{tikzpicture}[scale=0.9, every node/.style={transform shape}]
% Primaire
\node (p1) at (0,0) {\chemfig{[:30]R-[:-30]CH_2-[:-90]OH}};
\node[below=0.5cm of p1, align=center] {\textbf{Primaire ($1^\circ$)}\\ \chemfig{R-CH_2-OH}\\ $C_\alpha$ lié à 1 C};
% Secondaire
\node (p2) at (5,0) {\chemfig{[:30]R-[:-30]CH(-[:-90]OH)-[:-30]R'}};
\node[below=0.5cm of p2, align=center] {\textbf{Secondaire ($2^\circ$)}\\ \chemfig{R-CH(OH)-R'}\\ $C_\alpha$ lié à 2 C};
% Tertiaire
\node (p3) at (10,0) {\chemfig{[:30]R-[:-30]C(-[:-90]OH)(-[:90]R'')-[:-30]R'}};
\node[below=0.5cm of p3, align=center] {\textbf{Tertiaire ($3^\circ$)}\\ \chemfig{R-C(OH)(R'')-R'}\\ $C_\alpha$ lié à 3 C};
\end{tikzpicture}
\end{center}
\legende{Représentation de Cram des trois classes d'alcools. Le carbone $\alpha$ (porteur de l'OH) est en gras. Seule la nature des groupes R change.}
\end{popfigure}

\begin{exemplebox}[Classification d'alcools isomères \ce{C4H10O}]
Les 4 alcools isomères de formule brute \ce{C4H10O} :
\begin{enumerate}
\item \chemfig{CH_3-CH_2-CH_2-CH_2-OH} : Butan-1-ol $\rightarrow$ \textbf{Primaire} ($C_\alpha$ = CH2).
\item \chemfig{CH_3-CH_2-CH(-[2]OH)-CH_3} : Butan-2-ol $\rightarrow$ \textbf{Secondaire} ($C_\alpha$ = CH).
\item \chemfig{CH_3-CH(-[2]CH_3)-CH_2-OH} : 2-Méthylpropan-1-ol $\rightarrow$ \textbf{Primaire} ($C_\alpha$ = CH2).
\item \chemfig{CH_3-C(-[2]OH)(-[:-30]CH_3)-[:-90]CH_3} : 2-Méthylpropan-2-ol $\rightarrow$ \textbf{Tertiaire} ($C_\alpha$ = C quaternaire).
\end{enumerate}
\end{exemplebox}

\begin{methode}[Déterminer la classe d'un alcool à partir de sa formule]
\begin{enumerate}
\item Identifier le carbone porteur du $-\ce{OH}$ (carbone $\alpha$).
\item Compter le nombre d'atomes de \textbf{carbone} directement liés à ce carbone $\alpha$ (ne pas compter l'H ni l'O).
\item \textbf{1 carbone voisin} $\Rightarrow$ Primaire ; \textbf{2} $\Rightarrow$ Secondaire ; \textbf{3} $\Rightarrow$ Tertiaire.
\end{enumerate}
\end{methode}

\courssub{Les polyalcools (diols, triols) : glycol et glycérol}

Lorsque la molécule porte plusieurs groupes hydroxyles, on parle de \cle{polyalcools} (ou polyols). Les deux exemples majeurs au programme sont :

\begin{definition}[Glycol et Glycérol]
\begin{itemize}
\item \textbf{Éthane-1,2-diol} (nom usuel : \cle{glycol}) : \ce{HO-CH2-CH2-OH}. Diol vicinal (deux OH sur deux carbones adjacents). Liquide visqueux, doux, \textbf{toxique} (intoxication mortelle par ingestion, métabolisé en acide oxalique). Utilisé comme \textbf{antigel} dans les radiateurs de voitures (baisse le point de solidification de l'eau).
\item \textbf{Propane-1,2,3-triol} (nom usuel : \cle{glycérol} ou glycérine) : \ce{HO-CH2-CH(OH)-CH2-OH}. Triol. Liquide très visqueux, hygroscopique (absorbe l'humidité de l'air), doux, \textbf{non toxique}. Produit lors de la saponification (fabrication du savon à partir d'huile de palme chez les savonneries artisanales ou industrielles comme \emph{Savonnerie du Cameroun}). Utilisé en pharmacie, cosmétique (hydratant), alimentaire (E422).
\end{itemize}
\end{definition}

\begin{propriete}[Propriétés physiques des polyalcools]
La multiplicité des groupes $-\ce{OH}$ renforce considérablement les liaisons hydrogène intermoléculaires par rapport aux mono-alcools de masse molaire voisine :
\begin{itemize}
\item Températures d'ébullition très élevées (glycol : \SI{197}{\celsius} ; glycérol : \SI{290}{\celsius} avec décomposition).
\item Miscibilité totale avec l'eau (nombreux sites donneurs/accepteurs de liaisons H).
\item Viscosité forte (glycérol) due à l'enchevêtrement du réseau de liaisons hydrogène.
\end{itemize}
\end{propriete}

\begin{aretenir}[L'essentiel de la section 2]
\begin{itemize}
\item Alcool saturé acyclique : formule brute $C_nH_{2n+2}O$ (isomères des éthers).
\item Nomenclature IUPAC : chaîne la plus longue portant l'OH, numérotation donnant le plus bas locant à l'OH, suffixe \textbf{-ol} (ou -diol, -triol).
\item Trois classes selon le nombre de carbones liés au $C_\alpha$ : \textbf{Primaire} ($\ce{-CH2OH}$), \textbf{Secondaire} ($\ce{-CHOH-}$), \textbf{Tertiaire} ($\ce{-C(OH)R2}$).
\item Polyalcools : \textbf{Glycol} (éthane-1,2-diol, antigel, toxique) et \textbf{Glycérol} (propane-1,2,3-triol, savonnerie, cosmétique, non toxique).
\end{itemize}
\end{aretenir}

% ============================================================
% SECTION 3 : PROPRIÉTÉS PHYSIQUES DES ALCOOLS ET POLYALCOOLS
% ============================================================
\coursec{Propriétés physiques des alcools et polyalcools}

\courssub{Influence du groupe hydroxyle : liaisons hydrogène}

Le groupe $-\ce{OH}$ est \textbf{polaire} (différence d'électronégativité O=3,5, H=2,1). L'atome d'oxygène porte une charge partielle négative ($\delta^-$), l'hydrogène une charge partielle positive ($\delta^+$). Cela permet la formation de \cle{liaisons hydrogène} intermoléculaires : l'H $\delta^+$ d'une molécule est attiré par l'O $\delta^-$ d'une molécule voisine.

\begin{propriete}[Conséquences macroscopiques des liaisons hydrogène]
Comparés aux alcanes de masse molaire voisine (qui n'ont que des forces de Van der Waals faibles) :
\begin{itemize}
\item \textbf{Températures d'ébullition ($T_{eb}$) et de fusion ($T_f$) beaucoup plus élevées.}
\item \textbf{Solubilité dans l'eau} favorisée pour les chaînes courtes (l'alcool peut s'insérer dans le réseau H de l'eau).
\end{itemize}
\end{propriete}

\begin{popfigure}
\begin{center}
\begin{tikzpicture}
\begin{axis}[
    width=\linewidth, height=8cm,
    xlabel={Nombre d'atomes de carbone $n$},
    ylabel={Température d'ébullition (\si{\celsius})},
    grid=major, legend pos=north west,
    xmin=0, xmax=6, ymin=-100, ymax=150,
    tick label style={font=\small}, label style={font=\small}, legend style={font=\small}
]
\addplot[popcurve, mark=*, mark size=2, popblue] coordinates {
(1, -161) (2, -89) (3, -42) (4, 0) (5, 36)
};
\addlegendentry{Alcanes $C_nH_{2n+2}$}
\addplot[popcurve, mark=square*, mark size=2, poporange] coordinates {
(1, 65) (2, 78) (3, 97) (4, 118) (5, 138)
};
\addlegendentry{Alcools $C_nH_{2n+2}O$ (linéaires)}
\addplot[popcurve, mark=triangle*, mark size=2, popgreen] coordinates {
(1, -24) (2, 35) (3, 76)
};
\addlegendentry{Éthers $C_nH_{2n+2}O$}
\end{axis}
\end{tikzpicture}
\end{center}
\legende{Évolution des températures d'ébullition. Les alcools (courbe orange) ont des $T_{eb}$ bien supérieures aux alcanes (bleu) et aux éthers (vert) de même masse molaire, preuve de la force des liaisons hydrogène. Données approximatives : Méthane -161, Éthane -89, Propane -42, Butane 0, Pentane 36 ; Méthanol 65, Éthanol 78, Propan-1-ol 97, Butan-1-ol 118, Pentan-1-ol 138 ; Méthoxyméthane -24, Éthoxyéthane 35, Méthoxypropane 76.}
\end{popfigure}

\courssub{Solubilité dans l'eau : le combat chaîne carbonée vs groupe OH}

La solubilité résulte de l'équilibre entre :
\begin{itemize}
\item La partie \textbf{hydrophile} (groupe $-\ce{OH}$) qui « aime » l'eau (liaisons H donneur/accepteur).
\item La partie \textbf{hydrophobe} (chaîne carbonée $R$) qui « fuit » l'eau (effet hydrophobe, rupture du réseau H de l'eau).
\end{itemize}

\begin{propriete}[Règle empirique de solubilité des alcools dans l'eau à \SI{20}{\celsius}]
\begin{itemize}
\item $n \le 3$ (Méthanol, Éthanol, Propan-1-ol, Propan-2-ol) : \textbf{Miscibles} en toutes proportions.
\item $n = 4$ (Butan-1-ol) : \textbf{Faiblement soluble} (\approx \SI{7,5}{g/100 mL}). Deux phases au-delà.
\item $n = 5$ (Pentan-1-ol) : \textbf{Très peu soluble} (\approx \SI{2,2}{g/100 mL}).
\item $n \ge 6$ : \textbf{Pratiquement insolubles} (couche séparée).
\end{itemize}
\end{propriete}

\begin{attention}[Isomérie et solubilité]
Pour une même formule brute, les alcools \textbf{ramifiés} sont plus solubles que les alcools \textbf{linéaires} (la chaîne ramifiée est plus compacte, moins hydrophobe). Exemple : le 2-méthylpropan-2-ol (tert-butanol, $n=4$ ramifié) est miscible à l'eau, alors que le butan-1-ol (linéaire) ne l'est pas.
\end{attention}

\courssub{Propriétés physiques des polyalcools : glycol et glycérol}

Le tableau ci-dessous compare les propriétés physiques avec des alcools mono-fonctionnés de masse molaire proche.

\begin{center}
\begin{tabularx}{\linewidth}{|l|c|c|c|c|}
\hline
\rowcolor{popblueL}
\textbf{Composé} & \textbf{Formule brute} & \textbf{Masse molaire (g/mol)} & \textbf{$T_{eb}$ (\si{\celsius})} & \textbf{Solubilité eau} \\
\hline
Propan-1-ol & \ce{C3H8O} & 60 & 97 & Miscible \\
\hline
\textbf{Glycol (Éthane-1,2-diol)} & \ce{C2H6O2} & 62 & \textbf{197} & \textbf{Miscible} \\
\hline
Butan-1-ol & \ce{C4H10O} & 74 & 118 & 7,5 g/100mL \\
\hline
\textbf{Glycérol (Propane-1,2,3-triol)} & \ce{C3H8O3} & 92 & \textbf{290 (dec.)} & \textbf{Miscible} \\
\hline
\end{tabularx}
\end{center}

\begin{savaistu}[Le glycol, tueur silencieux des chiens... et des radiateurs]
Le glycol (éthane-1,2-diol) a un goût **sucré** qui attire les animaux (chiens, chats) s'il fuit d'un radiateur ou d'un bidon mal fermé. Son ingestion provoque une insuffisance rénale aiguë mortelle par formation de cristaux d'oxalate de calcium dans les tubules rénaux (métabolisation en acide glycolique puis oxalique). Au Cameroun, où les véhicules anciens sont nombreux et les fuites de radiateur fréquentes, c'est un risque réel. L'antidote est l'éthanol (perfusion IV) qui sature l'enzyme (alcool déshydrogénase) et empêche la transformation toxique. Le glycérol, lui, est inoffensif et utilisé dans les sirops pour la toux (ex. : sirop de bissap glycériné) pour son effet émollient et sucrant.
\end{savaistu}

\begin{aretenir}[L'essentiel de la section 3]
\begin{itemize}
\item Le groupe $-\ce{OH}$ crée des \textbf{liaisons hydrogène} intermoléculaires fortes.
\item $T_{eb}$ et $T_f$ des alcools $\gg$ alcanes/éthers de même masse molaire.
\item Solubilité dans l'eau : miscibles pour $n \le 3$, décroît rapidement pour $n \ge 4$ (partie hydrophobe dominante).
\item Polyalcools (glycol, glycérol) : $T_{eb}$ très hautes, miscibilité totale, forte viscosité (multiplicité des sites H).
\end{itemize}
\end{aretenir}

% ============================================================
% SECTION 4 : ALDÉHYDES ET CÉTONES : FORMULES ET NOMENCLATURE
% ============================================================
\coursec{Aldéhydes et cétones : formules et nomenclature}

\courssub{Formules générales et brutes}

Les aldéhydes et cétones saturés acycliques (un seul carbonyle, pas de C=C) sont des isomères de fonction. Leur formule brute commune est :

\[ \textbf{$C_nH_{2n}O$} \quad (n \ge 1 \text{ pour aldéhydes, } n \ge 3 \text{ pour cétones}) \]

\begin{exemplebox}[Vérification]
\begin{itemize}
\item Méthanal \ce{H-CHO} : $n=1 \Rightarrow \ce{CH2O}$.
\item Éthanal \ce{CH3-CHO} : $n=2 \Rightarrow \ce{C2H4O}$.
\item Propanal \ce{CH3-CH2-CHO} : $n=3 \Rightarrow \ce{C3H6O}$.
\item Propanone \ce{CH3-CO-CH3} : $n=3 \Rightarrow \ce{C3H6O}$ (isomère du propanal).
\item Butan-2-one \ce{CH3-CO-CH2-CH3} : $n=4 \Rightarrow \ce{C4H8O}$.
\end{itemize}
\end{exemplebox}

\courssub{Nomenclature systématique (IUPAC) des aldéhydes}

Le groupe $-\ce{CHO}$ est **toujours en bout de chaîne** (carbone n°1 par définition). Le suffixe est \textbf{-al}.

\begin{methode}[Nommer un aldéhyde saturé acyclique]
\begin{enumerate}
\item Choisir la chaîne la plus longue **contenant le carbone du carbonyle** (celui-ci est C1).
\item Le radical est celui de l'alcane correspondant (méth-, éth-, prop-, but-...).
\item Le suffixe est \textbf{-al} (remplace le -e final de l'alcane : méthane $\rightarrow$ méthanal).
\item Numéroter les substituants sur la chaîne à partir du carbonyle (C1).
\item Assembler : \emph{locants-substituants-radical-\textbf{al}}.
\end{enumerate}
\end{methode}

\begin{exemplebox}[Exemples d'aldéhydes]
\begin{enumerate}
\item \ce{H-CHO} : \textbf{Méthanal} (formol).
\item \ce{CH3-CHO} : \textbf{Éthanal}.
\item \ce{CH3-CH2-CHO} : \textbf{Propanal}.
\item \ce{CH3-CH(CH3)-CHO} : chaîne 3 C (propanal), méthyle sur C2 $\Rightarrow$ \textbf{2-Méthylpropanal}.
\item \ce{OHC-CH2-CH2-CHO} : 2 groupes $-\ce{CHO}$ $\Rightarrow$ \textbf{Butane-1,4-dial} (suffix -dial).
\end{enumerate}
\end{exemplebox}

\courssub{Nomenclature systématique (IUPAC) des cétones}

Le carbonyle est **à l'intérieur de la chaîne**. Le suffixe est \textbf{-one}.

\begin{methode}[Nommer une cétone saturée acyclique]
\begin{enumerate}
\item Choisir la chaîne la plus longue **contenant le carbone du carbonyle**.
\item Numéroter la chaîne de façon à donner le **plus bas locant possible au carbonyle** (règle du plus bas locant pour la fonction).
\item Le radical est celui de l'alcane correspondant.
\item Indiquer la position du carbonyle par son locant placé avant le suffixe \textbf{-one} (ex. : butan-\textbf{2}-one).
\item Nommer et numéroter les substituants.
\item Assembler : \emph{locants-substituants-radical-\textbf{locant}-one}.
\end{enumerate}
\end{methode}

\begin{exemplebox}[Exemples de cétones]
\begin{enumerate}
\item \ce{CH3-CO-CH3} : chaîne 3 C, CO sur C2 $\Rightarrow$ \textbf{Propanone} (acétone). Locant 2 obligatoire (seule place possible).
\item \ce{CH3-CO-CH2-CH3} : chaîne 4 C, CO sur C2 $\Rightarrow$ \textbf{Butan-2-one} (méthyléthylcétone).
\item \ce{CH3-CH2-CO-CH2-CH3} : chaîne 5 C, CO sur C3 $\Rightarrow$ \textbf{Pentan-3-one} (diéthylcétone).
\item \ce{CH3-CO-CH(CH3)2} : chaîne 4 C (butane), CO sur C2, méthyle sur C3 $\Rightarrow$ \textbf{3-Méthylbutan-2-one}.
\end{enumerate}
\end{exemplebox}

\begin{attention}[Pièges nomenclature aldéhydes/cétones]
\begin{itemize}
\item \textbf{Aldéhyde :} le carbonyle est **toujours en C1**. On n'écrit **jamais** le locant « 1 » devant -al (ex: \textbf{propanal}, pas propan-1-al).
\item \textbf{Cétone :} le locant du carbonyle est **obligatoire** (sauf ambiguïté impossible comme propanone). Ex: butan-2-one, pas butanone.
\item \textbf{Priorité fonctionnelle :} Si une molécule porte à la fois un $-\ce{OH}$ (alcool) et un $-\ce{CHO}$ (aldéhyde), l'aldéhyde est prioritaire $\Rightarrow$ suffixe \textbf{-al}, l'alcool devient préfixe \textbf{hydroxy-}. Ex: \ce{HO-CH2-CHO} = 2-hydroxyéthanal (glycolaldéhyde).
\end{itemize}
\end{attention}

\courssub{Isomérie dans la famille des aldéhydes et cétones}

Trois types d'isomérie sont possibles pour la formule $C_nH_{2n}O$ :
\begin{enumerate}
\item \textbf{Isomérie de chaîne :} squelette carboné différent. Ex: Butanal (linéaire) vs 2-Méthylpropanal (ramifié).
\item \textbf{Isomérie de position (cétones seulement) :} position du carbonyle sur la chaîne. Ex: Butan-2-one vs Butan-3-one (symétrique, identique à butan-2-one) ; Pentan-2-one vs Pentan-3-one.
\item \textbf{Isomérie de fonction :} Aldéhyde vs Cétone (même $n$). Ex: Propanal vs Propanone (\ce{C3H6O}) ; Butanal vs Butan-2-one (\ce{C4H8O}).
\end{enumerate}

\begin{exoresolu}[Nomenclature et isomérie pour \ce{C4H8O}]
\textbf{Énoncé.} Écrire les formules semi-développées et donner les noms IUPAC de tous les aldéhydes et cétones saturés acycliques de formule brute \ce{C4H8O}. Les classer par type d'isomérie.
\tcblower
\textbf{Solution.}
\textbf{Aldéhydes (carbonyle en C1) :}
\begin{enumerate}
\item \chemfig{CH_3-CH_2-CH_2-CHO} : \textbf{Butanal} (chaîne linéaire).
\item \chemfig{CH_3-CH(-[2]CH_3)-CHO} : \textbf{2-Méthylpropanal} (chaîne ramifiée).
\end{enumerate}
$\rightarrow$ Isomérie de **chaîne** entre 1 et 2.

\textbf{Cétones (carbonyle interne) :}
\begin{enumerate}
\setcounter{enumi}{2}
\item \chemfig{CH_3-CO-CH_2-CH_3} : \textbf{Butan-2-one} (ou méthyléthylcétone).
\item \chemfig{CH_3-CH_2-CO-CH_3} : Identique à 3 (symétrie), pas d'isomère butan-3-one distinct.
\end{enumerate}
Pas d'isomérie de position pour $n=4$ (seule position 2 possible).

\textbf{Isomérie de fonction :} Butanal (1) et 2-Méthylpropanal (2) sont isomères de fonction de Butan-2-one (3).
\end{exoresolu}

\begin{aretenir}[L'essentiel de la section 4]
\begin{itemize}
\item Aldéhydes et cétones saturés acycliques : formule brute commune $C_nH_{2n}O$.
\item Aldéhyde : suffixe \textbf{-al}, carbonyle en C1 (locant jamais écrit). Ex: Éthanal, Propanal.
\item Cétone : suffixe \textbf{-one}, carbonyle interne (locant obligatoire). Ex: Propanone, Butan-2-one.
\item Isoméries : chaîne, position (cétones), fonction (aldéhyde/cétone).
\end{itemize}
\end{aretenir}

% ============================================================
% SECTION 5 : TESTS D'IDENTIFICATION DES ALDÉHYDES ET CÉTONES
% ============================================================
\coursec{Tests d'identification des aldéhydes et cétones}

\courssub{Problématique : distinguer aldéhydes et cétones}

Les aldéhydes et les cétones partagent le groupe carbonyle $\ce{C=O}$. Ils donnent **tous les deux** le test de la 2,4-DNPH (précipité jaune/orange). Pour les distinguer, il faut exploiter la différence fondamentale : l'aldéhyde a un **H sur le carbone carbonyle** ($R-\ce{CHO}$), ce qui le rend **oxydable** (réducteur), alors que la cétone n'a pas cet H ($R-\ce{CO}-R'$) et résiste à l'oxydation ménagée.

\courssub{Test commun : la 2,4-DNPH (réactif de Brady)}

\begin{experience}[Test à la 2,4-dinitrophénylhydrazine (2,4-DNPH)]
\textbf{Dispositif :} Tube à essai, quelques gouttes de composé à tester, \SI{2}{mL} de solution de 2,4-DNPH dans l'éthanol acide (réactif de Brady).
\textbf{Protocole :} Mélanger, chauffer légèrement si nécessaire, laisser refroidir.
\textbf{Observation :}
\begin{itemize}
\item \textbf{Précipité cristallisé jaune-orangé à rouge} (hydrazone) $\Rightarrow$ **Présence d'un groupe carbonyle** (aldéhyde \textbf{OU} cétone).
\item Pas de précipité $\Rightarrow$ Pas de carbonyle (ou stéricité trop forte).
\end{itemize}
\textbf{Équation-bilan (générale) :}
\[ R_2C=O + H_2N-NH-C_6H_3(NO_2)_2 \rightarrow R_2C=N-NH-C_6H_3(NO_2)_2 \downarrow + H_2O \]
($R$ ou $H$ pour aldéhyde ; $R, R'$ pour cétone).
\end{experience}

\begin{attention}[2,4-DNPH ne distingue PAS aldéhyde de cétone]
Ce test est **positif pour les deux familles**. Il sert à *détecter* la fonction carbonyle, pas à la *classer*. Pour distinguer, il faut les tests d'oxydation (Fehling, Tollens, Schiff).
\end{attention}

\courssub{Test au réactif de Schiff (détection des aldéhydes)}

Le réactif de Schiff est une solution de fuchsine décolorée par le $\ce{SO2}$ (incolore). Les aldéhydes restaurent la couleur rose/magenta de la fuchsine par addition sur le carbonyle. Les cétones ne le font pas (ou très lentement/faiblement pour les cétones $\alpha$-hydroxylées, hors programme).

\begin{experience}[Test au réactif de Schiff]
\textbf{Dispositif :} Tube à essai, \SI{1}{mL} de réactif de Schiff + quelques gouttes de composé.
\textbf{Observation :}
\begin{itemize}
\item \textbf{Coloration rose vif / magenta} (apparition rapide, \textless 1 min) $\Rightarrow$ \textbf{Aldéhyde}.
\item \textbf{Reste incolore} (ou trouble jaunâtre lent) $\Rightarrow$ \textbf{Cétone} (négatif).
\end{itemize}
\textbf{Interprétation :} L'aldéhyde réagit avec le groupement sulfite de la fuchsine décolorée pour reformer le système conjugué coloré.
\end{experience}

\courssub{Test à la liqueur de Fehling (oxydation par \ce{Cu^2+})}

La liqueur de Fehling contient des ions \ce{Cu^2+} (bleu) complexés par l'ion tartrate en milieu basique (soude). L'aldéhyde réduit \ce{Cu^2+} en \ce{Cu+} (précipité rouge brique \ce{Cu2O}). La cétone ne réagit pas.

\begin{experience}[Test à la liqueur de Fehling]
\textbf{Dispositif :} Tube à essai, \SI{1}{mL} Fehling A (\ce{CuSO4}) + \SI{1}{mL} Fehling B (tartrate + \ce{NaOH}) $\rightarrow$ solution bleue homogène. Ajouter quelques gouttes de composé. Chauffer au bain-marie \SI{60}{\celsius} pendant quelques minutes.
\textbf{Observation :}
\begin{itemize}
\item \textbf{Précipité rouge brique} (\ce{Cu2O}, oxyde de cuivre(I)) $\Rightarrow$ \textbf{Aldéhyde}.
\item \textbf{Solution reste bleue} (pas de précipité) $\Rightarrow$ \textbf{Cétone}.
\end{itemize}
\textbf{Équation-bilan (milieu basique) :}
\[ R-CHO + 2 Cu^{2+} + 4 OH^- \rightarrow R-COO^- + Cu_2O \downarrow \text{ (rouge) } + 3 H_2O \]
\textit{Note :} L'aldéhyde est oxydé en ion carboxylate (\ce{R-COO-}), le cuivre(II) est réduit en cuivre(I) insoluble.
\end{experience}

\begin{attention}[Fehling : conditions strictes]
\begin{itemize}
\item Milieu **basique** indispensable (tartrate empêche la précipitation de \ce{Cu(OH)2} bleu pâle).
\item **Chauffer** (cinétique lente à froid).
\item Ne pas confondre le précipité bleu pâle de \ce{Cu(OH)2} (si pas assez de tartrate ou trop de soude) avec le précipité rouge brique \ce{Cu2O} caractéristique.
\end{itemize}
\end{attention}

\courssub{Test au réactif de Tollens (oxydation par \ce{Ag+}) : le miroir d'argent}

Le réactif de Tollens contient l'ion diamminargent(I) \ce{[Ag(NH3)2]+} (solution d'\ce{AgNO3} + \ce{NH3} jusqu'à redissolution du précipité \ce{AgOH}). L'aldéhyde réduit \ce{Ag+} en \ce{Ag} métallique (miroir ou précipité noir). La cétone ne réagit pas.

\begin{experience}[Test au réactif de Tollens]
\textbf{Dispositif :} Tube à essai **propre** (rincé à l'acide nitrique dilué puis eau, sinon miroir parasite), \SI{2}{mL} réactif de Tollens frais + quelques gouttes de composé. Chauffer **doucement** au bain-marie (\SI{35}{\celsius} max, pas de flamme directe).
\textbf{Observation :}
\begin{itemize}
\item \textbf{Miroir d'argent} (paroi interne argentée) ou **précipité noir** (argent colloïdal) $\Rightarrow$ \textbf{Aldéhyde}.
\item \textbf{Solution reste incolore} $\Rightarrow$ \textbf{Cétone}.
\end{itemize}
\textbf{Équation-bilan (milieu basique/ammoniacal) :}
\[ R-CHO + 2 [Ag(NH_3)_2]^+ + 3 OH^- \rightarrow R-COO^- + 2 Ag \downarrow \text{ (miroir) } + 4 NH_3 + 2 H_2O \]
\textit{Note :} L'argent métallique se dépose sur le verre (miroir) si le tube est très propre, sinon en poudre noire.
\end{experience}

\begin{attention}[Sécurité Tollens : danger d'explosion !]
Le réactif de Tollens **doit être préparé frais** et **détruit immédiatement après usage** en le diluant et en ajoutant de l'acide nitrique dilué ou de l'eau de Javel. S'il est stocké, il forme du fulminate d'argent \ce{Ag-C#N} (ou azoture d'argent), **hautement explosif** au sec ou au choc. \textbf{Interdiction formelle de garder un tube de Tollens usagé.}
\end{attention}

\courssub{Tableau récapitulatif des tests de distinction}

\begin{center}
\begin{tabularx}{\linewidth}{|l|X|X|}
\hline
\rowcolor{popblueL}
\textbf{Test} & \textbf{Aldéhyde ($R-\ce{CHO}$)} & \textbf{Cétone ($R-\ce{CO}-R'$)} \\
\hline
\textbf{2,4-DNPH (Brady)} & Précipité jaune-orangé \textbf{(+)} & Précipité jaune-orangé \textbf{(+)} \\
\hline
\textbf{Réactif de Schiff} & \textbf{Rose magenta} \textbf{(+)} & Incolore \textbf{(-)} \\
\hline
\textbf{Liqueur de Fehling} & Précipité \textbf{rouge brique} \ce{Cu2O} \textbf{(+)} & Bleu, pas de ppt \textbf{(-)} \\
\hline
\textbf{Réactif de Tollens} & \textbf{Miroir d'argent} \ce{Ag} \textbf{(+)} & Incolore \textbf{(-)} \\
\hline
\end{tabularx}
\end{center}

\begin{methode}[Stratégie d'identification complète d'un carbonylé inconnu]
\begin{enumerate}
\item Faire le test **2,4-DNPH**. Si négatif $\rightarrow$ pas de carbonyle (autre famille). Si positif $\rightarrow$ aldéhyde **OU** cétone.
\item Faire un test d'**oxydation** (Fehling **OU** Tollens **OU** Schiff).
      \begin{itemize}
      \item Positif (rouge Fehling, miroir Tollens, rose Schiff) $\rightarrow$ \textbf{Aldéhyde}.
      \item Négatif $\rightarrow$ \textbf{Cétone}.
      \end{itemize}
\end{enumerate}
\end{methode}

\begin{exoresolu}[Identification d'un carbonylé inconnu X]
\textbf{Énoncé.} On réalise trois tests sur un composé organique liquide X de formule brute \ce{C3H6O} :
\begin{itemize}
\item Test 2,4-DNPH : précipité jaune-orangé.
\item Test de Fehling (chauffé) : solution reste bleue.
\item Test de Tollens : solution reste incolore.
\end{itemize}
Déterminer la famille de X et proposer son nom et sa formule semi-développée.
\tcblower
\textbf{Solution.}
\begin{enumerate}
\item 2,4-DNPH positif $\Rightarrow$ X possède un groupe carbonyle $\Rightarrow$ Aldéhyde **ou** Cétone.
\item Fehling négatif (pas de \ce{Cu2O} rouge) $\Rightarrow$ X n'est **pas** oxydé par \ce{Cu^2+} basique $\Rightarrow$ Pas d'aldéhyde.
\item Tollens négatif (pas de miroir) $\Rightarrow$ Confirme l'absence d'aldéhyde.
\item Conclusion : X est une \textbf{cétone}.
\item Formule brute \ce{C3H6O} ($n=3$). Seule cétone possible : \textbf{Propanone} (acétone).
\item Formule semi-développée : \ce{CH3-CO-CH3}.
\end{enumerate}
\end{exoresolu}

\begin{aretenir}[L'essentiel de la section 5]
\begin{itemize}
\item \textbf{2,4-DNPH :} détecte **tout** carbonyle (aldéhyde + cétone). Précipité jaune-orangé.
\item \textbf{Schiff / Fehling / Tollens :} tests d'**oxydation** (réduction du réactif par l'aldéhyde).
      \begin{itemize}
      \item \textbf{Aldéhyde :} Positif (Rose / Rouge brique \ce{Cu2O} / Miroir \ce{Ag}).
      \item \textbf{Cétone :} Négatif (Incolore / Bleu / Incolore).
      \end{itemize}
\item Équations-bilans Fehling et Tollens à savoir écrire (oxydation aldéhyde $\rightarrow$ carboxylate).
\item Sécurité : Tollens frais, détruit après usage (risque explosion).
\end{itemize}
\end{aretenir}

% ============================================================
% SECTION 6 : SYNTHÈSE ET APPLICATIONS AU QUOTIDIEN
% ============================================================
\coursec{Synthèse et applications au quotidien}

\courssub{Tableau de synthèse global des familles oxygénées}

\begin{center}
\small
\begin{tabularx}{\linewidth}{|l|X|X|X|X|}
\hline
\rowcolor{popblueL}
\textbf{Critère} & \textbf{Alcool} & \textbf{Aldéhyde} & \textbf{Cétone} & \textbf{Éther-oxyde} \\
\hline
Groupe fonctionnel & $-\ce{OH}$ (Hydroxyle) & $-\ce{CHO}$ (Carbonyle terminal) & $-\ce{CO}-$ (Carbonyle interne) & $-\ce{O}-$ (Pont) \\
Formule générale & $R-\ce{OH}$ & $R-\ce{CHO}$ & $R-\ce{CO}-R'$ & $R-\ce{O}-R'$ \\
Formule brute (sat. acycl.) & $C_nH_{2n+2}O$ & $C_nH_{2n}O$ & $C_nH_{2n}O$ & $C_nH_{2n+2}O$ \\
Nomenclature (suffixe) & \textbf{-ol} (diol, triol) & \textbf{-al} (dial) & \textbf{-one} & \textbf{alkoxy-alcane} \\
Classes / Isomérie & Primaire, Secondaire, Tertiaire & Isomérie chaîne, fonction & Isomérie chaîne, position, fonction & Isomérie chaîne, fonction \\
\hline
\textbf{Test caractéristique} & \textbf{Na(s) : $\ce{H2}$ gazeux} & \textbf{Fehling/Tollens/Schiff : (+)} & \textbf{Fehling/Tollens/Schiff : (-)} & \textbf{Peu réactif} \\
& (Alcoolates) & \ce{Cu2O} rouge, Miroir \ce{Ag}, Rose & 2,4-DNPH : (+) & (Solvant inerte) \\
\hline
2,4-DNPH & (-) & \textbf{(+)} Jaune/Orange & \textbf{(+)} Jaune/Orange & (-) \\
\hline
Propriétés physiques & Liaisons H fortes. $T_{eb}$ haute. Sol. eau : $n\le3$ miscibles. & Polaires. $T_{eb}$ modérées. Sol. eau : petits miscibles. & Polaires. $T_{eb}$ modérées. Sol. eau : petits miscibles. & Faiblement polaires. Pas de liaison H donneur. $T_{eb}$ basses. Insolubles eau. \\
\hline
Exemple camerounais & Éthanol (Vin de palme, Bili-bili, Désinfectant, Carburant SONARA) & Méthanal (Formol, conservation labo, \textbf{interdit alimentaire}) & Propanone (Acétone, dissolvant, nettoyage industriel) & Éthoxyéthane (Ancien anesthésique, solvant labo) \\
\hline
\end{tabularx}
\end{center}

\courssub{Applications concrètes au Cameroun}

\begin{savaistu}[L'éthanol, de la tradition à l'industrie]
\begin{itemize}
\item \textbf{Vin de palme / Bili-bili / Ndoké :} Fermentation traditionnelle (levures) du sucre de sève de palmier (rafia, palmier à huile) ou de maïs/millet. Teneur en éthanol variable (3-10\% vol). Enjeu sanitaire : méthanol (toxique) parfois présent si fermentation mal maîtrisée ou addition frauduleuse.
\item \textbf{Désinfectant \SI{70}{\%} :} Concentration optimale pour dénaturer les protéines bactériennes (l'eau aide la pénétration). Produit localement par les laboratoires pharmaceutiques (ex. Cinpharm).
\item \textbf{Carburant (Supercarburant) :} La SONARA (Société Nationale de Raffinage) incorpore de l'éthanol (bioéthanol importé ou local futur) dans l'essence (E10 : 10\% éthanol) pour réduire les émissions de \ce{CO} et l'importation de pétrole.
\item \textbf{Solvant universel :} Extraction de principes actifs (pharmacopée traditionnelle : écorces de quinquina, feuilles d'artemisia), parfumerie, encres.
\end{itemize}
\end{savaistu}

\begin{savaistu}[Aldéhydes et cétones dans l'économie locale]
\begin{itemize}
\item \textbf{Méthanal (Formol) :} Au-delà des laboratoires (conservation spécimens SVT), utilisé dans la fabrication de résines (colles pour bois/contreplaqué - industrie du bois au Sud), désinfectant d'élevages (fumigation). \textbf{Surveillance sanitaire stricte} : fraudes sur poisson/viande (aspect frais artificiel).
\item \textbf{Propanone (Acétone) :} Solvant industriel majeur (peintures, vernis, colles, nettoyage pièces mécaniques). Production locale limitée, surtout importée.
\item \textbf{Aldéhydes aromatiques :} Benzaldéhyde (amande amère), vanilline (vanille), cinnamaldéhyde (cannelle) : arômes alimentaires, parfumerie. Le Cameroun produit de la vanille (région du Sud-Ouest) : extraction naturelle vs synthèse.
\end{itemize}
\end{savaistu}

\begin{savaistu}[Glycérol : le co-produit précieux de la savonnerie]
La réaction de **saponification** de l'huile de palme (triglycérides) par la soude \ce{NaOH} produit du savon (carboxylates de sodium) \textbf{ET} du glycérol (propane-1,2,3-triol).
\[ \text{TAG (huile)} + 3 \ce{NaOH} \rightarrow 3 \text{ Savon } (\ce{R-COO- Na+}) + \text{Glycérol} \]
Au Cameroun, les savonneries artisanelles (région de l'Ouest, du Littoral) et industrielles (Savonnerie du Cameroun, Azur) récupèrent ce glycérol. Il est purifié et revendu pour :
\begin{itemize}
\item Cosmétiques (crèmes hydratantes, savons surgras « surgraissés » au glycérol).
\item Pharmaceutique (sirops, suppositoires, excipient).
\item Alimentaire (humectant E422, confiserie).
\item Explosifs (nitroglycérine) — usage industriel strict.
\end{itemize}
C'est un bel exemple d'économie circulaire : déchet $\rightarrow$ ressource.
\end{savaistu}

\courssub{Exercices d'entraînement (niveau Probatoire)}

\begin{exercice}[Nomenclature et classes]
Donner le nom IUPAC et la classe (primaire, secondaire, tertiaire) des alcools suivants :
\begin{enumerate}
\item \ce{CH3-CH2-CH2-CH2-OH}
\item \ce{CH3-CH2-CH(OH)-CH3}
\item \ce{(CH3)3C-OH}
\item \ce{CH3-CH(CH3)-CH2-OH}
\item \ce{HO-CH2-CH2-CH2-CH2-OH}
\end{enumerate}
\end{exercice}

\begin{exercice}[Isomérie et formules]
\begin{enumerate}
\item Écrire les formules semi-développées et les noms des 3 aldéhydes de formule brute \ce{C5H10O}.
\item Écrire les formules semi-développées et les noms des 3 cétones de formule brute \ce{C5H10O}.
\item Parmi ces 6 composés, identifier les paires d'isomères de fonction.
\end{enumerate}
\end{exercice}

\begin{exercice}[Tests d'identification]
On dispose de trois tubes étiquetés A, B, C contenant chacun un liquide incolore : un aldéhyde, une cétone, un alcool. On réalise les tests suivants :
\begin{itemize}
\item Tube A : 2,4-DNPH $\rightarrow$ précipité jaune. Fehling $\rightarrow$ rouge brique.
\item Tube B : 2,4-DNPH $\rightarrow$ précipité jaune. Fehling $\rightarrow$ reste bleu. Tollens $\rightarrow$ incolore.
\item Tube C : 2,4-DNPH $\rightarrow$ pas de précipité. Sodium $\rightarrow$ dégagement gazeux.
\end{itemize}
Identifier le contenu de chaque tube (famille + exemple possible pour \ce{C3H6O} ou \ce{C3H8O}).
\end{exercice}

\begin{exercice}[Calcul et équation-bilan]
On fait réagir \SI{2,3}{\gram} d'éthanol avec un excès de sodium métallique.
\begin{enumerate}
\item Calculer la quantité de matière d'éthanol introduit ($M(\ce{C2H5OH}) = \SI{46}{\gram\per\mol}$).
\item Écrire l'équation-bilan de la réaction.
\item Calculer le volume de dihydrogène dégagé aux CNTP ($V_m = \SI{24}{\liter\per\mol}$).
\end{enumerate}
\end{exercice}

\begin{corrige}[Exercice 1 : Nomenclature et classes]
\begin{enumerate}
\item Butan-1-ol (Primaire : $\ce{-CH2OH}$)
\item Butan-2-ol (Secondaire : $\ce{-CHOH-}$)
\item 2-Méthylpropan-2-ol (Tertiaire : $\ce{-C(OH)(CH3)2}$)
\item 2-Méthylpropan-1-ol (Primaire : $\ce{-CH2OH}$ sur chaîne ramifiée)
\item Butane-1,4-diol (Diol primaire, deux fonctions primaires)
\end{enumerate}
\end{corrige}

\begin{corrige}[Exercice 2 : Isomérie \ce{C5H10O}]
\textbf{Aldéhydes (C1=CHO) :}
1. Pentanal \ce{CH3-CH2-CH2-CH2-CHO} (linéaire).
2. 2-Méthylbutanal \ce{CH3-CH2-CH(CH3)-CHO}.
3. 3-Méthylbutanal \ce{(CH3)2CH-CH2-CHO}.
4. 2,2-Diméthylpropanal \ce{(CH3)3C-CHO} (Pivaldehyde).
$\rightarrow$ 4 aldéhydes (isomérie de chaîne).

\textbf{Cétones (CO interne) :}
1. Pentan-2-one \ce{CH3-CO-CH2-CH2-CH3}.
2. Pentan-3-one \ce{CH3-CH2-CO-CH2-CH3} (symétrique).
3. 3-Méthylbutan-2-one \ce{CH3-CO-CH(CH3)2}.
$\rightarrow$ 3 cétones (isomérie de chaîne + position).

\textbf{Isomères de fonction :} Chaque aldéhyde est isomère de fonction de chaque cétone(12 paires au total).
\end{corrige}

\begin{corrige}[Exercice 3 : Tests identification]
\begin{itemize}
\item \textbf{Tube A :} 2,4-DNPH (+) + Fehling (+) $\Rightarrow$ \textbf{Aldéhyde}. Exemple \ce{C3H6O} : Propanal \ce{CH3CH2CHO}.
\item \textbf{Tube B :} 2,4-DNPH (+) + Fehling (-) + Tollens (-) $\Rightarrow$ \textbf{Cétone}. Exemple \ce{C3H6O} : Propanone \ce{CH3COCH3}.
\item \textbf{Tube C :} 2,4-DNPH (-) + Na (+) $\Rightarrow$ \textbf{Alcool}. Exemple \ce{C3H8O} : Propan-1-ol \ce{CH3CH2CH2OH} (ou Propan-2-ol).
\end{itemize}
\end{corrige}

\begin{corrige}[Exercice 4 : Calcul réaction Na/Éthanol]
\begin{enumerate}
\item $n(\ce{C2H5OH}) = \frac{m}{M} = \frac{2,3}{46} = \SI{0,050}{mol}$.
\item $\ce{2 C2H5OH + 2 Na -> 2 C2H5ONa + H2}$.
\item D'après l'équation : $2$ mol éthanol $\rightarrow$ $1$ mol $\ce{H2}$.
      $n(\ce{H2}) = \frac{n(\ce{C2H5OH})}{2} = \frac{0,050}{2} = \SI{0,025}{mol}$.
      $V(\ce{H2}) = n \times V_m = 0,025 \times 24 = \SI{0,60}{L} = \SI{600}{mL}$.
\end{enumerate}
\end{corrige}

\begin{aretenir}[Synthèse finale de la leçon 03 - Composés oxygénés]
\begin{itemize}
\item \textbf{4 familles :} Alcool ($-\ce{OH}$), Aldéhyde ($-\ce{CHO}$), Cétone ($-\ce{CO}-$), Éther-oxyde ($-\ce{O}-$).
\item \textbf{Diagnostic visuel :} Oxygène + environnement (H ? $\ce{C=O}$ terminal/interne ? Pont ?).
\item \textbf{Nomenclature IUPAC :} Suffixe prioritaire (-ol, -al, -one), numérotation plus bas locant pour la fonction.
\item \textbf{Alcools :} 3 classes (1$^\circ$, 2$^\circ$, 3$^\circ$) selon substitution du $C_\alpha$. Polyalcools (glycol, glycérol).
\item \textbf{Propriétés physiques :} Liaisons H $\Rightarrow$ $T_{eb}$ hautes, solubilité eau pour chaînes courtes ($n\le3$).
\item \textbf{Tests carbonyle :} 2,4-DNPH (+) pour les deux. Distinction Aldéhyde/Cétone : \textbf{Fehling, Tollens, Schiff} positifs \textbf{uniquement pour aldéhydes} (réducteurs).
\item \textbf{Équations-bilans :} Fehling ($\ce{Cu^2+ -> Cu2O}$), Tollens ($\ce{Ag+ -> Ag}$), Na/Alcool ($\ce{H2}$).
\item \textbf{Contexte Cameroun :} Éthanol (boissons, carburant, désinfectant), Formol (labo, fraude interdite), Acétone (solvant), Glycérol (savonnerie huile de palme).
\end{itemize}
\end{aretenir}

\coursec{Les alcools : formules, nomenclature et classes}

Dans la section précédente, nous avons découvert que les \cle{composés oxygénés} sont des molécules organiques qui contiennent au moins un atome d'oxygène, et nous avons appris à repérer trois groupes fonctionnels : le groupe \cle{hydroxyle} $-\ce{OH}$ des alcools, le groupe \cle{carbonyle} $\ce{C=O}$ des aldéhydes et des cétones. Nous allons maintenant nous concentrer sur la première grande famille : les \cle{alcools}. C'est une famille que tu connais déjà sans le savoir : l'éthanol est présent dans le vin de palme et le bili-bili, le glycérol entre dans la fabrication des savons artisanaux, et le méthanol est malheureusement célèbre pour les drames des alcools frelatés. Apprenons à les écrire, à les nommer et à les classer.

\courssub{Formule générale d'un alcool}

\begin{definition}[Alcool et formule générale]
Un \cle{alcool} est un composé organique oxygéné dont la molécule possède un \cle{groupe hydroxyle} $-\ce{OH}$ porté par un atome de carbone \emph{tétragonal} (carboné saturé, c'est-à-dire lié à quatre atomes par des liaisons simples).\par
Un alcool \emph{saturé} (dont la chaîne carbonée ne contient que des liaisons simples $\ce{C-C}$) répond à la formule générale :
\[ \text{$C_nH_{2n+1}OH$} \qquad \text{ou, de manière équivalente,} \qquad \text{$C_nH_{2n+2}O$} \]
où $n \geq 1$ est le nombre d'atomes de carbone.
\end{definition}

D'où vient cette formule ? Un alcane saturé a pour formule brute $C_nH_{2n+2}$. Pour obtenir un alcool, on remplace un atome d'hydrogène $-\ce{H}$ par un groupe hydroxyle $-\ce{OH}$ : on enlève un H, on ajoute un O et un H, le bilan est donc $C_nH_{2n+2}O$. Par exemple, pour $n = 2$ : $C_2H_6O$, c'est l'éthanol \ce{CH3-CH2-OH}.

\begin{attention}[Ne pas confondre alcool et phénol]
Le groupe $-\ce{OH}$ doit être porté par un carbone \emph{saturé}. Si le groupe $-\ce{OH}$ est directement fixé sur un cycle benzénique (carbone trigonal, hybridation $sp^2$), le composé est un \cle{phénol}, qui n'est \textbf{pas} un alcool et possède des propriétés très différentes (acidité marquée). De même, si le $-\ce{OH}$ est porté par un carbone d'une double liaison $\ce{C=C}$ (énol), le composé est généralement instable. Au programme de Première : alcool = $-\ce{OH}$ sur carbone tétragonal.
\end{attention}

\courssub{Nomenclature des alcools}

Nommer correctement une molécule, c'est lui donner une identité unique que tous les chimistes du monde comprennent. La nomenclature officielle (IUPAC) suit des règles simples et logiques.

\begin{methode}[Nommer un alcool en 3 étapes]
\begin{enumerate}
\item \textbf{Identifier la chaîne principale} : c'est la chaîne carbonée \emph{la plus longue} qui contient le \cle{carbone fonctionnel} (le carbone lié au $-\ce{OH}$). Elle donne le nom de l'alcane correspondant.
\item \textbf{Numéroter la chaîne} dans le sens qui donne au groupe $-\ce{OH}$ l'\emph{indice de position le plus petit possible}. Les ramifications reçoivent ensuite leurs indices.
\item \textbf{Construire le nom} : on cite d'abord les ramifications par ordre alphabétique avec leurs indices, puis le nom de l'alcane dont on supprime le « e » final, suivi de l'indice du $-\ce{OH}$ et du suffixe \cle{-ol} : \textbf{préfixes + alcane + indice + « ol »}.
\end{enumerate}
\end{methode}

\begin{exemplebox}[Appliquons la méthode]
\begin{itemize}
\item \ce{CH3-CH2-OH} : chaîne de 2 carbones $\to$ éthane ; le $-\ce{OH}$ est forcément en position 1 (inutile de le préciser pour 2 carbones) : \textbf{éthanol}.
\item \ce{CH3-CH2-CH2-OH} : chaîne de 3 carbones, $-\ce{OH}$ en position 1 : \textbf{propan-1-ol}.
\item \chemfig{H_3C-CH(-[2]OH)-CH_3} : chaîne de 3 carbones, $-\ce{OH}$ sur le carbone central : \textbf{propan-2-ol}.
\item \chemfig{CH_3-CH_2-CH(-[2]OH)-CH_3} : chaîne de 4 carbones, $-\ce{OH}$ en position 2 : \textbf{butan-2-ol}.
\item \chemfig{CH_3-CH(-[6]CH_3)-CH_2-OH} : chaîne la plus longue contenant le $-\ce{OH}$ : 3 carbones (propane), un méthyle en position 2, $-\ce{OH}$ en position 1 : \textbf{2-méthylpropan-1-ol}.
\end{itemize}
\end{exemplebox}

\begin{attention}[Erreur fréquente : la numérotation]
Une erreur classique au Probatoire consiste à écrire « propan-3-ol » pour \ce{CH3-CH2-CH2-OH}. C'est \textbf{faux} : on numérote toujours la chaîne de façon à donner au groupe $-\ce{OH}$ l'indice \emph{le plus faible}. Ici, en partant de la droite, le $-\ce{OH}$ est en position 1 : le nom correct est \textbf{propan-1-ol}. Retiens : le groupe fonctionnel est prioritaire, c'est lui qui commande le sens de numérotation.
\end{attention}

\courssub{Les trois classes d'alcools}

Deux alcools peuvent avoir la même formule brute et pourtant des comportements chimiques très différents (nous le verrons avec l'oxydation ménagée). Ce qui les distingue, c'est l'environnement du carbone fonctionnel. D'où la notion de \cle{classe}.

\begin{definition}[Classe d'un alcool]
La \cle{classe} d'un alcool est le \emph{nombre de groupes alkyles} (ou d'atomes de carbone) directement liés au \cle{carbone fonctionnel}, c'est-à-dire le carbone qui porte le groupe $-\ce{OH}$.
\begin{itemize}
\item \cle{Alcool primaire} (I) : le carbone fonctionnel est lié à \textbf{1 seul} atome de carbone (et donc à 2 hydrogènes) : $\ce{R-CH2-OH}$.
\item \cle{Alcool secondaire} (II) : le carbone fonctionnel est lié à \textbf{2} atomes de carbone (et à 1 hydrogène) : $\ce{R1-CH(OH)-R2}$.
\item \cle{Alcool tertiaire} (III) : le carbone fonctionnel est lié à \textbf{3} atomes de carbone (et à aucun hydrogène) : $\ce{R1R2R3C-OH}$.
\end{itemize}
\end{definition}

\begin{popfigure}
\begin{center}
\begin{tikzpicture}[node distance=1.5cm]
% Butan-1-ol (Primaire)
\node (p) {\chemfig{CH_3-CH_2-CH_2-\textcolor{popink}{C}(-[2]H)(-[6]H)-OH}};
\node[below=0.5cm of p] {\small \textbf{Butan-1-ol (primaire)}};

% Butan-2-ol (Secondaire)
\node[right=4cm of p] (s) {\chemfig{CH_3-CH_2-\textcolor{popink}{C}(-[2]OH)(-[6]H)-CH_3}};
\node[below=0.5cm of s] {\small \textbf{Butan-2-ol (secondaire)}};

% 2-Méthylpropan-2-ol (Tertiaire)
\node[right=4cm of s] (t) {\chemfig{CH_3-\textcolor{popink}{C}(-[2]OH)(-[6]CH_3)-CH_3}};
\node[below=0.5cm of t] {\small \textbf{2-Méthylpropan-2-ol (tertiaire)}};
\end{tikzpicture}
\end{center}
\legende{Les trois classes d'alcools de formule brute \ce{C4H10O} : le carbone fonctionnel (en rouge) est lié à 1, 2 ou 3 atomes de carbone.}
\end{popfigure}

\begin{methode}[Déterminer la classe d'un alcool à partir de sa formule développée]
\begin{enumerate}
\item Repérer le groupe $-\ce{OH}$ et donc le \textbf{carbone fonctionnel}.
\item Compter le nombre d'atomes de \textbf{carbone} directement liés à ce carbone fonctionnel (les hydrogènes ne comptent pas).
\item Conclure : 1 carbone $\to$ primaire ; 2 carbones $\to$ secondaire ; 3 carbones $\to$ tertiaire.
\end{enumerate}
\end{methode}

\begin{popfigure}
\begin{center}
\begin{tikzpicture}[
node distance=0.55cm and 1.1cm,
qbox/.style={draw=popblue, fill=popblueL, rounded corners=3pt, text width=4.6cm, align=center, font=\small, thick},
rbox/.style={draw=popgreen, fill=popgreenL, rounded corners=3pt, text width=3.2cm, align=center, font=\small\bfseries, thick},
lbl/.style={font=\small\itshape, text=popdark}]
\node[qbox] (q1) {Combien d'atomes de carbone sont liés au carbone fonctionnel C--OH ?};
\node[rbox, below left=1.1cm and -1.2cm of q1] (p) {1 carbone\\ Alcool PRIMAIRE};
\node[rbox, below=1.1cm of q1] (s) {2 carbones\\ Alcool SECONDAIRE};
\node[rbox, below right=1.1cm and -1.2cm of q1] (t) {3 carbones\\ Alcool TERTIAIRE};
\draw[-{Stealth}, thick, popblue] (q1.south) -- ++(0,-0.35) -| (p.north) node[lbl, pos=0.25, left]{1};
\draw[-{Stealth}, thick, popblue] (q1.south) -- (s.north) node[lbl, midway, right]{2};
\draw[-{Stealth}, thick, popblue] (q1.south) -- ++(0,-0.35) -| (t.north) node[lbl, pos=0.25, right]{3};
\end{tikzpicture}
\end{center}
\legende{Arbre de décision : déterminer la classe d'un alcool en une seule question.}
\end{popfigure}

\begin{exoresolu}[Classe du 2-méthylpropan-1-ol]
On donne l'alcool de formule semi-développée \ce{(CH3)2CH-CH2-OH}. Nommer cet alcool et préciser sa classe.
\tcblower
\textbf{Solution détaillée.}\par
\textbf{Étape 1 — Chaîne principale.} La chaîne la plus longue contenant le carbone lié au $-\ce{OH}$ comporte 3 atomes de carbone : c'est un dérivé du \textbf{propane}. Il y a une ramification méthyle $-\ce{CH3}$ sur le carbone n°2.\par
\textbf{Étape 2 — Numérotation.} On numérote en partant de l'extrémité portant le $-\ce{OH}$ : le groupe hydroxyle est en position 1, le méthyle en position 2.\par
\textbf{Étape 3 — Nom.} \textbf{2-méthylpropan-1-ol}.\par
\textbf{Classe.} Le carbone fonctionnel est le carbone du groupe $-\ce{CH2-OH}$. Il est lié à \emph{un seul} atome de carbone (le carbone n°2) et à deux hydrogènes. C'est donc un alcool \textbf{primaire}.\par
\emph{Vérification :} la formule brute est \ce{C4H10O}, conforme à $C_nH_{2n+2}O$ avec $n = 4$.
\end{exoresolu}

\begin{attention}[Piège : ne pas compter les hydrogènes]
La classe se détermine en comptant les \textbf{carbones} voisins du carbone fonctionnel, jamais les hydrogènes. Ainsi \ce{CH3-OH} (méthanol) : son carbone fonctionnel n'est lié à \emph{aucun} carbone ; par convention, on le range parmi les alcools \emph{primaires} (il se comporte comme eux). Autre piège : dans \ce{(CH3)3C-OH}, le carbone central n'a \emph{aucun} hydrogène : c'est un alcool tertiaire, et cette absence d'hydrogène sur le carbone fonctionnel aura des conséquences majeures en oxydation (section suivante).
\end{attention}

\courssub{Isomérie chez les alcools ; les éthers-oxydes}

Des \cle{isomères} sont des composés qui ont la \emph{même formule brute} mais des formules développées différentes, donc des propriétés différentes. Chez les alcools, on distingue trois types d'isomérie.

\begin{definition}[Les trois types d'isomérie]
\begin{itemize}
\item \cle{Isomérie de chaîne} : le squelette carboné diffère (chaîne linéaire ou ramifiée). Exemple : butan-1-ol \ce{CH3-CH2-CH2-CH2-OH} et 2-méthylpropan-1-ol \ce{(CH3)2CH-CH2-OH}.
\item \cle{Isomérie de position} : la chaîne est identique mais la position du groupe $-\ce{OH}$ diffère. Exemple : butan-1-ol et butan-2-ol \ce{CH3-CH2-CH(OH)-CH3}.
\item \cle{Isomérie de fonction} : le groupe fonctionnel diffère. Un alcool $C_nH_{2n+2}O$ est isomère de fonction des \cle{éthers-oxydes} $R-O-R'$. Exemple : l'éthanol \ce{CH3-CH2-OH} et le méthoxyméthane \ce{CH3-O-CH3} ont tous deux la formule brute \ce{C2H6O}.
\end{itemize}
\end{definition}

Un \cle{éther-oxyde} est un composé oxygéné dans lequel l'atome d'oxygène est inséré \emph{entre deux groupes alkyles} : $R-O-R'$. Pour le nommer, on considère le plus petit groupe alkyle associé à l'oxygène comme un préfixe \cle{alcoxy-} (méthoxy, éthoxy\ldots) accolé au nom de l'alcane correspondant au plus grand groupe : ainsi \ce{CH3-O-CH2-CH3} se nomme \textbf{méthoxyéthane}.

\begin{exemplebox}[Exemple chiffré : les 7 isomères de \ce{C4H10O}]
La formule brute \ce{C4H10O} ($n = 4$) correspond à \textbf{4 alcools} et \textbf{3 éthers-oxydes}, soit 7 isomères au total :
\begin{center}
\begin{tabularx}{\linewidth}{|l|X|l|}
\hline
\rowcolor{popblueL} \textbf{Formule semi-développée} & \textbf{Nom} & \textbf{Classe / fonction} \\
\hline
\ce{CH3-CH2-CH2-CH2-OH} & Butan-1-ol & Alcool primaire \\
\hline
\ce{(CH3)2CH-CH2-OH} & 2-Méthylpropan-1-ol & Alcool primaire \\
\hline
\ce{CH3-CH2-CH(OH)-CH3} & Butan-2-ol & Alcool secondaire \\
\hline
\ce{(CH3)3C-OH} & 2-Méthylpropan-2-ol & Alcool tertiaire \\
\hline
\ce{CH3-O-CH2-CH2-CH3} & Méthoxypropane & Éther-oxyde \\
\hline
\ce{CH3-O-CH(CH3)2} & 2-Méthoxypropane & Éther-oxyde \\
\hline
\ce{CH3-CH2-O-CH2-CH3} & Éthoxyéthane & Éther-oxyde \\
\hline
\end{tabularx}
\end{center}
On compte donc \textbf{2 alcools primaires, 1 alcool secondaire, 1 alcool tertiaire} et \textbf{3 éthers-oxydes}.
\end{exemplebox}

\begin{methode}[Dénombrer les isomères alcools d'une formule brute]
\begin{enumerate}
\item Écrire tous les \textbf{squelettes carbonés} possibles (isomérie de chaîne) : pour 4 C, la chaîne linéaire butane et la chaîne ramifiée 2-méthylpropane.
\item Sur chaque squelette, placer le $-\ce{OH}$ sur chaque carbone \emph{non équivalent} (isomérie de position).
\item Nommer chaque alcool obtenu et vérifier qu'aucun n'est compté deux fois.
\item Compléter éventuellement avec les éthers-oxydes $R-O-R'$ (isomérie de fonction).
\end{enumerate}
\end{methode}

\begin{savaistu}[Le méthanol, un alcool mortel]
Le \cle{méthanol} \ce{CH3-OH}, alcool primaire à un seul carbone, est extrêmement toxique : l'organisme le transforme en méthanal puis en acide méthanoïque, qui détruisent le nerf optique. Quelques millilitres suffisent à provoquer la \emph{cécité}, et une trentaine de millilitres entraînent la mort. C'est pourquoi les alcools frelatés — boissons artisanales adultérées ou mal distillées — provoquent régulièrement des drames, y compris au Cameroun. L'éthanol des boissons et le méthanol ne diffèrent que par un seul carbone, mais leurs effets sont sans commune mesure : en chimie, la structure fait toute la différence !
\end{savaistu}

\begin{exercice}[À toi de jouer]
On considère les composés suivants : (a) \ce{CH3-CH2-CH2-OH} ; (b) \ce{CH3-CH(OH)-CH3} ; (c) \ce{CH3-O-CH2-CH3} ; (d) \ce{(CH3)2CH-OH}.
\begin{enumerate}
\item Nommer chaque composé.
\item Préciser la classe de chaque alcool.
\item Indiquer deux couples d'isomères et préciser le type d'isomérie.
\end{enumerate}
\end{exercice}

\begin{corrige}[Exercice 1]
\textbf{1. Noms :} (a) propan-1-ol ; (b) propan-2-ol ; (c) méthoxyéthane ; (d) propan-2-ol (identique à (b)).\par
\textbf{2. Classes :} (a) le carbone fonctionnel $-\ce{CH2-OH}$ est lié à 1 carbone : alcool \textbf{primaire}. (b) et (d) : le carbone fonctionnel $-\ce{CH(OH)}-$ est lié à 2 carbones : alcool \textbf{secondaire}. (c) n'est pas un alcool mais un éther-oxyde.\par
\textbf{3. Isomères :} (a) et (b) ont tous deux la formule brute \ce{C3H8O} avec le $-\ce{OH}$ en positions différentes : \textbf{isomérie de position}. (a) et (c) ont la même formule brute \ce{C3H8O} mais des fonctions différentes (alcool / éther-oxyde) : \textbf{isomérie de fonction}.
\end{corrige}

\begin{aretenir}[L'essentiel de la section]
\begin{itemize}
\item Un alcool saturé répond à la formule générale $C_nH_{2n+1}OH$, soit $C_nH_{2n+2}O$.
\item Nomenclature : chaîne la plus longue contenant le C--OH, numérotation donnant au $-\ce{OH}$ l'indice le plus petit, nom = préfixes + alcane + indice + « ol ».
\item La \cle{classe} d'un alcool est le nombre de carbones liés au carbone fonctionnel : 1 $\to$ primaire, 2 $\to$ secondaire, 3 $\to$ tertiaire.
\item Les alcools présentent des isoméries de chaîne, de position et de fonction (avec les éthers-oxydes $R-O-R'$, nommés alcoxy + alcane).
\item \ce{C4H10O} : 4 alcools (2 primaires, 1 secondaire, 1 tertiaire) et 3 éthers-oxydes.
\end{itemize}
\end{aretenir}

Maintenant que tu sais écrire, nommer et classer les alcools, une question naturelle se pose : pourquoi l'éthanol se mélange-t-il à l'eau alors que l'huile de palme ne s'y mélange pas ? Pourquoi les points d'ébullition des alcools sont-ils si élevés comparés à ceux des alcanes ? C'est ce que nous explorerons dans la section suivante, consacrée aux propriétés physiques des alcools et aux polyalcools comme le glycol et le glycérol.

\coursec{Propriétés physiques des alcools et polyalcools}

Dans la section précédente, nous avons appris à reconnaître, nommer et classer les alcools grâce à leur formule développée. Mais une question naturelle se pose : pourquoi l'éthanol est-il liquide à température ambiante alors que le propane, de masse molaire presque identique, est un gaz ? Pourquoi le vin de palme se mélange-t-il si bien à l'eau, alors que l'huile de palme, elle, flotte ? Toutes ces observations s'expliquent par la \cle{structure moléculaire} des alcools, et plus précisément par une petite liaison aux grands effets : la \cle{liaison hydrogène}. C'est l'objet de cette section.

\courssub{Polarité de la liaison O--H et liaison hydrogène}

\courssub{Une liaison O--H fortement polarisée}

Dans une molécule d'alcool, le groupe caractéristique est le groupe \cle{hydroxyle} $-\ce{OH}$. Or l'atome d'oxygène est beaucoup plus \cle{électronégatif} que l'atome d'hydrogène : il attire vers lui les électrons de la liaison covalente $\ce{O}-\ce{H}$. Il en résulte une \cle{liaison polarisée} : l'oxygène porte une charge partielle négative $\delta^-$ et l'hydrogène une charge partielle positive $\delta^+$.

\[
\chemfig{C-[2]O^{\delta^-}-[0]H^{\delta^+}}
\]

La molécule d'alcool est donc une molécule \cle{polaire} : elle possède un pôle négatif (l'oxygène) et un pôle positif (l'hydrogène du groupe hydroxyle).

\courssub{La liaison hydrogène : un pont entre molécules}

Intuitivement, imagine deux molécules d'éthanol voisines : le pôle positif de l'une (le H du $-\ce{OH}$) est attiré par le pôle négatif de l'autre (l'oxygène, qui possède des doublets non liants). Cette attraction électrostatique crée un véritable \emph{pont} entre les molécules.

\begin{definition}[Liaison hydrogène]
Une \cle{liaison hydrogène} est une attraction électrostatique qui s'établit entre l'atome d'hydrogène d'un groupe $-\ce{OH}$ (porteur d'une charge partielle $\delta^+$) et l'atome d'oxygène, riche en électrons, d'une molécule voisine. On la représente par des pointillés : $\ce{O-H}\cdots\ce{O}$. Elle est environ dix fois plus faible qu'une liaison covalente, mais bien plus forte que les simples attractions entre molécules non polaires (forces de Van der Waals).
\end{definition}

\begin{popfigure}
\begin{center}
\begin{tikzpicture}[scale=1.05, every node/.style={font=\small}]
% Molécule 1
\node (C1) at (0,0) {$\mathrm{CH_3}$};
\node (C2) at (1.5,0) {$\mathrm{CH_2}$};
\node (O1) at (3.1,0) {$\mathrm{O}^{\delta^-}$};
\node (H1) at (4.6,0) {$\mathrm{H}^{\delta^+}$};
\draw[thick,popdark] (C1)--(C2)--(O1)--(H1);
% Molécule 2
\node (O2) at (6.6,0) {$\mathrm{O}^{\delta^-}$};
\node (H2) at (8.1,0) {$\mathrm{H}^{\delta^+}$};
\node (C3) at (8.1,-1.4) {$\mathrm{CH_2}$};
\node (C4) at (8.1,-2.7) {$\mathrm{CH_3}$};
\draw[thick,popdark] (O2)--(H2)--(C3)--(C4);
% Liaison hydrogène
\draw[dashed,very thick,popink] (H1)--(O2) node[midway,above,popink]{liaison hydrogène};
% Doublets
\node[popblue,above] at (O1.north) {$\bullet\bullet$};
\node[popblue,above] at (O2.north) {$\bullet\bullet$};
% Légendes
\node[popmuted,below=0.1cm] at (0.75,-0.35) {molécule 1};
\node[popmuted,right=0.1cm] at (8.3,-1.4) {molécule 2};
\end{tikzpicture}
\end{center}
\legende{Liaison hydrogène (pointillés roses) entre deux molécules d'éthanol : le H$^{\delta^+}$ de l'une est attiré par l'O$^{\delta^-}$ de l'autre.}
\end{popfigure}

\begin{attention}[Ne pas confondre]
La liaison hydrogène n'est \textbf{pas} une liaison covalente : aucun électron n'est mis en commun. C'est une attraction \emph{entre molécules} (force intermoléculaire). Autre piège : elle n'existe de façon notable que lorsque l'hydrogène est lié à un atome très électronégatif ($\ce{O}$, $\ce{N}$ ou $\ce{F}$). Les alcanes, qui ne possèdent que des liaisons $\ce{C-H}$ faiblement polarisées, ne forment \textbf{pas} de liaisons hydrogène.
\end{attention}

\courssub{Conséquences sur les températures de fusion et d'ébullition}

\courssub{Des alcools beaucoup moins volatils que les alcanes}

Pour faire bouillir un liquide, il faut fournir assez d'énergie pour \emph{séparer} les molécules les unes des autres. Or, dans un alcool, les molécules sont retenues par les liaisons hydrogène : il faut donc davantage d'énergie, donc une température plus élevée.

\begin{propriete}[Températures d'ébullition élevées]
À masse molaire comparable, un alcool bout à une température \textbf{nettement plus élevée} qu'un alcane, à cause des liaisons hydrogène qui associent ses molécules. Les alcools courants (méthanol, éthanol, propanols, butanols) sont \textbf{liquides à température ambiante} (leurs températures de fusion sont inférieures à \SI{0}{\celsius} pour la plupart : celle de l'éthanol vaut \SI{-114}{\celsius}).
\end{propriete}

\begin{exemplebox}[Éthanol contre propane : le match des masses voisines]
Comparons deux molécules de masses molaires très proches :
\begin{itemize}
\item éthanol \ce{C2H5OH} : $M = \SI{46}{g.mol^{-1}}$, $\theta_{eb} = \SI{78}{\celsius}$ ;
\item propane \ce{C3H8} : $M = \SI{44}{g.mol^{-1}}$, $\theta_{eb} = \SI{-42}{\celsius}$.
\end{itemize}
Écart : \SI{120}{\celsius} ! Le propane est un gaz (celui des bouteilles de gaz domestiques), l'éthanol est un liquide. Toute la différence vient des liaisons hydrogène, présentes dans l'éthanol, absentes dans le propane.
\end{exemplebox}

\courssub{Influence de la longueur de la chaîne carbonée}

Lorsque le nombre d'atomes de carbone augmente dans la série des alcan-1-ols, la masse molaire augmente et les attractions entre les chaînes carbonées (forces de Van der Waals) s'ajoutent aux liaisons hydrogène : la température d'ébullition \cle{augmente régulièrement}.

\begin{popfigure}
\begin{center}
\begin{tikzpicture}
\begin{axis}[
width=0.85\linewidth, height=6.5cm,
xlabel={Nombre d'atomes de carbone $n$},
ylabel={$\theta_{eb}$ (\si{\celsius})},
xmin=0.5, xmax=8.5, ymin=-60, ymax=200,
xtick={1,2,3,4,5,6,7,8},
grid=both, grid style={popline!40},
legend style={at={(0.03,0.97)}, anchor=north west, font=\small}
]
\addplot[popcurve, mark=*, mark options={popblue}] coordinates {(1,65) (2,78) (3,97) (4,118) (5,138) (6,157) (7,176) (8,195)};
\addlegendentry{Alcan-1-ols $C_nH_{2n+1}OH$}
\addplot[popcurve, poporange, mark=square*, mark options={poporange}] coordinates {(1,-162) (2,-89) (3,-42) (4,-0.5) (5,36) (6,69) (7,98) (8,126)};
\addlegendentry{Alcanes $C_nH_{2n+2}$}
\end{axis}
\end{tikzpicture}
\end{center}
\legende{Températures d'ébullition des alcan-1-ols et des alcanes en fonction du nombre de carbones. La courbe des alcools est toujours bien au-dessus : c'est la signature des liaisons hydrogène.}
\end{popfigure}

\courssub{Influence de la classe et de la ramification}

À formule brute égale (isomères), tous les alcools ne bouillent pas à la même température. Plus la molécule est \emph{ramifiée}, plus elle est « compacte » : les contacts entre molécules sont moins étendus, les attractions de Van der Waals sont plus faibles, et l'ébullition est abaissée.

\begin{exemplebox}[Trois isomères en \ce{C4H10O}]
\begin{center}
\begin{tabularx}{\linewidth}{|l|X|c|}
\hline
\rowcolor{popblueL} \textbf{Alcool} & \textbf{Classe} & $\boldsymbol{\theta_{eb}}$ \\
\hline
Butan-1-ol & primaire & \SI{118}{\celsius} \\
\hline
Butan-2-ol & secondaire & \SI{100}{\celsius} \\
\hline
2-méthylpropan-2-ol & tertiaire & \SI{82}{\celsius} \\
\hline
\end{tabularx}
\end{center}
Conclusion générale : à formule brute égale, $\theta_{eb}(\text{primaire}) > \theta_{eb}(\text{secondaire}) > \theta_{eb}(\text{tertiaire})$. La ramification \cle{abaisse} la température d'ébullition.
\end{exemplebox}

\courssub{Solubilité des alcools dans l'eau}

\courssub{Le duel hydrophile / hydrophobe}

L'eau est le champion des liaisons hydrogène : c'est un solvant \cle{polaire}. Une molécule d'alcool possède deux parties de caractère opposé :
\begin{itemize}
\item le groupe $-\ce{OH}$, \cle{hydrophile} (qui « aime l'eau ») : il forme des liaisons hydrogène avec les molécules d'eau et favorise la dissolution ;
\item la chaîne carbonée, \cle{hydrophobe} (qui « fuit l'eau »), comme l'huile : elle s'oppose à la dissolution.
\end{itemize}
La solubilité résulte de ce duel : tout dépend de la taille relative des deux parties.

\begin{methode}[Prévoir la solubilité d'un alcool dans l'eau]
\begin{enumerate}
\item Repérer la partie hydrophile (le ou les groupes $-\ce{OH}$) et la partie hydrophobe (la chaîne carbonée).
\item Comparer leurs « poids » respectifs : compter le nombre d'atomes de carbone par groupe $-\ce{OH}$.
\item Conclure : 
\begin{itemize}
\item si la chaîne compte \textbf{au plus 3 carbones} par groupe $-\ce{OH}$, l'alcool est \textbf{miscible à l'eau en toutes proportions} ;
\item pour \textbf{4 carbones} (butanols), la solubilité est \textbf{limitée} (partiellement miscible, environ \SI{73}{g.L^{-1}} pour le butan-1-ol) ;
\item \textbf{au-delà de 4 carbones}, la partie hydrophobe domine et la solubilité devient faible (inférieure à \SI{30}{g.L^{-1}} pour le pentan-1-ol).
\end{itemize}
\end{enumerate}
\end{methode}

\begin{propriete}[Évolution de la solubilité]
Les alcools à courte chaîne (\cle{méthanol}, \cle{éthanol}, \cle{propan-1-ol}, \cle{propan-2-ol}) sont \textbf{miscibles à l'eau en toutes proportions}. La solubilité \textbf{diminue} quand la chaîne carbonée s'allonge : le butan-1-ol est encore partiellement soluble (environ \SI{73}{g.L^{-1}} à \SI{20}{\celsius}), mais à partir du pentan-1-ol (\SI{22}{g.L^{-1}}) la solubilité devient faible, et les alcools supérieurs (hexanol, heptanol...) sont pratiquement insolubles.
\end{propriete}

\begin{attention}[Erreur fréquente]
Ne dis jamais « tous les alcools sont solubles dans l'eau ». C'est faux ! Seuls les alcools à \textbf{courte chaîne} (1 à 3 carbones) sont miscibles en toutes proportions. Le butan-1-ol (4 carbones) a une solubilité limitée. À partir du \cle{pentanol} (5 carbones), la partie hydrophobe l'emporte : l'alcool forme une phase séparée, comme l'huile de palme sur l'eau. Au Probatoire, précise toujours : « miscible \emph{car} la chaîne carbonée est courte (≤ 3 C) ».
\end{attention}

\begin{exoresolu}[Solubilité et ébullition : deux prévisions]
\textbf{Énoncé.} On considère le méthanol \ce{CH3OH} et l'octan-1-ol \ce{CH3-(CH2)6-CH2OH}. (1) Lequel est miscible à l'eau ? Justifier. (2) Lequel a la température d'ébullition la plus élevée ? Justifier.
\tcblower
\textbf{Solution.} (1) Le méthanol possède un groupe $-\ce{OH}$ hydrophile et une chaîne d'un seul carbone : la partie hydrophile domine, il forme des liaisons hydrogène avec l'eau et y est \textbf{miscible en toutes proportions}. L'octan-1-ol possède une longue chaîne hydrophobe de 8 carbones qui masque l'effet du $-\ce{OH}$ : il est \textbf{pratiquement insoluble} dans l'eau (solubilité $\approx \SI{0,3}{g.L^{-1}}$). (2) L'octan-1-ol a la masse molaire la plus grande ($M = \SI{130}{g.mol^{-1}}$ contre \SI{32}{g.mol^{-1}}) : les forces de Van der Waals entre ses longues chaînes s'ajoutent aux liaisons hydrogène. Sa température d'ébullition (\SI{195}{\celsius}) est donc bien plus élevée que celle du méthanol (\SI{65}{\celsius}).
\end{exoresolu}

\courssub{Les polyalcools : glycol et glycérol}

Certains composés possèdent \emph{plusieurs} groupes hydroxyle : ce sont les \cle{polyalcools} (ou polyols). Chaque groupe $-\ce{OH}$ supplémentaire multiplie les liaisons hydrogène : ces composés sont très solubles dans l'eau, visqueux, et bouillent à des températures élevées.

\begin{definition}[Deux polyalcools de référence]
\begin{itemize}
\item L'\cle{éthane-1,2-diol}, appelé \emph{glycol} : \chemfig{HO-CH_2-CH_2-OH}, de formule brute \ce{C2H6O2}. C'est un diol (deux groupes $-\ce{OH}$), liquide, miscible à l'eau, $\theta_{eb} = \SI{197}{\celsius}$.
\item Le \cle{propane-1,2,3-triol}, appelé \emph{glycérol} (ou glycérine) : \chemfig{HO-CH_2-CH(-[2]OH)-CH_2-OH}, de formule brute \ce{C3H8O3}. C'est un triol, liquide visqueux et sirupeux, miscible à l'eau, $\theta_{eb} = \SI{290}{\celsius}$ (avec décomposition).
\end{itemize}
\end{definition}

\begin{exemplebox}[Des propriétés remarquables]
Grâce à leurs nombreux groupes $-\ce{OH}$, le glycol et le glycérol sont \textbf{hygroscopiques} (ils retiennent l'eau de l'air) et \textbf{très solubles} dans l'eau. Le glycol sert d'\cle{antigel} dans les radiateurs automobiles (il abaisse le point de congélation et relève le point d'ébullition de l'eau) ; le glycérol, non toxique et adoucissant, entre dans la composition des savons, des crèmes cosmétiques et de nombreux médicaments (sirops, suppositoires).
\end{exemplebox}

\begin{savaistu}[Glycol, glycérol et vie quotidienne au Cameroun]
Le glycol versé dans le radiateur d'une voiture abaisse la température de congélation de l'eau (et élève sa température d'ébullition) : le moteur est protégé du gel comme de la surchauffe, utile même sous nos latitudes pour éviter la surchauffe dans les embouteillages de Douala ou Yaoundé. Le glycérol, lui, est un sous-produit de la \emph{saponification} : lorsque les savonneries artisanales de Douala font réagir l'huile de palme avec la soude, elles obtiennent du savon \emph{et} du glycérol, que l'on retrouve dans les savons « glycérinés » doux pour la peau. La chimie des polyalcools est ainsi présente jusque dans les marchés camerounais !
\end{savaistu}

\begin{aretenir}[L'essentiel de la section]
\begin{itemize}
\item La liaison $\ce{O-H}$ est polarisée ($\ce{O}^{\delta^-}$, $\ce{H}^{\delta^+}$) : les molécules d'alcool s'associent par des \cle{liaisons hydrogène} (forces intermoléculaires fortes).
\item Les alcools bouillent à des températures \textbf{bien plus élevées} que les alcanes de masse voisine (éthanol \SI{78}{\celsius} contre propane \SI{-42}{\celsius}) et sont liquides à température ambiante.
\item L'ébullition \textbf{augmente} avec la longueur de la chaîne ; à formule brute égale : primaire $>$ secondaire $>$ tertiaire (la ramification abaisse l'ébullition).
\item La solubilité dans l'eau \textbf{diminue} quand la chaîne s'allonge : miscibles en toutes proportions jusqu'à 3 carbones (propanols), solubilité limitée pour 4 carbones (butanols), très faible au-delà (pentanols et supérieurs).
\item Les polyalcools (glycol \ce{C2H6O2}, glycérol \ce{C3H8O3}) cumulent les groupes $-\ce{OH}$ : très solubles, visqueux, aux applications nombreuses (antigel, savons, cosmétiques, pharmacie).
\end{itemize}
\end{aretenir}

\begin{exercice}[Pour t'entraîner]
On donne trois alcools : A = propan-1-ol, B = 2-méthylpropan-2-ol, C = hexan-1-ol. (1) Classer A, B, C par température d'ébullition croissante en justifiant. (2) Lequel est miscible à l'eau en toutes proportions ? (3) Expliquer pourquoi le glycérol est beaucoup plus soluble dans l'eau que l'hexan-1-ol, alors que sa chaîne carbonée compte 3 atomes de carbone.
\end{exercice}

\begin{corrige}[Exercice]
(1) B ($\SI{82}{\celsius}$, tertiaire ramifié, 4 C) $<$ A ($\SI{97}{\celsius}$, primaire, 3 C, chaîne linéaire mais plus courte que C) $<$ C ($\SI{157}{\celsius}$, primaire, 6 C, chaîne la plus longue). La ramification de B abaisse fortement son ébullition ; la longue chaîne de C élève la sienne par forces de Van der Waals. (2) Le propan-1-ol A : chaîne courte (3 C), la partie hydrophile domine, il est miscible en toutes proportions. (3) Le glycérol possède \emph{trois} groupes $-\ce{OH}$ pour seulement 3 carbones (ratio 1 OH / 1 C) : il forme de très nombreuses liaisons hydrogène avec l'eau, ce qui le rend miscible, tandis que l'hexan-1-ol n'a qu'un seul $-\ce{OH}$ pour 6 carbones hydrophobes (ratio 1 OH / 6 C), la partie hydrophobe l'emporte.
\end{corrige}

\coursec{Aldéhydes et cétones : formules et nomenclature}

Dans la section précédente, nous avons vu que les alcools possèdent le groupe hydroxyle \cle{$-$OH} et que ce groupe explique leurs propriétés physiques remarquables (liaisons hydrogène, solubilité dans l'eau, températures d'ébullition élevées). Mais le monde des composés oxygénés ne s'arrête pas là. Lorsque tu laisses une bouteille de vin de palme ouverte quelques jours, elle « tourne » : l'odeur devient piquante, le goût devient âcre. Ce qui s'est passé, c'est que l'éthanol a été transformé en un nouveau composé oxygéné : l'\cle{éthanal}, puis en acide éthanoïque (le vinaigre). Ce nouveau composé appartient à une famille que nous allons maintenant étudier : les \cle{aldéhydes}, et leurs cousines, les \cle{cétones}.

\courssub{Le groupe carbonyle : le cœur des aldéhydes et des cétones}

\textbf{Une intuition d'abord.} Dans un alcool, l'atome d'oxygène est relié au carbone par une \emph{liaison simple} ($\ce{C-O-H}$). Dans les aldéhydes et les cétones, l'oxygène est relié au carbone par une \emph{liaison double} $\ce{C=O}$. Ce groupe $\ce{C=O}$ s'appelle le \cle{groupe carbonyle}. C'est lui qui commande toute la chimie de ces deux familles.

\begin{definition}[Aldéhydes et cétones]
Un \cle{aldéhyde} est un composé organique qui possède le groupe carbonyle \cle{$\ce{C=O}$} en \textbf{extrémité} de chaîne carbonée. Son groupe caractéristique s'écrit \cle{$-$CHO} : le carbone du carbonyle est lié à \textbf{un atome d'hydrogène} et à \textbf{un seul groupe alkyle} R (ou à un hydrogène pour le méthanal). Formule générale : R$-$CHO.

Une \cle{cétone} est un composé organique qui possède le groupe carbonyle \cle{$\ce{C=O}$} à l'\textbf{intérieur} de la chaîne carbonée. Le carbone du carbonyle est lié à \textbf{deux groupes alkyles} R et R'. Formule générale : R$-$CO$-$R'.
\end{definition}

\begin{aretenir}[Formule brute générale]
Les aldéhydes et les cétones \textbf{saturés non cycliques} (une seule liaison double $\ce{C=O}$, toutes les liaisons C$-$C simples, chaîne ouverte) ont pour formule brute générale :
\[ C_nH_{2n}O \qquad (n \geq 1 \text{ pour un aldéhyde, } n \geq 3 \text{ pour une cétone}) \]
Remarque bien que c'est la \textbf{même} formule générale pour les deux familles : c'est l'origine de l'isomérie de fonction que nous verrons plus loin. Une cétone possède au minimum 3 atomes de carbone, car le carbonyle doit être entouré de deux groupes alkyles.
\end{aretenir}

\textbf{Comment distinguer visuellement les deux familles ?} Pose-toi une seule question : le carbone du $\ce{C=O}$ porte-t-il un hydrogène ?
\begin{itemize}
\item S'il porte un \textbf{H} (groupe $-$CHO en bout de chaîne) $\Rightarrow$ c'est un \textbf{aldéhyde} ;
\item s'il porte \textbf{deux carbones} (groupe $-$CO$-$ au milieu de la chaîne) $\Rightarrow$ c'est une \textbf{cétone}.
\end{itemize}

\begin{attention}[Le groupe $-$CHO ne s'écrit jamais $-$COH]
Écrire R$-$COH est une erreur classique : cela laisserait croire que l'oxygène est lié au carbone par une liaison simple et porte l'hydrogène, comme dans un alcool ! Dans un aldéhyde, l'hydrogène est lié \textbf{au carbone}, et l'oxygène est lié au carbone par une \textbf{double liaison} : R$-$CHO signifie R$-$CH$=$O. Retiens : \textbf{CHO = C doublement lié à O, H sur le C}.
\end{attention}

\courssub{Nomenclature des aldéhydes : le suffixe « -al »}

\textbf{L'idée directrice.} La nomenclature des aldéhydes suit celle des alcools, avec deux simplifications heureuses : le groupe $-$CHO est \emph{toujours} en bout de chaîne, donc le carbone du carbonyle est \emph{toujours} le carbone n°1. Il n'y a donc \textbf{jamais d'indice de position} à écrire.

\begin{methode}[Nommer un aldéhyde]
\begin{enumerate}
\item Repérer le groupe $-$CHO : son carbone est obligatoirement le carbone \textbf{n°1}.
\item Identifier la chaîne la plus longue \textbf{contenant} ce carbone ; elle donne la racine (méthan-, éthan-, propan-, butan-\ldots).
\item Remplacer le « e » final de l'alcane par le suffixe \cle{-al} : \textbf{éthane $\rightarrow$ éthanal}.
\item Indiquer les ramifications éventuelles avec leur indice de position (le CHO étant en position 1).
\end{enumerate}
\end{methode}

\begin{exemplebox}[Les quatre premiers aldéhydes]
\begin{center}
\begin{tabularx}{\linewidth}{|l|X|X|}
\hline
\rowcolor{popblueL} \textbf{Formule brute} & \textbf{Formule semi-développée} & \textbf{Nom} \\
\hline
\ce{CH2O} & \ce{HCHO} (ou H$-$CHO) & \textbf{méthanal} (formaldéhyde) \\
\hline
\ce{C2H4O} & \ce{CH3-CHO} & \textbf{éthanal} (acétaldéhyde) \\
\hline
\ce{C3H6O} & \ce{CH3-CH2-CHO} & \textbf{propanal} \\
\hline
\ce{C4H8O} & \ce{CH3-CH2-CH2-CHO} & \textbf{butanal} \\
\hline
\end{tabularx}
\end{center}
Vérifie la formule générale : pour l'éthanal, $n = 2$, donc $C_2H_{2\times 2}O = \ce{C2H4O}$. C'est bien cela.
\end{exemplebox}

\begin{attention}[Ne jamais écrire « éthan-1-al »]
Erreur extrêmement fréquente au Probatoire ! Certains élèves, par habitude des alcools (éthan-1-ol, propan-2-ol), écrivent « éthan-1-al » ou « butan-1-al ». C'est \textbf{faux} : comme le groupe $-$CHO est forcément en extrémité de chaîne, le carbone du carbonyle est \textbf{toujours} le carbone n°1, sans ambiguïté possible. L'indice de position est donc \textbf{inutile et interdit}. On écrit \textbf{éthanal}, \textbf{butanal}, tout simplement. En revanche, pour un aldéhyde ramifié comme \ce{(CH3)2CH-CHO}, on écrit bien \textbf{2-méthylpropanal} : les indices servent alors à localiser les \emph{ramifications}, pas le CHO.
\end{attention}

\courssub{Nomenclature des cétones : le suffixe « -one »}

\textbf{L'idée directrice.} Pour une cétone, le carbonyle peut se trouver à différentes positions \emph{à l'intérieur} de la chaîne. Il faut donc préciser sa position par un indice, en choisissant le sens de numérotation qui donne au carbonyle \textbf{l'indice le plus petit possible}.

\begin{methode}[Nommer une cétone]
\begin{enumerate}
\item Identifier la chaîne carbonée la plus longue \textbf{contenant} le groupe $\ce{C=O}$ ; elle donne la racine.
\item \textbf{Numéroter} la chaîne de façon à donner au carbone du carbonyle l'indice \textbf{le plus petit possible}.
\item Remplacer le « e » final de l'alcane par le suffixe \cle{-one}, précédé de l'indice de position : \textbf{pentan-2-one}.
\item Citer ensuite les ramifications éventuelles par ordre alphabétique.
\end{enumerate}
\end{methode}

\begin{exemplebox}[Les premières cétones]
\begin{center}
\begin{tabularx}{\linewidth}{|l|X|X|}
\hline
\rowcolor{popgreenL} \textbf{Formule brute} & \textbf{Formule semi-développée} & \textbf{Nom} \\
\hline
\ce{C3H6O} & \ce{CH3-CO-CH3} & \textbf{propanone} (acétone) \\
\hline
\ce{C4H8O} & \ce{CH3-CO-CH2-CH3} & \textbf{butanone} \\
\hline
\ce{C5H10O} & \ce{CH3-CO-CH2-CH2-CH3} & \textbf{pentan-2-one} \\
\hline
\ce{C5H10O} & \ce{CH3-CH2-CO-CH2-CH3} & \textbf{pentan-3-one} \\
\hline
\end{tabularx}
\end{center}

\textbf{Deux remarques essentielles sur ces exemples.}
\begin{itemize}
\item La propanone n'a \textbf{pas besoin d'indice} : sur une chaîne de 3 carbones, le carbonyle ne peut être qu'en position 2 (les positions 1 et 3 donneraient un aldéhyde). De même, « butanone » suffit car butan-2-one et « butan-3-one » désignent la même molécule (il suffit de numéroter dans l'autre sens). Par convention, on omet l'indice lorsqu'il n'y a qu'une seule position possible.
\item À partir de 5 carbones, l'indice devient \textbf{indispensable} : pentan-2-one et pentan-3-one sont deux molécules différentes (ce sont des isomères de position).
\end{itemize}
\end{exemplebox}

\begin{attention}[Pentan-2-one, jamais pentan-4-one]
Prenons la molécule \ce{CH3-CH2-CH2-CO-CH3}. Si tu numérotes la chaîne à partir de la gauche, le carbonyle porte l'indice 4 ; à partir de la droite, il porte l'indice 2. La règle est sans appel : \textbf{on choisit toujours l'indice le plus petit}. Le nom correct est donc \cle{pentan-2-one}. Écrire « pentan-4-one » est une erreur de nomenclature sanctionnée à l'examen, même si la molécule décrite est la bonne.
\end{attention}

\courssub{Isomérie de fonction entre aldéhydes et cétones}

Rappelle-toi : les aldéhydes et les cétones saturés non cycliques partagent la même formule générale $C_nH_{2n}O$. Il existe donc, pour une même formule brute, des aldéhydes \emph{et} des cétones : ce sont des \cle{isomères de fonction}, c'est-à-dire des composés de même formule brute mais de \textbf{groupes fonctionnels différents}.

\begin{exemplebox}[Le couple propanal / propanone : \ce{C3H6O}]
La formule brute \ce{C3H6O} correspond à deux composés :
\begin{itemize}
\item le \textbf{propanal} \ce{CH3-CH2-CHO} : carbonyle en extrémité (aldéhyde) ;
\item la \textbf{propanone} \ce{CH3-CO-CH3} : carbonyle à l'intérieur (cétone).
\end{itemize}
\end{exemplebox}

\begin{popfigure}
\begin{center}
\chemfig{H_3C-CH_2-C(=[2]O)(-[6]H)}
\hspace{2.5cm}
\chemfig{H_3C-C(=[2]O)(-[6]CH_3)-CH_3}
\end{center}
\vspace{0.3cm}
\begin{center}
\textbf{Propanal (aldéhyde)} \hspace{2.2cm} \textbf{Propanone (cétone)}
\end{center}
\legende{Isomérie de fonction : le propanal et la propanone ont la même formule brute \ce{C3H6O}, mais le groupe carbonyle $\ce{C=O}$ est en extrémité de chaîne dans l'aldéhyde (le carbone du carbonyle porte un H) et à l'intérieur de la chaîne dans la cétone (le carbone du carbonyle porte deux groupes méthyle).}
\end{popfigure}

\begin{exoresolu}[Dénombrer les isomères de formule brute \ce{C4H8O}]
Énoncé. Donner les formules semi-développées et les noms de tous les aldéhydes et cétones de formule brute \ce{C4H8O}.
\tcblower
\textbf{Solution détaillée.}

\textbf{Étape 1 : vérifier la formule générale.} Pour $n = 4$ : $C_4H_{2\times 4}O = \ce{C4H8O}$. La formule correspond bien à des aldéhydes et cétones saturés non cycliques.

\textbf{Étape 2 : chercher les aldéhydes} (groupe $-$CHO en bout de chaîne, le carbone du CHO est le n°1).
\begin{itemize}
\item Chaîne linéaire de 4 carbones : \ce{CH3-CH2-CH2-CHO} $\rightarrow$ \textbf{butanal} ;
\item Chaîne ramifiée de 3 carbones + 1 méthyle : \ce{(CH3)2CH-CHO}, c'est-à-dire \chemfig{CH_3-CH(-[6]CH_3)-CHO} $\rightarrow$ \textbf{2-méthylpropanal}.
\end{itemize}
Il n'y a pas d'autre possibilité : avec 4 carbones au total, la chaîne principale contenant le CHO compte 4 ou 3 carbones.

\textbf{Étape 3 : chercher les cétones} (groupe $-$CO$-$ à l'intérieur de la chaîne).
\begin{itemize}
\item \ce{CH3-CO-CH2-CH3} $\rightarrow$ \textbf{butanone}.
\end{itemize}
C'est la seule : une chaîne de 4 carbones n'offre qu'une position interne (la position 2, équivalente à la position 3 par symétrie de numérotation), et une chaîne ramifiée de 3 carbones ne peut pas porter de carbonyle interne entouré de deux alkyles.

\textbf{Conclusion.} \ce{C4H8O} correspond à \textbf{2 aldéhydes} (butanal et 2-méthylpropanal) et \textbf{1 cétone} (butanone), soit 3 isomères au total. Le butanal et la butanone sont isomères de fonction ; le butanal et le 2-méthylpropanal sont isomères de chaîne.
\end{exoresolu}

\begin{methode}[Stratégie générale pour dénombrer les isomères $C_nH_{2n}O$]
\begin{enumerate}
\item \textbf{Aldéhydes} : fixer $-$CHO en position 1, puis faire varier le squelette des $n-1$ carbones restants (isomérie de chaîne).
\item \textbf{Cétones} : placer le groupe $-$CO$-$ à chaque position interne possible de chaque squelette (isomérie de position), en éliminant les doublons par symétrie.
\item Vérifier que chaque formule développée compte bien $n$ carbones, $2n$ hydrogènes et 1 oxygène.
\end{enumerate}
\end{methode}

\courssub{Du quotidien au laboratoire : pourquoi ces molécules comptent}

\begin{savaistu}[L'acétone et l'éthanal, deux stars du quotidien]
La \textbf{propanone}, connue sous le nom courant d'\cle{acétone}, est la cétone la plus simple et l'un des solvants organiques les plus utilisés au monde : c'est elle qui dissout le vernis à ongles dans les dissolvants que l'on trouve dans tous les marchés de Yaoundé et de Douala. Elle sert aussi à nettoyer la verrerie de laboratoire car elle évapore très vite.

L'\textbf{éthanal}, lui, est le premier produit de l'oxydation de l'éthanol. C'est lui qui donne son odeur piquante caractéristique au vin de palme ou au vin qui « tourne » avant de devenir vinaigre : l'éthanol est d'abord oxydé en éthanal, puis l'éthanal est oxydé à son tour en acide éthanoïque. Chez l'être humain, c'est aussi l'éthanal, formé dans le foie lors de la dégradation de l'alcool, qui est en grande partie responsable des effets désagréables de l'excès de boisson. Le \textbf{méthanal} (formaldéhyde), enfin, est utilisé en solution aqueuse (le formol) pour conserver les pièces anatomiques et les spécimens de biologie.
\end{savaistu}

\begin{aretenir}[L'essentiel de la section]
\begin{itemize}
\item Aldéhydes et cétones possèdent le \textbf{groupe carbonyle} $\ce{C=O}$ ; formule générale commune $C_nH_{2n}O$ (composés saturés non cycliques).
\item \textbf{Aldéhyde} R$-$CHO : carbonyle en \textbf{extrémité} de chaîne, le carbone du CHO porte un H. Suffixe \textbf{-al}, \textbf{jamais d'indice de position} (le CHO est toujours le carbone n°1) : méthanal, éthanal, propanal, butanal.
\item \textbf{Cétone} R$-$CO$-$R' : carbonyle à l'\textbf{intérieur} de la chaîne. Suffixe \textbf{-one} avec l'\textbf{indice le plus petit possible} : propanone, butanone, pentan-2-one, pentan-3-one.
\item Aldéhydes et cétones de même formule brute sont des \textbf{isomères de fonction} (ex. propanal et propanone, \ce{C3H6O}).
\end{itemize}
\end{aretenir}

Maintenant que tu sais reconnaître et nommer ces deux familles, une question se pose naturellement : comment, au laboratoire, distinguer concrètement un aldéhyde d'une cétone versés dans deux tubes non étiquetés ? C'est précisément l'objet de la section suivante, consacrée aux tests d'identification : la 2,4-DNPH, le réactif de Schiff, la liqueur de Fehling et le réactif de Tollens.

\coursec{Tests d'identification des aldéhydes et cétones}

Dans la section précédente, nous avons appris à reconnaître et à nommer les aldéhydes et les cétones à partir de leurs formules. Mais dans un laboratoire réel — par exemple au contrôle qualité de la SONARA ou dans une brasserie — on reçoit parfois un flacon sans étiquette. Comment savoir si le liquide contient un aldéhyde, une cétone, ou aucun des deux ? La réponse passe par des \cle{tests d'identification} : des réactions chimiques simples, rapides, qui produisent un changement visible (coloration, précipité, dépôt métallique) caractéristique d'une fonction. C'est l'objet de cette section, capitale pour le Probatoire.

\courssub{Le test à la 2,4-DNPH : caractériser le groupe carbonyle}

\textbf{Intuition.} Aldéhydes et cétones possèdent tous les deux un groupe \cle{carbonyle} $\ce{C=O}$. Il existe donc un réactif qui « reconnaît » ce groupe, sans distinguer les deux familles : c'est la 2,4-DNPH.

\begin{definition}[La 2,4-dinitrophénylhydrazine (2,4-DNPH)]
La \cle{2,4-DNPH} (ou 2,4-dinitrophénylhydrazine) est un réactif qui donne, avec tout composé porteur d'un groupe carbonyle $\ce{C=O}$ (aldéhyde \textbf{ou} cétone), un \cle{précipité jaune-orangé} de 2,4-dinitrophénylhydrazone. Elle \textbf{caractérise le groupe carbonyle} mais \textbf{ne distingue pas} un aldéhyde d'une cétone.
\end{definition}

\textbf{Protocole.} On introduit quelques gouttes du composé à tester dans un tube à essai, puis on ajoute quelques gouttes de solution de 2,4-DNPH. L'apparition d'un précipité jaune-orangé (souvent instantané, parfois après léger chauffage) prouve la présence d'un groupe carbonyle. Si rien ne se passe, le composé n'est ni un aldéhyde ni une cétone (c'est peut-être un alcool, un acide carboxylique, un ester...).

\begin{attention}[Ce que la DNPH ne dit pas]
Un précipité jaune-orangé avec la 2,4-DNPH signifie seulement : « il y a un groupe carbonyle ». Il est \textbf{impossible} de conclure « c'est un aldéhyde » ou « c'est une cétone » sur la base de ce seul test. Il faut un deuxième test (Schiff, Fehling ou Tollens) pour trancher.
\end{attention}

\courssub{Le réactif de Schiff : un test spécifique des aldéhydes}

\textbf{Intuition.} Pour distinguer aldéhydes et cétones, il faut exploiter une différence de réactivité. Le réactif de Schiff est le premier test qui le permet.

\begin{propriete}[Test au réactif de Schiff]
Le \cle{réactif de Schiff} (fuchsine décolorée par l'acide sulfureux $\ce{SO2}$) est incolore. En présence d'un \textbf{aldéhyde}, il se colore en \cle{rose-violet} (magenta). Avec une \textbf{cétone}, le réactif \textbf{reste incolore} : le test est négatif.
\end{propriete}

\textbf{Protocole.} Quelques gouttes du composé à tester + quelques gouttes de réactif de Schiff, à température ambiante (à froid). Coloration rose-violette $\Rightarrow$ aldéhyde ; pas de coloration $\Rightarrow$ cétone (ou absence de carbonyle, déjà écartée par la DNPH).

\begin{attention}[Sensibilité du réactif de Schiff]
Le réactif de Schiff est très sensible : il peut rosir légèrement au contact de l'air (oxydation lente) ou avec certains composés non carbonylés (acides forts, $\ce{SO2}$ résiduel). Un test positif doit être une coloration \textbf{franche et rapide} (quelques secondes). Une coloration lente et diffuse après plusieurs minutes n'est pas significative.
\end{attention}

\courssub{La liqueur de Fehling : l'oxydation ménagée des aldéhydes}

\textbf{Intuition.} Les aldéhydes sont des espèces \cle{réductrices} : ils peuvent être oxydés facilement en acides carboxyliques (ou en ions carboxylates en milieu basique). Les cétones, elles, ne subissent pas cette oxydation ménagée. La liqueur de Fehling met en évidence cette différence de comportement rédox.

\begin{propriete}[Test à la liqueur de Fehling]
La \cle{liqueur de Fehling} est une solution \textbf{bleue} contenant des ions cuivre(II) $\ce{Cu^2+}$ complexés par l'ion tartrate (sel de Rochelle) en milieu basique (soude). Chauffée en présence d'un \textbf{aldéhyde}, elle donne un \cle{précipité rouge brique} d'oxyde de cuivre(I) $\ce{Cu2O}$. Avec une \textbf{cétone}, la solution \textbf{reste bleue} : test négatif.
\end{propriete}

\textbf{Interprétation rédox.} L'aldéhyde $\ce{R-CHO}$ est oxydé en ion carboxylate $\ce{R-COO-}$ (couple $\ce{R-COO-}/\ce{R-CHO}$), tandis que les ions $\ce{Cu^2+}$ sont réduits en $\ce{Cu2O}$ (le cuivre passe du degré d'oxydation $+II$ à $+I$). L'équation-bilan en milieu basique s'écrit :

\[ \ce{R-CHO + 2Cu^2+ + 5OH- -> R-COO- + Cu2O + 3H2O} \]

\begin{methode}[Démonstration de l'équilibrage (rappel méthodologique)]
\begin{enumerate}
\item \textbf{Demi-équation d'oxydation} (aldéhyde $\to$ carboxylate) : $\ce{R-CHO + 3OH- -> R-COO- + 2H2O + 2e-}$
\item \textbf{Demi-équation de réduction} ($\ce{Cu^2+} \to \ce{Cu2O}$) : $\ce{2Cu^2+ + 2OH- + 2e- -> Cu2O + H2O}$
\item \textbf{Somme} : on additionne membre à membre et on simplifie l'eau et les $\ce{OH-}$.
\end{enumerate}
\end{methode}

\textbf{Conditions :} chauffage au bain-marie (environ \SI{60}{\celsius} à \SI{80}{\celsius}), milieu basique. Le précipité rouge brique de $\ce{Cu2O}$ est la signature du test positif.

\begin{attention}[Test négatif $\neq$ « sans action » mal formulée]
Dire que la liqueur de Fehling est « sans action » sur une cétone signifie que le \textbf{test est négatif} : la solution \textbf{reste bleue}, aucun précipité n'apparaît, rien d'observable ne se produit. Au Probatoire, rédige toujours l'observation concrète (« la solution reste bleue, pas de précipité rouge brique ») plutôt qu'une formule vague.
\end{attention}

\courssub{Le réactif de Tollens : le miroir d'argent}

\begin{propriete}[Test au réactif de Tollens]
Le \cle{réactif de Tollens} (nitrate d'argent ammoniacal) contient des ions diammineargent(I) $\ce{[Ag(NH3)2]+}$. En présence d'un \textbf{aldéhyde}, à tiède, il se forme un \cle{dépôt d'argent métallique} $\ce{Ag}$ qui tapisse la paroi du tube : c'est le fameux \cle{miroir d'argent}. Avec une \textbf{cétone}, le test est négatif : pas de dépôt.
\end{propriete}

\textbf{Interprétation rédox.} Là encore, l'aldéhyde est oxydé en ion carboxylate, et l'ion $\ce{Ag+}$ (complexé) est réduit en argent métallique $\ce{Ag}$ :

\[ \ce{R-CHO + 2[Ag(NH3)2]+ + 3OH- -> R-COO- + 2Ag + 4NH3 + 2H2O} \]

\begin{methode}[Démonstration de l'équilibrage (rappel méthodologique)]
\begin{enumerate}
\item \textbf{Oxydation} : $\ce{R-CHO + 3OH- -> R-COO- + 2H2O + 2e-}$
\item \textbf{Réduction} : $\ce{2[Ag(NH3)2]+ + 2e- -> 2Ag + 4NH3}$
\item \textbf{Somme} : on ajoute $3\ce{OH-}$ à gauche pour neutraliser l'acidité latente et former l'eau (équilibrage O/H en milieu basique).
\end{enumerate}
\end{methode}

\textbf{Conditions :} chauffage doux au bain-marie, milieu ammoniacal basique. Le miroir gris-brillant sur la paroi du tube est inratable.

\begin{savaistu}[Pourquoi « miroir » d'argent ?]
Avant l'industrialisation moderne, le test de Tollens était réellement utilisé pour \textbf{argenter les miroirs} : on versait sur une plaque de verre un mélange de réactif de Tollens et d'un aldéhyde (comme le glucose, qui possède une fonction aldéhyde), et une fine couche d'argent métallique se déposait uniformément, rendant le verre réfléchissant. C'est de cette application artisanale que vient le nom « miroir d'argent ». Le même principe sert encore à argenter les boules de Noël et certains thermos !
\end{savaistu}

\courssub{Tableau récapitulatif et stratégie d'identification}

\begin{aretenir}[Résultats des quatre tests]
\begin{center}
\begin{tabularx}{\linewidth}{|l|X|X|}
\hline
\rowcolor{popblueL} \textbf{Test} & \textbf{Aldéhyde} & \textbf{Cétone} \\
\hline
2,4-DNPH & Précipité jaune-orangé & Précipité jaune-orangé \\
\hline
Réactif de Schiff & Coloration rose-violette & Reste incolore (négatif) \\
\hline
Liqueur de Fehling (à chaud) & Précipité rouge brique $\ce{Cu2O}$ & Solution bleue inchangée (négatif) \\
\hline
Réactif de Tollens & Miroir d'argent $\ce{Ag}$ & Pas de dépôt (négatif) \\
\hline
\end{tabularx}
\end{center}
\textbf{Conclusion clé :} les aldéhydes sont des \cle{réducteurs} (ils réduisent $\ce{Cu^2+}$ et $\ce{Ag+}$) ; les cétones ne le sont pas.
\end{aretenir}

\begin{popfigure}
\begin{tikzpicture}[scale=0.95]
% Tube 1 : DNPH
\draw[thick,popdark] (0,0) -- (0,2.6) (0.9,0) -- (0.9,2.6);
\draw[thick,popdark] (0,0) arc (180:360:0.45);
\fill[popgold!80] (0.02,0.02) arc (180:360:0.43) -- (0.88,0.55) -- (0.02,0.55) -- cycle;
\fill[popgold!30] (0.02,0.55) rectangle (0.88,1.5);
\node[align=center,font=\small] at (0.45,-0.7) {DNPH\\ précipité\\ jaune-orangé};
% Tube 2 : Schiff
\begin{scope}[xshift=2.2cm]
\draw[thick,popdark] (0,0) -- (0,2.6) (0.9,0) -- (0.9,2.6);
\draw[thick,popdark] (0,0) arc (180:360:0.45);
\fill[poppink!60] (0.02,0.02) arc (180:360:0.43) -- (0.88,1.5) -- (0.02,1.5) -- cycle;
\node[align=center,font=\small] at (0.45,-0.7) {Schiff\\ coloration\\ rose-violette};
\end{scope}
% Tube 3 : Fehling
\begin{scope}[xshift=4.4cm]
\draw[thick,popdark] (0,0) -- (0,2.6) (0.9,0) -- (0.9,2.6);
\draw[thick,popdark] (0,0) arc (180:360:0.45);
\fill[poporange!85] (0.02,0.02) arc (180:360:0.43) -- (0.88,0.5) -- (0.02,0.5) -- cycle;
\fill[popblue!40] (0.02,0.5) rectangle (0.88,1.5);
\node[align=center,font=\small] at (0.45,-0.7) {Fehling\\ précipité\\ rouge brique};
\end{scope}
% Tube 4 : Tollens
\begin{scope}[xshift=6.6cm]
\draw[thick,popdark] (0,0) -- (0,2.6) (0.9,0) -- (0.9,2.6);
\draw[thick,popdark] (0,0) arc (180:360:0.45);
\fill[popmuted!40] (0.02,0.02) arc (180:360:0.43) -- (0.88,1.5) -- (0.02,1.5) -- cycle;
\draw[line width=1.5pt,popmuted] (0.06,0.1) -- (0.06,1.45);
\draw[line width=1.5pt,popmuted] (0.84,0.1) -- (0.84,1.45);
\node[align=center,font=\small] at (0.45,-0.7) {Tollens\\ miroir\\ d'argent};
\end{scope}
\end{tikzpicture}
\legende{Aspect des quatre tubes à essai lors de tests positifs sur un aldéhyde : précipité jaune-orangé (DNPH), coloration rose-violette (Schiff), précipité rouge brique de $\ce{Cu2O}$ sous solution bleue (Fehling), dépôt d'argent sur les parois (Tollens).}
\end{popfigure}

\begin{methode}[Identifier un composé inconnu : la stratégie en deux étapes]
\begin{enumerate}
\item \textbf{Test à la 2,4-DNPH.} Précipité jaune-orangé $\Rightarrow$ le composé contient un groupe carbonyle (aldéhyde ou cétone). Pas de précipité $\Rightarrow$ ni aldéhyde ni cétone : inutile d'aller plus loin.
\item \textbf{Test au réactif de Schiff (ou Fehling, ou Tollens).} Test positif $\Rightarrow$ c'est un \textbf{aldéhyde}. Test négatif $\Rightarrow$ c'est une \textbf{cétone}.
\item \textbf{Conclure} en citant les deux observations, puis, si la formule brute est connue, proposer la formule semi-développée et le nom du composé.
\end{enumerate}
\textbf{Piège :} un seul test positif à la DNPH ne suffit jamais à conclure entre aldéhyde et cétone.
\end{methode}

\begin{popfigure}
\begin{tikzpicture}[
  node distance=0.55cm,
  boite/.style={rectangle, rounded corners=3pt, draw=popblue, fill=popblueL, thick, align=center, font=\small, text width=3.4cm, minimum height=0.9cm},
  test/.style={rectangle, rounded corners=3pt, draw=poporange, fill=poporangeL, thick, align=center, font=\small, text width=3.4cm, minimum height=0.9cm},
  concl/.style={rectangle, rounded corners=3pt, draw=popgreen, fill=popgreenL, thick, align=center, font=\small, text width=3.4cm, minimum height=0.9cm},
  fleche/.style={-{Stealth[length=2.5mm]}, thick, popdark}
]
\node[boite] (inc) {Composé inconnu à identifier};
\node[test, below=of inc] (dnph) {Test à la 2,4-DNPH};
\node[concl, below left=0.7cm and -1.2cm of dnph] (pas) {Pas de précipité :\\ ni aldéhyde ni cétone};
\node[test, below right=0.7cm and -1.2cm of dnph] (schiff) {Précipité jaune-orangé :\\ test au réactif de Schiff (ou Fehling/Tollens)};
\node[concl, below left=0.7cm and -1.4cm of schiff] (ald) {Test positif :\\ \textbf{aldéhyde}};
\node[concl, below right=0.7cm and -1.4cm of schiff] (cet) {Test négatif :\\ \textbf{cétone}};
\draw[fleche] (inc) -- (dnph);
\draw[fleche] (dnph) -- (pas) node[midway, above left, font=\scriptsize]{négatif};
\draw[fleche] (dnph) -- (schiff) node[midway, above right, font=\scriptsize]{positif};
\draw[fleche] (schiff) -- (ald) node[midway, above left, font=\scriptsize]{positif};
\draw[fleche] (schiff) -- (cet) node[midway, above right, font=\scriptsize]{négatif};
\end{tikzpicture}
\legende{Organigramme d'identification d'un composé carbonylé inconnu : la DNPH détecte le groupe carbonyle, puis le réactif de Schiff (ou Fehling, ou Tollens) distingue l'aldéhyde de la cétone.}
\end{popfigure}

\begin{exoresolu}[Identification d'un composé de formule brute \ce{C3H6O}]
Un composé organique A a pour formule brute \ce{C3H6O}. Il donne un précipité jaune-orangé avec la 2,4-DNPH. Chauffé avec la liqueur de Fehling, il provoque l'apparition d'un précipité rouge brique. Identifier A, écrire sa formule semi-développée et l'équation-bilan du test de Fehling.
\tcblower
\textbf{Étape 1 — Exploiter le test à la DNPH.} Le précipité jaune-orangé prouve la présence d'un groupe carbonyle : A est un aldéhyde ou une cétone. Avec \ce{C3H6O}, les deux candidats sont le \cle{propanal} \ce{CH3-CH2-CHO} et la \cle{propanone} \ce{CH3-CO-CH3}.

\textbf{Étape 2 — Exploiter le test de Fehling.} Le précipité rouge brique de $\ce{Cu2O}$ est un test \textbf{positif} : A est un réducteur, donc un \textbf{aldéhyde}. La propanone (cétone) aurait laissé la solution bleue inchangée.

\textbf{Conclusion :} A est le \cle{propanal} : \chemfig{CH_3-CH_2-C(=[1]O)-[7]H}.

\textbf{Équation-bilan du test (milieu basique, chauffage) :}
\[ \ce{CH3-CH2-CHO + 2Cu^2+ + 5OH- -> CH3-CH2-COO- + Cu2O + 3H2O} \]
L'aldéhyde est oxydé en ion propanoate ; les ions $\ce{Cu^2+}$ sont réduits en oxyde de cuivre(I) rouge brique.
\end{exoresolu}

\begin{exercice}[À toi de jouer]
Un liquide B de formule brute \ce{C4H8O} donne un précipité jaune-orangé avec la 2,4-DNPH, mais ne rosit pas le réactif de Schiff et ne réagit pas avec la liqueur de Fehling. Donner la nature de B, sa formule semi-développée et son nom.
\end{exercice}

\begin{corrige}[Exercice 1]
Le test à la DNPH est positif : B contient un groupe carbonyle. Les tests de Schiff et de Fehling sont négatifs : B n'est pas un aldéhyde, c'est donc une \textbf{cétone}. La seule cétone acyclique saturée de formule brute \ce{C4H8O} est la \cle{butan-2-one} (anciennement butanone) : \chemfig{CH_3-C(=[1]O)-[7]CH_2-CH_3}.
\end{corrige}

\begin{aretenir}[L'essentiel de la section]
\begin{itemize}
\item La \cle{2,4-DNPH} caractérise le groupe carbonyle : précipité jaune-orangé avec les aldéhydes \textbf{et} les cétones.
\item Le \cle{réactif de Schiff} (rose-violet), la \cle{liqueur de Fehling} (précipité rouge brique $\ce{Cu2O}$ à chaud) et le \cle{réactif de Tollens} (miroir d'argent) sont positifs avec les \textbf{aldéhydes uniquement}.
\item Les aldéhydes sont des \cle{réducteurs} (oxydables en carboxylates) ; les cétones ne le sont pas.
\item Stratégie : DNPH d'abord, puis Schiff, Fehling ou Tollens pour trancher.
\end{itemize}
\end{aretenir}

\coursec{Synthèse et applications au quotidien}

Dans la section précédente, tu as appris à distinguer expérimentalement les aldéhydes des cétones grâce à la 2,4-DNPH, au réactif de Schiff, à la liqueur de Fehling et au réactif de Tollens. Ces tests ne sont pas de simples recettes de laboratoire : ils sont au cœur d'une démarche plus vaste qui consiste à \cle{identifier} un composé organique inconnu, démarche très appréciée au Probatoire. Mais une question demeure : d'où viennent ces aldéhydes et ces cétones ? La réponse est élégante : ils naissent de l'\cle{oxydation ménagée} des alcools. Cette dernière section boucle la boucle : elle relie les trois familles oxygénées entre elles, dresse un bilan complet de la leçon, puis montre comment ces molécules irriguent la vie quotidienne camerounaise, de la pharmacie au savon artisanal, en passant par les carburants.

\courssub{Bilan général : les familles oxygénées en un coup d'œil}

Avant d'aller plus loin, prenons de la hauteur. Trois familles de composés oxygénés ont été étudiées dans cette leçon : les \cle{alcools}, les \cle{aldéhydes} et les \cle{cétones}. Chacune possède un groupe fonctionnel caractéristique, une formule générale, un suffixe de nomenclature et des tests d'identification propres. Le tableau suivant condense l'essentiel de la leçon : apprends-le par cœur, il résume à lui seul la moitié des questions de Probatoire sur ce chapitre.

\begin{center}
\begin{tabularx}{\linewidth}{|l|X|X|X|}
\hline
\rowcolor{popblueL} \textbf{Critère} & \textbf{Alcool} & \textbf{Aldéhyde} & \textbf{Cétone} \\
\hline
Groupe fonctionnel & hydroxyle $-\ce{OH}$ & carbonyle $-\ce{CHO}$ (en bout de chaîne) & carbonyle $\ce{C=O}$ (à l'intérieur) \\
\hline
Formule générale & $C_nH_{2n+1}OH$ (ou $C_nH_{2n+2}O$) & $C_nH_{2n}O$ & $C_nH_{2n}O$ \\
\hline
Suffixe / préfixe & -ol & -al & -one \\
\hline
Exemple & éthanol \ce{CH3-CH2-OH} & éthanal \ce{CH3-CHO} & propanone \ce{CH3-CO-CH3} \\
\hline
Test 2,4-DNPH & négatif & positif (précipité jaune-orangé) & positif (précipité jaune-orangé) \\
\hline
Réactif de Schiff & négatif & positif (coloration rose-violette) & négatif \\
\hline
Liqueur de Fehling & négatif & positif (précipité rouge brique) & négatif \\
\hline
Réactif de Tollens & négatif & positif (miroir d'argent) & négatif \\
\hline
\end{tabularx}
\end{center}

\begin{aretenir}[La stratégie d'identification en deux temps]
La 2,4-DNPH détecte la \textbf{famille des composés carbonylés} (aldéhydes \emph{et} cétones confondus). Les réactifs de Schiff, de Fehling et de Tollens tranchent ensuite \textbf{entre aldéhyde et cétone} : seul l'aldéhyde, qui possède un hydrogène porté par le carbone du carbonyle, est un réducteur capable de les faire réagir.
\end{aretenir}

\courssub{Complément Probatoire : de l'alcool au composé carbonylé, la chaîne d'oxydation}

Voici l'intuition. Un alcool possède un atome d'hydrogène « fragile » : celui du groupe $-\ce{OH}$, mais aussi celui (ou ceux) porté(s) par le carbone fonctionnel. Un \cle{oxydant ménagé} (dichromate de potassium \ce{K2Cr2O7} en milieu acide, ou permanganate de potassium \ce{KMnO4} dilué) peut arracher ces hydrogènes \emph{sans casser la chaîne carbonée}. Le résultat dépend alors crucialement de la \textbf{classe} de l'alcool, c'est-à-dire du nombre de groupes alkyle portés par le carbone fonctionnel.

\begin{propriete}[Oxydation ménagée des alcools — complément Probatoire]
\begin{itemize}
\item Un \textbf{alcool primaire} $R-CH_2-OH$ s'oxyde d'abord en \textbf{aldéhyde} $R-CHO$ ; si l'oxydant est en excès ou si l'aldéhyde n'est pas éliminé du milieu au fur et à mesure, l'oxydation se poursuit jusqu'à l'\textbf{acide carboxylique} $R-COOH$.
\item Un \textbf{alcool secondaire} $R_1-CH(OH)-R_2$ s'oxyde en \textbf{cétone} $R_1-CO-R_2$, et l'oxydation s'arrête là.
\item Un \textbf{alcool tertiaire} ne subit \textbf{pas d'oxydation ménagée} : son carbone fonctionnel ne porte aucun atome d'hydrogène à arracher. (En conditions brutales, la chaîne carbonée se rompt : ce n'est plus une oxydation ménagée.)
\end{itemize}
\end{propriete}

\begin{popfigure}
\begin{center}
\begin{tikzpicture}[node distance=0.5cm]
\node (a) {\chemfig{CH_3-CH_2-OH}};
\node[below=0.15cm of a, font=\small\itshape, text=popdark] {éthanol (primaire)};
\node[right=2.6cm of a] (b) {\chemfig{CH_3-C(=[1]O)-[7]H}};
\node[below=0.15cm of b, font=\small\itshape, text=popdark] {éthanal (aldéhyde)};
\node[right=2.6cm of b] (c) {\chemfig{CH_3-C(=[1]O)-[7]OH}};
\node[below=0.15cm of c, font=\small\itshape, text=popdark] {acide éthanoïque};
\draw[-{Stealth[length=3mm]}, line width=1.2pt, poporange] ($(a.east)+(0.25,0.35)$) -- node[above, font=\small, text=popdark]{oxydation ménagée} node[below, font=\small, text=popmuted]{\ce{Cr2O7^2-}/\ce{H+}} ($(b.west)+(-0.25,0.35)$);
\draw[-{Stealth[length=3mm]}, line width=1.2pt, poporange] ($(b.east)+(0.25,0.35)$) -- node[above, font=\small, text=popdark]{oxydation} node[below, font=\small, text=popmuted]{oxydant en excès} ($(c.west)+(-0.25,0.35)$);
\end{tikzpicture}
\end{center}
\legende{Chaîne d'oxydation ménagée d'un alcool primaire : alcool $\rightarrow$ aldéhyde $\rightarrow$ acide carboxylique. Pour un alcool secondaire, la chaîne s'arrête à la cétone ; pour un alcool tertiaire, elle ne démarre pas.}
\end{popfigure}

Écrivons les équations-bilans avec l'oxydant le plus classique, l'ion dichromate \ce{Cr2O7^2-} en milieu acide (couple \ce{Cr2O7^2-}/\ce{Cr^3+}). Pour l'alcool primaire, par exemple l'éthanol :

\[ \ce{3 CH3-CH2-OH + Cr2O7^2- + 8 H+ -> 3 CH3-CHO + 2 Cr^3+ + 7 H2O} \]

L'aldéhyde formé peut être oxydé à son tour :

\[ \ce{3 CH3-CHO + Cr2O7^2- + 8 H+ -> 3 CH3-COOH + 2 Cr^3+ + 4 H2O} \]

Pour un alcool secondaire, par exemple le propan-2-ol :

\[ \ce{3 CH3-CH(OH)-CH3 + Cr2O7^2- + 8 H+ -> 3 CH3-CO-CH3 + 2 Cr^3+ + 7 H2O} \]

\textbf{Observation expérimentale caractéristique :} la solution passe de l'\textbf{orange} (couleur de l'ion dichromate) au \textbf{vert} (couleur de l'ion chrome \ce{Cr^3+}). Ce virage de couleur est d'ailleurs le principe de l'\emph{éthylotest} chimique autrefois utilisé par les forces de l'ordre : l'éthanol expiré par un conducteur ayant bu décolore le dichromate.

\begin{attention}[Ne confonds pas oxydation ménagée et combustion]
L'oxydation ménagée \textbf{conserve le squelette carboné} : elle transforme seulement le groupe fonctionnel. La combustion, elle, détruit totalement la molécule en \ce{CO2} et \ce{H2O}. Dire « l'alcool tertiaire ne s'oxyde pas » est un abus de langage : il ne subit pas d'oxydation \emph{ménagée}, mais il brûle parfaitement !
\end{attention}

\begin{attention}[Le piège classique du test DNPH]
Erreur fréquente au Probatoire : un élève réalise le test à la 2,4-DNPH, observe un précipité jaune-orangé et conclut « le composé est un aldéhyde ». C'est \textbf{faux} : la DNPH réagit avec les aldéhydes \emph{et} les cétones. Le test DNPH positif prouve seulement la présence d'un \textbf{groupe carbonyle}. Pour conclure, il faut un second test : Schiff, Fehling ou Tollens. Positif $\rightarrow$ aldéhyde ; négatif $\rightarrow$ cétone.
\end{attention}

\courssub{Applications au quotidien : des molécules qui changent la vie}

\courssub{L'éthanol, l'alcool aux mille visages}

L'\cle{éthanol} \ce{CH3-CH2-OH} est sans conteste le composé oxygéné le plus présent dans la vie quotidienne.

\begin{itemize}
\item \textbf{Désinfection et hygiène.} L'alcool à \SI{70}{\%} (v/v) (70 mL d'éthanol pour 100 mL de solution) est un antiseptique de référence dans les hôpitaux et les pharmacies de Yaoundé ou de Douala : il détruit les membranes des bactéries et de nombreux virus. Les solutions hydro-alcooliques en contiennent de \SI{60}{\%} à \SI{80}{\%} en volume.
\item \textbf{Boissons.} Obtenu par fermentation alcoolique des sucres (action des levures sur le glucose : $\ce{C6H12O6 -> 2 C2H5OH + 2 CO2}$), il est présent dans le vin de palme, le bili-bili (mil fermenté) ou la bière industrielle. La distillation du vin de palme donne l'odontol ou l'afofo, eau-de-vie très concentrée en éthanol.
\item \textbf{Biocarburant.} Mélangé à l'essence (carburant E10 : \SI{10}{\%} d'éthanol), il réduit la consommation de pétrole et les émissions polluantes.
\end{itemize}

\begin{savaistu}[De la canne à sucre au réservoir : la filière SOSUCAM]
La SOSUCAM (Société Sucrière du Cameroun), installée à Nkoteng et Mbandjock, produit du sucre à partir de la canne. La mélasse, résidu sucré de la cristallisation, peut être fermentée puis distillée pour produire de l'éthanol. Dans plusieurs pays (Brésil en tête), cet éthanol « vert » est mélangé à l'essence comme \cle{biocarburant} : il est dit renouvelable car le \ce{CO2} dégagé lors de sa combustion a été préalablement absorbé par la canne pendant sa croissance. Le bilan carbone est donc presque neutre, contrairement à celui de l'essence fossile.
\end{savaistu}

\courssub{Le glycérol et la savonnerie artisanale}

Le \cle{glycérol} (propane-1,2,3-triol, \ce{HO-CH2-CH(OH)-CH2-OH}) est un polyalcool sirupeux, soluble dans l'eau, hygroscopique. Il intervient dans la fabrication du savon : la \textbf{saponification} est la réaction entre un corps gras (ester de glycérol et d'acides gras, issu par exemple de l'huile de palme) et la soude à chaud. Elle produit du savon et libère du glycérol :

\[ \text{corps gras} + 3\,\ce{NaOH} \xrightarrow{\Delta} 3\,\text{savon} + \text{glycérol} \]

Dans les savonneries artisanales qui utilisent l'huile de palme rouge, le glycérol formé reste souvent dans le savon et lui confère ses propriétés adoucissantes. Industriellement, il est récupéré et revendu aux industries pharmaceutique (sirops, suppositoires) et cosmétique (crèmes hydratantes), car il retient l'eau.

\courssub{La propanone (acétone), la cétone industrielle}

La \cle{propanone} \ce{CH3-CO-CH3}, appelée couramment acétone, est la cétone la plus simple. C'est un excellent \textbf{solvant} organique : elle dissout les vernis à ongles (d'où son usage comme dissolvant), les colles, les résines et de nombreuses peintures. L'industrie l'emploie dans la fabrication de plastiques et de médicaments. Au laboratoire, elle sert à rincer et sécher rapidement la verrerie, car elle est très volatile ($\theta_{eb} = \SI{56}{\celsius}$).

\courssub{Le méthanal (formaldéhyde) et le formol}

Le \cle{méthanal} \ce{HCHO} est le plus simple des aldéhydes. Sa solution aqueuse à environ \SI{37}{\%} en masse, appelée \textbf{formol}, est un puissant désinfectant et un agent de conservation : dans les laboratoires d'anatomie et les facultés de médecine, elle permet de conserver les pièces anatomiques et les spécimens biologiques en durcissant les protéines. Le formol sert aussi à la fabrication de colles et de résines (mélaminé des tables de laboratoire).

\begin{attention}[Sécurité : le méthanol, un poison mortel]
Le \textbf{méthanol} \ce{CH3-OH}, alcool à un seul carbone, est \textbf{toxique} : son oxydation dans l'organisme produit du méthanal puis de l'acide méthanoïque, qui attaquent le nerf optique. L'ingestion de quelques millilitres provoque la cécité ; une trentaine de millilitres peut être mortelle. Les \textbf{alcools frelatés} (boissons artisanales mal distillées, ou coupées avec de l'alcool industriel dénaturé contenant du méthanol) provoquent régulièrement des intoxications graves. Ne jamais consommer d'alcool de provenance douteuse, et ne jamais confondre alcool à brûler et alcool alimentaire.
\end{attention}

\begin{attention}[Inflammabilité : alcools et acétone]
L'éthanol, le méthanol et la propanone sont des liquides \textbf{très inflammables} dont les vapeurs forment avec l'air des mélanges explosifs. Au laboratoire : jamais de flamme nue à proximité, chauffage au bain-marie plutôt qu'au bec Bunsen, bouchons refermés immédiatement, stockage loin des sources de chaleur. C'est pourquoi l'alcool à brûler se manipule avec la plus grande prudence dans les cuisines.
\end{attention}

\courssub{Méthodologie : identifier un composé oxygéné au Probatoire}

\begin{methode}[Démarche d'identification d'un composé oxygéné]
Face à un exercice du type « un composé A de formule brute \ce{C_xH_yO}... », procède toujours ainsi :
\begin{enumerate}
\item \textbf{Exploite la formule brute.} Calcule le nombre d'insaturations. Si la formule est $C_nH_{2n}O$, il y a une insaturation : le composé est un carbonylé (aldéhyde ou cétone) à chaîne saturée, ou un alcool/éther insaturé — l'énoncé précise généralement le contexte. Si elle est $C_nH_{2n+2}O$, c'est un alcool ou un éther-oxyde saturé.
\item \textbf{Test à la 2,4-DNPH.} Positif $\rightarrow$ composé carbonylé (aldéhyde \emph{ou} cétone). Négatif $\rightarrow$ ni aldéhyde ni cétone.
\item \textbf{Test de Schiff, Fehling ou Tollens.} Positif $\rightarrow$ aldéhyde ; négatif $\rightarrow$ cétone.
\item \textbf{Exploite les données complémentaires} (oxydation d'un alcool, nombre d'isomères, nom de l'alcool de départ) pour déterminer la formule semi-développée exacte et nommer le composé.
\item \textbf{Conclus proprement} : formule semi-développée + nom systématique.
\end{enumerate}
\end{methode}

\begin{exoresolu}[Identification d'un composé carbonylé (type Probatoire)]
L'oxydation ménagée d'un alcool A de formule brute \ce{C4H10O} donne un composé B. Le composé B donne un précipité jaune-orangé avec la 2,4-DNPH et un précipité rouge brique avec la liqueur de Fehling.
\begin{enumerate}
\item Que peut-on déduire des tests réalisés sur B ?
\item Donner la classe de l'alcool A.
\item Sachant que la chaîne carbonée de A est ramifiée, donner les formules semi-développées et les noms de A et de B.
\item Écrire l'équation-bilan de la réaction de B avec la liqueur de Fehling.
\end{enumerate}
\tcblower
\textbf{1. Interprétation des tests.} Le test positif à la 2,4-DNPH prouve que B possède un \textbf{groupe carbonyle} : B est un aldéhyde ou une cétone. Le test positif à la liqueur de Fehling (précipité rouge brique d'oxyde de cuivre(I) \ce{Cu2O}) prouve que B est un \textbf{réducteur} : c'est donc un \textbf{aldéhyde}.

\textbf{2. Classe de A.} Seul un alcool \textbf{primaire} donne un aldéhyde par oxydation ménagée. A est donc un alcool primaire.

\textbf{3. Formules et noms.} La formule \ce{C4H10O} correspond à $C_nH_{2n+2}O$ avec $n = 4$ : alcool saturé. A étant primaire et à chaîne \emph{ramifiée}, sa formule semi-développée est \chemfig{CH_3-CH(-[2]CH_3)-CH_2-OH} : c'est le \textbf{2-méthylpropan-1-ol}. (L'autre alcool primaire de formule \ce{C4H10O}, le butan-1-ol, a une chaîne linéaire : il est exclu.) L'oxydation conserve le squelette carboné, donc B est \chemfig{CH_3-CH(-[2]CH_3)-CHO} : le \textbf{2-méthylpropanal}.

\textbf{4. Équation-bilan avec la liqueur de Fehling.} En milieu basique, l'ion cuivre(II) \ce{Cu^2+} (complexé, responsable de la couleur bleue) est réduit en oxyde de cuivre(I) \ce{Cu2O} rouge brique, tandis que l'aldéhyde est oxydé en ion carboxylate :

\[ \ce{CH3-CH(CH3)-CHO + 2 Cu^2+ + 5 OH- -> CH3-CH(CH3)-COO- + Cu2O + 3 H2O} \]

On note $R = \ce{CH3-CH(CH3)-}$ pour alléger : $\ce{R-CHO + 2 Cu^2+ + 5 OH- -> R-COO- + Cu2O + 3 H2O}$. Le virage du bleu au rouge brique est la signature expérimentale de l'aldéhyde.
\end{exoresolu}

\begin{exercice}[À toi de jouer]
Un composé organique C a pour formule brute \ce{C3H6O}. Il donne un précipité jaune avec la 2,4-DNPH, mais ne réagit ni avec le réactif de Schiff ni avec la liqueur de Fehling. Par ailleurs, C est obtenu par oxydation ménagée d'un alcool D.
\begin{enumerate}
\item Déterminer la fonction chimique de C, sa formule semi-développée et son nom.
\item Donner la formule semi-développée, le nom et la classe de D.
\item Écrire l'équation-bilan de l'oxydation de D par l'ion dichromate en milieu acide.
\end{enumerate}
\end{exercice}

\begin{corrige}[Exercice : identification de C et D]
\textbf{1.} La formule \ce{C3H6O} est du type $C_nH_{2n}O$ : une insaturation. DNPH positive $\rightarrow$ carbonylé ; Schiff et Fehling négatifs $\rightarrow$ ce n'est \emph{pas} un aldéhyde. C est donc une \textbf{cétone} : \chemfig{CH_3-C(=[1]O)-[7]CH_3}, la \textbf{propanone} (acétone).

\textbf{2.} Une cétone provient de l'oxydation ménagée d'un alcool \textbf{secondaire} : D est \chemfig{CH_3-CH(-[2]OH)-CH_3}, le \textbf{propan-2-ol}, alcool secondaire.

\textbf{3.} Équation-bilan :
\[ \ce{3 CH3-CH(OH)-CH3 + Cr2O7^2- + 8 H+ -> 3 CH3-CO-CH3 + 2 Cr^3+ + 7 H2O} \]
La solution passe de l'orange au vert.
\end{corrige}

\begin{aretenir}[L'essentiel de la section]
\begin{itemize}
\item La 2,4-DNPH détecte les carbonylés (aldéhydes \emph{et} cétones) ; Schiff, Fehling et Tollens distinguent l'aldéhyde (positif) de la cétone (négatif).
\item Oxydation ménagée : alcool \textbf{primaire} $\rightarrow$ aldéhyde $\rightarrow$ acide carboxylique ; alcool \textbf{secondaire} $\rightarrow$ cétone ; alcool \textbf{tertiaire} $\rightarrow$ pas d'oxydation ménagée.
\item Éthanol : antiseptique, boissons fermentées, biocarburant ; glycérol : coproduit de la saponification ; acétone : solvant industriel ; formol : conservateur.
\item Le méthanol est un poison (cécité, mort) ; alcools et acétone sont très inflammables.
\item Méthode d'identification : formule brute $\rightarrow$ DNPH $\rightarrow$ Schiff/Fehling/Tollens $\rightarrow$ conclusion nommée.
\end{itemize}
\end{aretenir}

Tu disposes désormais d'une vision complète des composés oxygénés : leurs structures, leurs noms, leurs propriétés et leurs tests. La leçon suivante te fera découvrir la famille qui prolonge naturellement la chaîne d'oxydation que tu viens d'étudier : les \textbf{acides carboxyliques} et leurs dérivés, les \textbf{esters}, à l'origine des arômes de fruits et des savons.

\newpage

\coursec{Travaux pratiques}

\courssub{Objectifs du TP}

Ce travail pratique te place dans la peau d'un \cle{chimiste analyste} : trois flacons ont perdu leur étiquette dans le laboratoire du lycée, et tu dois retrouver la nature de leur contenu en utilisant uniquement des tests chimiques. C'est exactement la démarche suivie chaque jour dans les laboratoires de contrôle de qualité, par exemple pour vérifier qu'un produit vendu comme « alcool » ne contient pas de solvant dangereux.

À l'issue de ce TP, tu seras capable de :

\begin{itemize}
\item réaliser dans les règles de sécurité les trois tests caractéristiques des composés carbonylés : test à la \cle{2,4-DNPH}, test au \cle{réactif de Schiff} et test à la \cle{liqueur de Fehling} ;
\item interpréter les observations (précipité jaune-orangé, coloration rose-violette, précipité rouge brique) pour distinguer un \cle{aldéhyde} d'une \cle{cétone} et d'un \cle{alcool} ;
\item identifier deux échantillons inconnus à partir d'un tableau d'observations ;
\item écrire l'équation-bilan de la réaction d'oxydation ménagée d'un aldéhyde par la liqueur de Fehling.
\end{itemize}

\begin{aretenir}[Rappels indispensables avant de commencer]
\begin{itemize}
\item La 2,4-DNPH réagit avec \textbf{tous} les composés carbonylés (aldéhydes \emph{et} cétones) : précipité jaune-orangé.
\item Le réactif de Schiff et la liqueur de Fehling ne réagissent qu'avec les \textbf{aldéhydes} (coloration rose-violette ; précipité rouge brique à chaud).
\item Les alcools ne réagissent avec \textbf{aucun} de ces trois réactifs.
\end{itemize}
\end{aretenir}

\courssub{Matériel et produits}

\begin{center}
\begin{tabularx}{\linewidth}{|l|X|}
\hline
\rowcolor{popblueL} \textbf{Matériel} & \textbf{Produits} \\
\hline
9 tubes à essai propres et secs & Échantillon A : éthanal (ou formol dilué à 5 \%) \\
\hline
3 pipettes compte-gouttes & Échantillon B : propanone (acétone commerciale) \\
\hline
Portoir à tubes & Échantillon C : éthanol \\
\hline
Bécher de \SI{250}{mL} (bain-marie) & Solution de 2,4-DNPH (quelques mL) \\
\hline
Bec Bunsen ou réchaud à gaz & Réactif de Schiff (quelques mL) \\
\hline
Thermomètre (facultatif) & Liqueur de Fehling fraîchement préparée \\
\hline
Pissette d'eau distillée & Eau distillée \\
\hline
\end{tabularx}
\end{center}

\begin{popfigure}
\begin{center}
\begin{tikzpicture}[scale=0.85, every node/.style={font=\small}]
% Molécules avec chemfig
\node (ethanal) at (0,0) {\chemfig{CH_3-CH=O}};
\node (propanone) at (5,0) {\chemfig{CH_3-C(=[2]O)-CH_3}};
\node (ethanol) at (10,0) {\chemfig{CH_3-CH_2-OH}};
% Noms
\node[below=3pt, popdark, font=\bfseries\small] at (ethanal.south) {Éthanal (aldéhyde)};
\node[below=3pt, popdark, font=\bfseries\small] at (propanone.south) {Propanone (cétone)};
\node[below=3pt, popdark, font=\bfseries\small] at (ethanol.south) {Éthanol (alcool primaire)};
% Groupes fonctionnels mis en évidence
\draw[poporange, thick, decorate, decoration={brace, amplitude=5pt}] (ethanal.north east) -- node[above=5pt, poporange, font=\small] {Groupe \cle{carbonyle} \ce{C=O}} (ethanal.north west);
\draw[poporange, thick, decorate, decoration={brace, amplitude=5pt}] (propanone.north east) -- node[above=5pt, poporange, font=\small] {Groupe \cle{carbonyle} \ce{C=O}} (propanone.north west);
\draw[popgreen, thick, decorate, decoration={brace, amplitude=5pt}] (ethanol.north east) -- node[above=5pt, popgreen, font=\small] {Groupe \cle{hydroxyle} \ce{-OH}} (ethanol.north west);
\end{tikzpicture}
\end{center}
\legende{Formules semi-développées des trois échantillons : l'éthanal et la propanone portent un groupe carbonyle (test DNPH +), l'éthanol porte un groupe hydroxyle (tests carbonylés -).}
\end{popfigure}

\begin{savaistu}[Et si un réactif manque au laboratoire ?]
Dans de nombreux lycées camerounais, les réactifs sont coûteux. Quelques alternatives :
\begin{itemize}
\item La \textbf{liqueur de Fehling} peut être préparée au laboratoire en mélangeant, au moment de l'emploi, une solution de sulfate de cuivre(II) (\ce{CuSO4}) et une solution de potasse (\ce{KOH}) additionnée de sel de Seignette (tartrate double de sodium et de potassium, disponible en pharmacie) : le mélange doit être bleu et limpide.
\item L'\textbf{éthanal} peut être remplacé par du formol (solution commerciale de méthanal) dilué, ou par du glucose en poudre (le glucose possède une fonction aldéhyde et donne un test de Fehling positif).
\item L'\textbf{acétone} vendue en quincaillerie ou en pharmacie comme dissolvant convient parfaitement.
\item À défaut de bain-marie, une casserole contenant de l'eau chauffée sur un réchaud à gaz ou à bois fait l'affaire.
\end{itemize}
\end{savaistu}

\courssub{Sécurité}

\begin{attention}[Consignes de sécurité — à lire avant toute manipulation]
\begin{itemize}
\item \textbf{Pictogramme « inflammable »} (flamme) : l'éthanol, l'acétone et l'éthanal sont très inflammables. Aucune flamme nue à proximité des flacons ; le chauffage se fait uniquement au \cle{bain-marie}, jamais le tube directement dans la flamme.
\item \textbf{Pictogramme « toxique / nocif »} (tête de mort ou point d'exclamation) : la 2,4-DNPH est toxique et tache durablement la peau en jaune ; le formol est nocif par inhalation. Porte des gants si disponibles, des lunettes de protection et une blouse ; travaille près d'une fenêtre ouverte.
\item \textbf{Pictogramme « corrosif »} : la liqueur de Fehling contient une base concentrée (\ce{KOH}). En cas de contact avec la peau, rincer abondamment à l'eau froide pendant plusieurs minutes et prévenir le professeur.
\item Ne jamais pipeter avec la bouche ; utiliser exclusivement les pipettes compte-gouttes.
\item Ne jamais goûter ni sentir directement les produits : pour percevoir une odeur, agiter doucement l'air au-dessus du tube vers le nez avec la main.
\item Les déchets contenant du cuivre et de la DNPH ne se jettent pas à l'évier : les verser dans le bidon de récupération prévu par le professeur.
\end{itemize}
\end{attention}

\courssub{Schéma du dispositif}

\begin{popfigure}
\begin{center}
\begin{tikzpicture}[scale=0.95, every node/.style={font=\small}]
% Bécher bain-marie
\draw[thick, popblue] (-1.2,0) -- (-1.2,2.6) -- (-0.9,2.6);
\draw[thick, popblue] (1.2,0) -- (1.2,2.6) -- (0.9,2.6);
\draw[thick, popblue] (-1.2,0) -- (1.2,0);
% Eau
\fill[popblue!25] (-1.15,0.05) rectangle (1.15,1.5);
\draw[popblue!60, decorate, decoration={snake, amplitude=1.2pt, segment length=8pt}] (-1.1,1.5) -- (1.1,1.5);
\node[popblue!80!black] at (0,0.9) {eau chaude};
% Tube à essai
\draw[thick, popdark] (-0.25,3.4) -- (-0.25,0.6) arc[start angle=180, end angle=360, radius=0.25] -- (0.25,3.4);
\fill[poporange!40] (-0.22,0.65) arc[start angle=180, end angle=360, radius=0.22] -- (0.22,1.6) -- (-0.22,1.6) -- cycle;
\node[popdark, align=center] at (0,1.15) {\scriptsize Fehling\\ \scriptsize + échantillon};
% Flamme
\draw[poporange, fill=poporange!60] (0,-0.9) .. controls (0.35,-0.55) and (0.2,-0.25) .. (0,-0.05) .. controls (-0.2,-0.25) and (-0.35,-0.55) .. (0,-0.9);
\draw[popgold, fill=popgold!70] (0,-0.75) .. controls (0.15,-0.5) .. (0,-0.2) .. controls (-0.15,-0.5) .. (0,-0.75);
% Support
\draw[thick, popmuted] (-1.6,-1.0) -- (1.6,-1.0);
% Annotations
\draw[->, popdark] (2.6,1.2) -- (1.25,1.2);
\node[popdark, anchor=west] at (2.6,1.2) {bain-marie à \SI{60}{\celsius}--\SI{70}{\celsius}};
\draw[->, popdark] (2.6,2.6) -- (0.3,2.9);
\node[popdark, anchor=west] at (2.6,2.6) {tube à essai};
\draw[->, popdark] (2.6,-0.5) -- (0.4,-0.55);
\node[popdark, anchor=west] at (2.6,-0.5) {flamme du bec Bunsen};
\end{tikzpicture}
\end{center}
\legende{Dispositif de chauffage au bain-marie pour le test à la liqueur de Fehling : le tube contenant le mélange réactionnel est plongé dans l'eau chaude, jamais chauffé directement à la flamme.}
\end{popfigure}

\courssub{Protocole}

\begin{methode}[Déroulement des trois tests]
\textbf{Préparation.} Numéroter neuf tubes à essai : A1, A2, A3 pour l'échantillon A ; B1, B2, B3 pour B ; C1, C2, C3 pour C. Introduire dans chaque tube environ \SI{1}{mL} (soit 20 gouttes) de l'échantillon correspondant, avec une pipette propre dédiée à chaque flacon.

\begin{enumerate}
\item \textbf{Test à la 2,4-DNPH} (tubes A1, B1, C1) : ajouter 5 gouttes de solution de 2,4-DNPH dans chaque tube. Agiter doucement. Observer l'apparition éventuelle d'un précipité jaune-orangé. Noter l'observation immédiatement.
\item \textbf{Test au réactif de Schiff} (tubes A2, B2, C2) : ajouter environ \SI{1}{mL} de réactif de Schiff dans chaque tube. Agiter et attendre 2 à 3 minutes \emph{à froid}. Observer l'apparition éventuelle d'une coloration rose-violette.
\item \textbf{Test à la liqueur de Fehling} (tubes A3, B3, C3) : ajouter environ \SI{2}{mL} de liqueur de Fehling (solution bleue) dans chaque tube. Plonger les trois tubes dans le bain-marie maintenu entre \SI{60}{\celsius} et \SI{70}{\celsius} pendant 3 à 5 minutes. Observer l'apparition éventuelle d'un précipité rouge brique et le changement de teinte de la solution.
\item \textbf{Consignation} : reporter toutes les observations dans le tableau de résultats, en précisant la couleur, la nature (précipité ou coloration) et la rapidité d'apparition.
\item \textbf{Conclusion} : déduire de l'ensemble des tests la nature (aldéhyde, cétone ou alcool) de chacun des échantillons A, B et C.
\end{enumerate}
\end{methode}

\courssub{Observations et résultats attendus}

Voici le tableau de résultats type, tel qu'il doit apparaître dans ton compte rendu :

\begin{center}
\begin{tabularx}{\linewidth}{|l|X|X|X|}
\hline
\rowcolor{popblueL} \textbf{Test} & \textbf{Échantillon A} & \textbf{Échantillon B} & \textbf{Échantillon C} \\
\hline
2,4-DNPH & Précipité jaune-orangé abondant, rapide & Précipité jaune-orangé & Aucun changement \\
\hline
Réactif de Schiff & Coloration rose-violette en 2 min & Aucune coloration & Aucune coloration \\
\hline
Liqueur de Fehling (à chaud) & Précipité rouge brique, solution verdit & Solution bleue inchangée & Solution bleue inchangée \\
\hline
\rowcolor{popgreenL} \textbf{Conclusion} & Aldéhyde (éthanal) & Cétone (propanone) & Alcool (éthanol) \\
\hline
\end{tabularx}
\end{center}

\begin{attention}[Pièges d'observation]
\begin{itemize}
\item Un léger trouble avec la DNPH sur l'échantillon C peut provenir d'une pipette mal rincée : toujours utiliser une pipette par flacon.
\item Le réactif de Schiff rosit lentement à l'air : ne conclure qu'après comparaison avec un tube témoin contenant uniquement le réactif.
\item Un chauffage trop prolongé du test de Fehling peut faire brunir la solution même sans aldéhyde : respecter les 3 à 5 minutes.
\end{itemize}
\end{attention}

\courssub{Exploitation}

\begin{exoresolu}[Questions d'exploitation]
\textbf{1.} Pourquoi le test à la 2,4-DNPH, à lui seul, ne permet-il pas de distinguer A de B ?

\textbf{2.} Quels tests permettent de trancher entre A et B ? Justifier.

\textbf{3.} L'échantillon C ne réagit avec aucun des trois réactifs. Que peut-on en conclure ?

\textbf{4.} Écrire l'équation-bilan de la réaction entre l'éthanal (A) et la liqueur de Fehling, sachant que l'ion \ce{Cu^{2+}} (complexé, bleu) est réduit en oxyde de cuivre(I) \ce{Cu2O} (rouge brique) et que l'aldéhyde est oxydé en ion carboxylate \ce{CH3COO-} en milieu basique.

\textbf{5.} On réalise les trois tests sur un quatrième flacon D : DNPH positive, Schiff négatif, Fehling négatif. Quelle est la famille de D ? Proposer un exemple de composé.
\tcblower
\textbf{1.} La 2,4-DNPH réagit avec le groupe carbonyle \ce{C=O} présent dans les aldéhydes \emph{et} dans les cétones. A et B possèdent tous deux ce groupe : le précipité jaune-orangé apparaît donc dans les deux cas. Ce test prouve seulement que A et B sont des \cle{composés carbonylés}.

\textbf{2.} Ce sont le réactif de Schiff et la liqueur de Fehling : ce sont des oxydants ménagés qui n'oxydent que les aldéhydes. A donne une coloration rose-violette avec Schiff et un précipité rouge brique avec Fehling : A est un \textbf{aldéhyde}. B ne réagit ni à Schiff ni à Fehling : B est une \textbf{cétone}.

\textbf{3.} C n'est ni un aldéhyde ni une cétone : il ne possède pas de groupe carbonyle. Sa réponse négative aux trois tests est compatible avec un \textbf{alcool} (ici l'éthanol), ce qui est confirmé par son odeur caractéristique et son inflammabilité.

\textbf{4.} L'équation-bilan s'écrit :
\[ \ce{CH3CHO + 2 Cu^{2+} + 5 OH- -> CH3COO- + Cu2O + 3 H2O} \]
L'éthanal est oxydé en ion éthanoate ; les ions cuivre(II), responsables de la couleur bleue, sont réduits en oxyde de cuivre(I) rouge brique qui précipite.

\textbf{5.} D est une \textbf{cétone} : elle possède un groupe carbonyle (DNPH positive) mais n'est pas oxydable par les oxydants ménagés. Exemple : la propanone \ce{CH3COCH3} ou la butanone \ce{CH3COCH2CH3}.
\end{exoresolu}

\courssub{Conclusion}

Ce TP illustre la stratégie générale d'\cle{identification} des composés organiques : on combine plusieurs tests complémentaires, car aucun test isolé ne suffit. La 2,4-DNPH détecte la famille des composés carbonylés ; le réactif de Schiff et la liqueur de Fehling distinguent ensuite, au sein de cette famille, les aldéhydes (oxydables) des cétones (non oxydables par les oxydants ménagés). L'absence de réponse aux trois tests oriente vers un alcool. Cette démarche par élimination, rigoureuse et économique en réactifs, est celle utilisée dans les laboratoires de contrôle : elle montre que la chimie est une science expérimentale où l'observation soigneuse, consignée dans un tableau, conduit à une conclusion argumentée. Retiens la phrase-clé : \emph{« DNPH pour le carbonyle, Schiff et Fehling pour l'aldéhyde »}.

\newpage

\coursec{Méthodes et exercices résolus}

\courssub{Fiche méthode 1 : Nommer un alcool}

\begin{methode}[Nommer un alcool en nomenclature systématique]
\begin{enumerate}
\item \textbf{Repérer} le groupe caractéristique hydroxyle \ce{-OH} dans la formule développée ou semi-développée.
\item \textbf{Identifier la chaîne principale} : c'est la chaîne carbonée la plus longue contenant le carbone fonctionnel (celui qui porte le \ce{-OH}).
\item \textbf{Numéroter} la chaîne de sorte que le carbone fonctionnel ait l'indice de position \emph{le plus petit possible}.
\item \textbf{Construire le nom} : nom de l'alcane correspondant, dont on remplace le \emph{e} final par le suffixe \cle{-ol}, précédé de l'indice de position : \emph{alcane} $\to$ \emph{alcan-}indice\emph{-ol}.
\item \textbf{Citer les ramifications} (alkyles) en préfixe, avec leurs indices, par ordre alphabétique.
\item \textbf{Vérifier} : le nom ne commence jamais par « 1-ol » si la chaîne n'a qu'un seul carbone fonctionnel en bout de chaîne... sauf si c'est nécessaire pour lever une ambiguïté (ex. propan-1-ol / propan-2-ol).
\end{enumerate}
\textbf{Réflexe} : si le \ce{-OH} est sur un carbone en bout de chaîne, l'indice est 1 ; s'il est au milieu, compare les deux sens de numérotation.
\end{methode}

\begin{exoresolu}[Nommer des alcools]
Nommer les alcools suivants :
\[
\text{a) } \ce{CH3-CH2-CH2-CH2-OH} \qquad
\text{b) } \chemfig{CH_3-CH(-[2]OH)-CH_2-CH_3} \qquad
\text{c) } \chemfig{CH_3-CH(-[6]CH_3)-CH_2-CH_2-OH}
\]
\tcblower
\textbf{a)} La chaîne la plus longue contenant le \ce{-OH} compte 4 carbones : butane. Le \ce{-OH} est sur le carbone 1 (bout de chaîne). Pas de ramification.

\textbf{Nom : butan-1-ol.}

\textbf{b)} Chaîne principale : 4 carbones (butane). Numérotation : depuis la gauche, le \ce{-OH} est en position 2 ; depuis la droite, en position 3. On retient l'indice le plus petit : 2.

\textbf{Nom : butan-2-ol.}

\textbf{c)} La chaîne la plus longue contenant le \ce{-OH} compte 4 carbones (attention : le \ce{CH3} en ramification ne fait pas partie de la chaîne principale car on doit inclure le carbone fonctionnel). Numérotation à partir du carbone portant le \ce{-OH} : \ce{-OH} en position 1, méthyle en position 3.

\textbf{Nom : 3-méthylbutan-1-ol.}

\textbf{Conclusion} : a) butan-1-ol ; b) butan-2-ol ; c) 3-méthylbutan-1-ol.
\end{exoresolu}

\courssub{Fiche méthode 2 : Distinguer les trois classes d'alcools}

\begin{methode}[Déterminer la classe d'un alcool]
\begin{enumerate}
\item \textbf{Repérer le carbone fonctionnel} : c'est le carbone qui porte le groupe \ce{-OH}.
\item \textbf{Compter les atomes de carbone directement liés} à ce carbone fonctionnel (les groupes alkyles $R$).
\item \textbf{Conclure} selon le tableau :
\begin{center}
\begin{tabularx}{\linewidth}{|l|X|X|}
\hline
\rowcolor{popblueL} \textbf{Nombre de C liés} & \textbf{Classe} & \textbf{Structure} \\
\hline
1 (ou 0 pour le méthanol) & Alcool \cle{primaire} (I) & $R-CH_2-OH$ \\
\hline
2 & Alcool \cle{secondaire} (II) & $R-CH(OH)-R'$ \\
\hline
3 & Alcool \cle{tertiaire} (III) & $R-C(OH)(R')-R''$ \\
\hline
\end{tabularx}
\end{center}
\item \textbf{Piège à éviter} : ne compte que les \emph{carbones} voisins, jamais les hydrogènes.
\end{enumerate}
\textbf{Réflexe} : un alcool primaire a toujours le motif $-CH_2-OH$ ; un alcool secondaire le motif $-CH(OH)-$ ; un alcool tertiaire le motif $-C(OH)R_3$.
\end{methode}

\begin{exoresolu}[Classes des isomères de formule brute $\ce{C4H10O}$]
On donne les quatre alcools isomères de formule brute \ce{C4H10O} :
\begin{itemize}
\item A : \ce{CH3-CH2-CH2-CH2-OH}
\item B : \chemfig{CH_3-CH(-[2]OH)-CH_2-CH_3}
\item C : \chemfig{CH_3-CH(-[6]CH_3)-CH_2-OH}
\item D : \chemfig{CH_3-C(-[2]OH)(-[6]CH_3)-CH_3}
\end{itemize}
Donner le nom et la classe de chacun.
\tcblower
On repère dans chaque cas le carbone fonctionnel et on compte ses voisins carbonés.

\textbf{A} : carbone fonctionnel \ce{CH2} lié à \textbf{un seul} carbone. Nom : \textbf{butan-1-ol}, alcool \textbf{primaire}.

\textbf{B} : carbone fonctionnel \ce{CH} lié à \textbf{deux} carbones (\ce{CH3} et \ce{CH2CH3}). Nom : \textbf{butan-2-ol}, alcool \textbf{secondaire}.

\textbf{C} : carbone fonctionnel \ce{CH2} lié à \textbf{un seul} carbone. Nom : \textbf{2-méthylpropan-1-ol}, alcool \textbf{primaire}.

\textbf{D} : carbone fonctionnel \ce{C} lié à \textbf{trois} carbones (trois \ce{CH3}). Nom : \textbf{2-méthylpropan-2-ol}, alcool \textbf{tertiaire}.

\textbf{Conclusion} : A et C sont primaires, B est secondaire, D est tertiaire. Ces quatre composés sont des isomères de constitution (même formule brute \ce{C4H10O}, structures différentes).
\end{exoresolu}

\courssub{Fiche méthode 3 : Expliquer les propriétés physiques des alcools}

\begin{methode}[Relier structure moléculaire et propriétés physiques]
\begin{enumerate}
\item \textbf{Repérer le groupe hydroxyle} \ce{-OH} : il est polaire et permet la formation de \cle{liaisons hydrogène} entre molécules d'alcool.
\item \textbf{Températures de fusion et d'ébullition} : les liaisons hydrogène renforcent la cohésion entre molécules $\Rightarrow$ les alcools ont des températures d'ébullition \emph{nettement plus élevées} que les alcanes de masse molaire voisine. Plus la chaîne carbonée est longue, plus la température d'ébullition augmente.
\item \textbf{Solubilité dans l'eau} : le groupe \ce{-OH} (partie \emph{hydrophile}) forme des liaisons hydrogène avec l'eau ; la chaîne carbonée (partie \emph{hydrophobe}) s'y oppose. Les alcools à chaîne courte ($n \leq 4$) sont miscibles ou très solubles ; la solubilité \emph{diminue} quand la chaîne s'allonge.
\item \textbf{Polyalcools} : plus il y a de groupes \ce{-OH} (glycol : 2 ; glycérol : 3), plus les liaisons hydrogène sont nombreuses $\Rightarrow$ températures d'ébullition et solubilité dans l'eau encore plus élevées, liquides visqueux.
\end{enumerate}
\textbf{Réflexe} : toujours citer la \emph{liaison hydrogène} dans la justification, c'est le mot-clé attendu par le correcteur.
\end{methode}

\begin{exoresolu}[Comparaison de températures d'ébullition et de solubilités]
On donne les températures d'ébullition : propane $\theta_{eb} = \SI{-42}{\celsius}$ ; éthanol $\theta_{eb} = \SI{78}{\celsius}$ ; glycérol (propane-1,2,3-triol) $\theta_{eb} = \SI{290}{\celsius}$.
\begin{enumerate}
\item Expliquer pourquoi l'éthanol bout à une température bien plus élevée que le propane, alors que leurs masses molaires sont voisines ($M(\ce{C3H8}) = \SI{44}{g.mol^{-1}}$ ; $M(\ce{C2H5OH}) = \SI{46}{g.mol^{-1}}$).
\item Expliquer pourquoi le glycérol est miscible à l'eau en toutes proportions alors que l'hexan-1-ol y est très peu soluble.
\end{enumerate}
\tcblower
\textbf{1.} Les molécules de propane sont apolaires : seules de faibles interactions de Van der Waals les unissent, faciles à rompre. L'éthanol possède un groupe hydroxyle \ce{-OH} polaire : ses molécules s'associent par \textbf{liaisons hydrogène} (\ce{O-H ... O}). Rompre ces liaisons exige beaucoup plus d'énergie, d'où une température d'ébullition très supérieure ($\SI{78}{\celsius}$ contre $\SI{-42}{\celsius}$) malgré des masses molaires quasi identiques.

\textbf{2.} Le glycérol \ce{HO-CH2-CH(OH)-CH2-OH} possède \textbf{trois groupes hydroxyle} pour seulement trois carbones : sa partie hydrophile domine largement. Il forme de nombreuses liaisons hydrogène avec l'eau, d'où une miscibilité totale. L'hexan-1-ol ne possède qu'\textbf{un seul} groupe \ce{-OH} pour une longue chaîne de six carbones, hydrophobe : la partie carbonée l'emporte, la solubilité dans l'eau devient très faible.

\textbf{Conclusion} : les liaisons hydrogène expliquent à la fois les températures d'ébullition élevées et la solubilité dans l'eau ; le rapport (nombre de \ce{-OH}) / (longueur de la chaîne carbonée) est le critère décisif.
\end{exoresolu}

\courssub{Fiche méthode 4 : Nommer un aldéhyde et une cétone}

\begin{methode}[Nomenclature des aldéhydes et des cétones]
\textbf{Pour un aldéhyde} $R-\ce{CHO}$ :
\begin{enumerate}
\item Le groupe carbonyle \ce{C=O} est \emph{toujours en extrémité} de chaîne : le carbone fonctionnel est le carbone \textbf{n°1}.
\item Chaîne principale la plus longue contenant le \ce{-CHO} ; nom de l'alcane + suffixe \cle{-al} : \emph{alcanal} (pas d'indice, il est sous-entendu 1).
\item Citer les ramifications en préfixe.
\end{enumerate}
\textbf{Pour une cétone} $R-\ce{CO}-R'$ :
\begin{enumerate}
\item Le groupe carbonyle est \emph{à l'intérieur} de la chaîne.
\item Chaîne principale la plus longue contenant le \ce{C=O} ; numéroter pour donner au carbonyle l'indice \emph{le plus petit}.
\item Nom de l'alcane + suffixe \cle{-one} précédé de l'indice : \emph{alcan-}indice\emph{-one}.
\end{enumerate}
\textbf{Réflexe} : aldéhyde = \ce{-CHO} en bout (suffixe \emph{-al}) ; cétone = \ce{-CO-} au milieu (suffixe \emph{-one} + indice obligatoire).
\end{methode}

\begin{exoresolu}[Nommer des aldéhydes et des cétones]
Nommer les composés suivants et préciser leur famille :
\[
\text{a) } \ce{CH3-CH2-CHO} \qquad
\text{b) } \ce{CH3-CO-CH2-CH3} \qquad
\text{c) } \chemfig{CH_3-CH(-[6]CH_3)-CHO} \qquad
\text{d) } \ce{CH3-CH2-CO-CH2-CH3}
\]
\tcblower
\textbf{a)} Groupe \ce{-CHO} en extrémité : \textbf{aldéhyde}. Chaîne de 3 carbones $\Rightarrow$ \textbf{propanal}.

\textbf{b)} Groupe \ce{C=O} à l'intérieur de la chaîne : \textbf{cétone}. Chaîne de 4 carbones ; numérotation donnant l'indice le plus petit au carbonyle : position 2 $\Rightarrow$ \textbf{butan-2-one} (ou butanone).

\textbf{c)} Groupe \ce{-CHO} en extrémité : \textbf{aldéhyde}. Chaîne principale de 3 carbones contenant le \ce{-CHO}, méthyle en position 2 $\Rightarrow$ \textbf{2-méthylpropanal}.

\textbf{d)} \textbf{Cétone}. Chaîne de 5 carbones, carbonyle en position 3 (identique dans les deux sens) $\Rightarrow$ \textbf{pentan-3-one}.

\textbf{Conclusion} : a) propanal (aldéhyde) ; b) butan-2-one (cétone) ; c) 2-méthylpropanal (aldéhyde) ; d) pentan-3-one (cétone).
\end{exoresolu}

\courssub{Fiche méthode 5 : Distinguer aldéhydes et cétones par les tests}

\begin{methode}[Choisir et interpréter les tests d'identification]
\begin{enumerate}
\item \textbf{Test commun — 2,4-DNPH} : aldéhydes \emph{et} cétones donnent un \textbf{précipité jaune-orangé}. Ce test prouve la présence du groupe carbonyle \ce{C=O} mais ne distingue pas les deux familles.
\item \textbf{Tests discriminants} (oxydation ménagée, positifs uniquement pour les aldéhydes) :
\begin{center}
\begin{tabularx}{\linewidth}{|l|X|X|}
\hline
\rowcolor{popblueL} \textbf{Réactif} & \textbf{Aldéhyde} & \textbf{Cétone} \\
\hline
2,4-DNPH & Précipité jaune-orangé & Précipité jaune-orangé \\
\hline
Réactif de Schiff & Coloration \textbf{rose-violette} & Incolore (négatif) \\
\hline
Liqueur de Fehling (à chaud) & Précipité \textbf{rouge brique} \ce{Cu2O} & Reste bleue (négatif) \\
\hline
Réactif de Tollens & \textbf{Miroir d'argent} & Négatif \\
\hline
\end{tabularx}
\end{center}
\item \textbf{Stratégie} : 2,4-DNPH d'abord (carbonyle ?), puis Schiff ou Fehling (aldéhyde ou cétone ?).
\item \textbf{Justification} : l'aldéhyde est un \emph{réducteur} (il s'oxyde en acide carboxylique) ; la cétone ne subit pas l'oxydation ménagée.
\end{enumerate}
\textbf{Réflexe} : « test positif » = changement visible observé ; toujours préciser la couleur attendue.
\end{methode}

\begin{exoresolu}[Identification de deux composés carbonylés]
Deux flacons sans étiquette A et B contiennent l'un le propanal, l'autre la propanone. On réalise les tests suivants :
\begin{itemize}
\item Avec la 2,4-DNPH, A et B donnent tous deux un précipité jaune-orangé.
\item Avec la liqueur de Fehling chauffée, A donne un précipité rouge brique ; B ne réagit pas.
\item Avec le réactif de Schiff, A donne une coloration rose-violette ; B reste incolore.
\end{itemize}
Identifier A et B en justifiant chaque étape.
\tcblower
\textbf{Étape 1 — 2,4-DNPH} : les deux composés donnent un précipité jaune-orangé $\Rightarrow$ A et B possèdent tous deux un groupe \textbf{carbonyle} \ce{C=O}. Ce test ne permet pas encore de les distinguer : c'est cohérent, le propanal et la propanone sont tous deux des composés carbonylés.

\textbf{Étape 2 — Liqueur de Fehling} : le précipité rouge brique de \ce{Cu2O} obtenu avec A montre que A est un \textbf{réducteur} : c'est un \textbf{aldéhyde}. B ne réduit pas la liqueur de Fehling : c'est une \textbf{cétone}.

\textbf{Étape 3 — Réactif de Schiff} : la coloration rose-violette avec A confirme que A est un aldéhyde ; le test négatif avec B confirme la cétone.

\textbf{Conclusion} : \textbf{A est le propanal} \ce{CH3-CH2-CHO} et \textbf{B est la propanone} \ce{CH3-CO-CH3}. Les tests de Fehling et de Schiff, concordants, lèvent toute ambiguïté laissée par la 2,4-DNPH.
\end{exoresolu}

\courssub{Fiche méthode 6 : Écrire les équations-bilans des réactions de test}

\begin{methode}[Rédiger l'équation d'un test d'oxydation ménagée]
\begin{enumerate}
\item \textbf{Écrire le couple oxydant/réducteur} mis en jeu : l'aldéhyde $R-\ce{CHO}$ s'oxyde en acide carboxylique $R-\ce{COOH}$ (ou ion carboxylate $R-\ce{COO-}$ en milieu basique).
\item \textbf{Identifier l'oxydant du test} : ion \ce{Cu^2+} complexé (Fehling, réduit en \ce{Cu2O} rouge brique) ou ion \ce{Ag+} complexé (Tollens, réduit en Ag métallique).
\item \textbf{Écrire les demi-équations} en précisant le milieu (Fehling et Tollens s'utilisent en milieu basique) :
\begin{itemize}
\item Oxydation : $\ce{R-CHO + 3OH- -> R-COO- + 2H2O + 2e-}$
\item Fehling : $\ce{2Cu^2+ + 2OH- + 2e- -> Cu2O + H2O}$
\item Tollens : $\ce{Ag+ + e- -> Ag}$
\end{itemize}
\item \textbf{Combiner} en égalisant les électrons échangés, puis simplifier les espèces communes.
\item \textbf{Vérifier} la conservation des atomes et des charges.
\end{enumerate}
\textbf{Réflexe} : le bilan global d'un test de Fehling fait apparaître le précipité \ce{Cu2O} ; celui de Tollens fait apparaître Ag solide (le miroir).
\end{methode}

\begin{exoresolu}[Équation du test de la liqueur de Fehling sur l'éthanal]
L'éthanal \ce{CH3-CHO} donne un test positif à la liqueur de Fehling chauffée : on observe un précipité rouge brique d'oxyde de cuivre(I) \ce{Cu2O}. En milieu basique, l'éthanal est oxydé en ion éthanoate \ce{CH3-COO-}.
\begin{enumerate}
\item Écrire les deux demi-équations électroniques.
\item En déduire l'équation-bilan de la réaction.
\end{enumerate}
On donne les couples : \ce{CH3-COO-}/\ce{CH3-CHO} et \ce{Cu^2+}/\ce{Cu2O}.
\tcblower
\textbf{1. Demi-équations} (milieu basique, on équilibre avec \ce{OH-} et \ce{H2O}) :

Oxydation de l'aldéhyde :
\[
\ce{CH3-CHO + 3OH- -> CH3-COO- + 2H2O + 2e-}
\]

Réduction des ions cuivre(II) :
\[
\ce{2Cu^2+ + 2OH- + 2e- -> Cu2O + H2O}
\]

\textbf{2. Équation-bilan} : les deux demi-équations échangent chacune 2 électrons ; on les additionne membre à membre :
\[
\ce{CH3-CHO + 3OH- + 2Cu^2+ + 2OH- -> CH3-COO- + 2H2O + Cu2O + H2O}
\]
En simplifiant (\ce{5OH-} à gauche, \ce{3H2O} à droite) :
\[
\boxed{\ce{CH3-CHO + 2Cu^2+ + 5OH- -> CH3-COO- + Cu2O + 3H2O}}
\]

\textbf{Vérification} : charges à gauche $2\times(+2) + 5\times(-1) = -1$ ; à droite $-1$. Atomes : C : 2/2 ; H : 9/9 ; O : 6/6 ; Cu : 2/2. L'équation est équilibrée.

\textbf{Conclusion} : le précipité rouge brique de \ce{Cu2O} observé est la signature de la réduction des ions \ce{Cu^2+} par l'éthanal, oxydé en ion éthanoate.
\end{exoresolu}

\begin{attention}[Les 5 erreurs qui coûtent des points au Probatoire]
\begin{enumerate}
\item \textbf{Oublier l'indice de position} dans le nom d'une cétone (écrire « butanone » est toléré, mais « pentanone » sans indice est fautif : pentan-2-one ou pentan-3-one ?). Pour un alcool, l'indice est obligatoire dès qu'il y a ambiguïté : propan-1-ol $\neq$ propan-2-ol.
\item \textbf{Mal compter les voisins du carbone fonctionnel} : la classe d'un alcool se détermine en comptant les \emph{carbones} liés au carbone porteur du \ce{-OH}, pas les hydrogènes ni la longueur totale de la chaîne.
\item \textbf{Croire que la 2,4-DNPH distingue aldéhyde et cétone} : elle donne un précipité jaune-orangé avec \emph{les deux}. Seuls Schiff, Fehling et Tollens sont discriminants (positifs pour l'aldéhyde seul).
\item \textbf{Écrire des couleurs vagues} : « le test est positif » ne suffit pas. Il faut préciser : précipité \emph{rouge brique} (Fehling), coloration \emph{rose-violette} (Schiff), \emph{miroir d'argent} (Tollens), précipité \emph{jaune-orangé} (DNPH).
\item \textbf{Déséquilibrer les équations de test} : oublier les \ce{OH-} du milieu basique, ou ne pas vérifier la conservation des charges. Pense aussi à simplifier les \ce{H2O} et \ce{OH-} présents des deux côtés avant de conclure.
\end{enumerate}
\end{attention}

\newpage

\coursec{Exercices}

\coursec{Leçon 03 : Les composés oxygénés : alcools, aldéhydes, cétones}

\courssub{1. Généralités sur les composés oxygénés}

Les \cle{composés oxygénés} sont des composés organiques dont la molécule contient au moins un atome d'oxygène lié à la chaîne carbonée. Ils sont omniprésents dans la vie quotidienne au Cameroun : l'\textbf{éthanol} dans le vin de palme, le bili-bili ou les boissons brassées ; l'\textbf{éthanal} dans les arômes de fruits mûrs ; l'\textbf{acide acétique} dans le vinaigre de palme ; le \textbf{glycérol} dans les savons artisanaux.

\begin{definition}[Composé oxygéné]
Un composé oxygéné est un dérivé d'hydrocarbure obtenu par substitution d'un ou plusieurs atomes d'hydrogène par un ou plusieurs groupes fonctionnels contenant de l'oxygène.
\end{definition}

Les familles étudiées dans cette leçon se distinguent par leur \cle{groupe fonctionnel} (GF) :
\begin{itemize}
    \item \textbf{Alcools} : groupe \cle{hydroxyle} $-\ce{OH}$ lié à un carbone \emph{sp³} (saturé).
    \item \textbf{Aldéhydes} : groupe \cle{carbonyle} $\ce{C=O}$ porté par un carbone terminal (lié à au moins un H).
    \item \textbf{Cétones} : groupe \cle{carbonyle} $\ce{C=O}$ porté par un carbone interne (lié à deux atomes de carbone).
\end{itemize}

\begin{aretenir}[Groupes fonctionnels et formules générales]
\begin{center}
\begin{tabularx}{\linewidth}{|l|X|c|X|}
\hline
\rowcolor{popblueL} \textbf{Famille} & \textbf{Groupe fonctionnel} & \textbf{Formule générale (saturé, acyclique)} & \textbf{Exemple courant} \\
\hline
Alcool & $-\ce{OH}$ (hydroxyle) & $C_nH_{2n+2}O$ \quad ($n \ge 1$) & Éthanol $\ce{CH3CH2OH}$ (vin de palme) \\
\hline
Aldéhyde & $\ce{-CHO}$ (formyle) & $C_nH_{2n}O$ \quad ($n \ge 1$) & Éthanal $\ce{CH3CHO}$ (arôme banane) \\
\hline
Cétone & $\ce{-CO-}$ (oxo) & $C_nH_{2n}O$ \quad ($n \ge 3$) & Propanone $\ce{CH3COCH3}$ (solvant) \\
\hline
\end{tabularx}
\end{center}
\end{aretenir}

\courssub{2. Les alcools : formules, nomenclature et classes}

\paragraph{Formules.} La formule brute générale des alcools saturés acycliques est \textbf{$C_nH_{2n+2}O$} ($n \in \mathbb{N}^*$). La formule semi-développée montre le groupe $-\ce{OH}$ sur la chaîne.

\paragraph{Nomenclature (IUPAC).}
\begin{methode}[Nommer un alcool]
\begin{enumerate}
    \item Identifier la \textbf{chaîne principale} : la plus longue chaîne carbonée portant le groupe $-\ce{OH}$.
    \item Numéroter la chaîne en donnant le \textbf{plus petit locant} au carbone portant le $-\ce{OH}$ (le carbone fonctionnel).
    \item Nommer les substituants (alkyles, halogènes...) avec leurs locants.
    \item Écrire le nom : \emph{locant-ol} (ex. \emph{butan-2-ol}) ou \emph{locant-ol-...} si substitué. Le suffixe est \textbf{-ol}.
\end{enumerate}
\end{methode}

\paragraph{Classes d'alcools (Isomérie de fonction et de chaîne).} La classe dépend du nombre d'atomes de carbone \textbf{directement liés au carbone fonctionnel} (celui qui porte le $-\ce{OH}$).

\begin{propriete}[Classification des alcools]
\begin{itemize}
    \item \textbf{Alcool primaire (1°)} : Le carbone fonctionnel est lié à \textbf{un seul} atome de carbone (et deux H). \quad $R-CH_2-OH$
    \item \textbf{Alcool secondaire (2°)} : Le carbone fonctionnel est lié à \textbf{deux} atomes de carbone (et un H). \quad $R-CH(OH)-R'$
    \item \textbf{Alcool tertiaire (3°)} : Le carbone fonctionnel est lié à \textbf{trois} atomes de carbone (pas de H). \quad $R-C(OH)(R')-R''$
\end{itemize}
\end{propriete}

\begin{exemplebox}[Exemples des trois classes (formule brute \ce{C4H10O})]
\begin{center}
\begin{tabular}{ccc}
\chemfig{CH_3-CH_2-CH_2-CH_2-OH} & \chemfig{CH_3-CH_2-CH(-[2]OH)-CH_3} & \chemfig{CH_3-C(-[2]CH_3)(-[6]OH)-CH_3} \\
Butan-1-ol \textbf{(1°)} & Butan-2-ol \textbf{(2°)} & 2-méthylpropan-2-ol \textbf{(3°)}
\end{tabular}
\end{center}
\end{exemplebox}

\paragraph{Éthers-oxydes.} Ils possèdent le groupe \cle{éther} $\mathrm{R-O-R'}$ (oxygène lié à deux carbones). Nomenclature : \emph{alcoxyalcane} (ex. \emph{méthoxypropane} pour \ce{CH3-O-CH2CH2CH3}) ou \emph{éther X-ylique et Y-lique}.

\courssub{3. Propriétés physiques des alcools et polyalcools}

\paragraph{Liaisons hydrogène.} Le groupe $-\ce{OH}$ est très polarisé ($\delta^-$ sur O, $\delta^+$ sur H). Les molécules d'alcool s'associent par \cle{liaisons hydrogène} intermoléculaires ($\ce{O-H\bond{...}O}$).

\begin{propriete}[Conséquences macroscopiques]
\begin{itemize}
    \item \textbf{Températures d'ébullition ($\theta_{\text{éb}}$) élevées} par rapport aux alcanes de masse molaire comparable (ex. éthanol $\SI{78}{\celsius}$ vs propane $\SI{-42}{\celsius}$).
    \item \textbf{Solubilité dans l'eau} : \textbf{Miscibilité totale} pour les alcools à chaîne courte ($n \le 3$) car ils forment des liaisons H avec l'eau. La solubilité \textbf{diminue} quand la chaîne carbonée hydrophobe s'allonge ($n \ge 4$).
    \item \textbf{Effet de ramification} : À formule brute égale, un alcool ramifié a une $\theta_{\text{éb}}$ \textbf{plus basse} (surface de contact réduite, forces de Van der Waals plus faibles) et une \textbf{meilleure solubilité} (encombrement stérique moindre pour l'hydratation).
\end{itemize}
\end{propriete}

\begin{experience}[Comparaison $\theta_{\text{éb}}$ : éthanol, butan-1-ol, 2-méthylpropan-2-ol, propane]
\begin{popfigure}
\begin{tikzpicture}[scale=0.8]
\begin{axis}[
    width=\linewidth, height=6cm,
    ybar, bar width=12pt,
    ylabel={$\theta_{\text{éb}}$ (\si{\celsius})},
    symbolic x coords={Propane, 2-méthylpropan-2-ol, Éthanol, Butan-1-ol},
    xtick=data,
    nodes near coords,
    ymin=-50, ymax=130,
    ymajorgrids=true,
    tick label style={font=\small},
    every node near coord/.append style={font=\small, anchor=south},
    fill=popblue!70
]
\addplot coordinates {(Propane,-42) (2-méthylpropan-2-ol,83) (Éthanol,78) (Butan-1-ol,118)};
\end{axis}
\end{tikzpicture}
\legende{Températures d'ébullition : le propane (gaz) n'a pas de liaison H ; l'éthanol (petite chaîne) bout plus bas que le butan-1-ol (chaîne longue) ; le tert-butanol (ramifié) bout plus bas que le butan-1-ol (linéaire).}
\end{popfigure}
\end{experience}

\paragraph{Polyalcools.} Molécules portant plusieurs groupes $-\ce{OH}$.
\begin{itemize}
    \item \textbf{Glycol (éthane-1,2-diol)} $\ce{HO-CH2-CH2-OH}$ : antigel, synthèse polymères.
    \item \textbf{Glycérol (propane-1,2,3-triol)} $\ce{HO-CH2-CH(OH)-CH2-OH}$ : \textbf{visqueux}, hygroscopique, miscible à l'eau. Produit lors de la \textbf{saponification} des huiles (huile de palme + soude $\rightarrow$ savon + glycérol) dans les savonneries artisanales.
\end{itemize}

\begin{savaistu}[Le glycérol au Cameroun]
Dans les savonneries artisanales (ex. savon de Marseille local, savon noir), l'hydrolyse basique de l'huile de palme (triglycérides) libère du glycérol. Ce dernier adoucit la peau et est récupéré dans l'industrie cosmétique. Le glycol, lui, est toxique par ingestion (danger pour les freins automobiles).
\end{savaistu}

\courssub{4. Aldéhydes et cétones : formules et nomenclature}

Le groupe \cle{carbonyle} $\ce{C=O}$ est plan, polarisé ($\delta^+$ sur C, $\delta^-$ sur O), site de réactivité majeur (addition nucléophile).

\begin{definition}[Aldéhyde et Cétone]
\begin{itemize}
    \item \textbf{Aldéhyde} : $R-\mathrm{CHO}$ (R = H ou alkyl). Le carbone carbonyle est terminal. Suffixe \textbf{-al}. Nomenclature : chaîne la plus longue contenant le $\mathrm{CHO}$, numérotée depuis le carbonyle (locant 1 omis).
    \item \textbf{Cétone} : $R-\mathrm{CO}-R'$ (R, R' = alkyl). Le carbone carbonyle est interne. Suffixe \textbf{-one}. Nomenclature : chaîne la plus longue contenant le $\mathrm{C=O}$, numérotée pour donner le plus petit locant au carbonyle.
\end{itemize}
\end{definition}

\begin{exemplebox}[Nomenclature]
\begin{itemize}
    \item $\ce{CH3CHO}$ : Éthanal (ancien : acétaldéhyde).
    \item $\ce{CH3CH2CHO}$ : Propanal.
    \item $\ce{CH3COCH3}$ : Propan-2-one (ancien : acétone).
    \item $\ce{CH3CH2COCH2CH3}$ : Pentan-3-one.
\end{itemize}
\end{exemplebox}

\paragraph{Isomérie.} Aldéhydes et cétones sont \textbf{isomères de fonction} (même formule brute $C_nH_{2n}O$, groupes fonctionnels différents). Exemple : Propanal ($\ce{C3H6O}$) et Propanone ($\ce{C3H6O}$).

\courssub{5. Tests d'identification des aldéhydes et cétones}

Ces tests exploitent la \textbf{réductibilité} du groupe carbonyle des aldéhydes (présence d'un H sur le C=O) vs l'absence de réductibilité facile pour les cétones.

\begin{experience}[Test à la 2,4-DNPH (2,4-dinitrophénylhydrazine)]
\textbf{Dispositif :} Quelques gouttes de composé carbonylé + solution de 2,4-DNPH dans l'éthanol acide.
\textbf{Observation :} \textbf{Précipité cristallisé jaune-orangé} (2,4-dinitrophénylhydrazone) pour \textbf{aldéhydes ET cétones}.
\textbf{Équation :} $\ce{R2C=O + H2N-NH-C6H3(NO2)2 -> R2C=N-NH-C6H3(NO2)2 v + H2O}$
\textbf{Utilité :} \textbf{Test de caractérisation} (point de fusion du dérivé) mais \textbf{PAS de distinction} aldéhyde/cétone.
\end{experience}

\begin{experience}[Test de Schiff (fuchsine décolorée par $\ce{SO2}$)]
\textbf{Dispositif :} Réactif incolore + composé carbonylé (à froid).
\textbf{Observation :} Coloration \textbf{rose-violet (magenta)} immédiate pour les \textbf{aldéhydes}. \textbf{Négatif} (reste incolore) pour les cétones.
\textbf{Mécanisme :} Addition du $\ce{SO2}$ sur le carbonyle, puis restauration du chromophore fuchsine par l'aldéhyde.
\textbf{Attention :} Réversible à chaud ou par excès d'acide fort.
\end{experience}

\begin{experience}[Test de Fehling (cuivre(II) complexé par tartrate, milieu basique)]
\textbf{Dispositif :} Mélange Fehling A ($\ce{CuSO4}$) + Fehling B (tartrate + $\ce{NaOH}$) + aldéhyde $\rightarrow$ chauffage.
\textbf{Observation :} \textbf{Précipité rouge brique} d'oxyde de cuivre(I) $\ce{Cu2O}$ pour les \textbf{aldéhydes}. \textbf{Négatif} (solution bleue inchangée) pour les cétones.
\textbf{Demi-équations :}
\begin{align*}
\text{Oxydation : } & \ce{R-CHO + 3 OH- -> R-COO- + 2 H2O + 2 e-} \\
\text{Réduction : } & \ce{2 Cu^2+ + 2 e- -> Cu2O v} \quad (\text{complexé tartrate})
\end{align*}
\textbf{Bilan :} $\ce{R-CHO + 2 Cu^2+ + 4 OH- -> R-COO- + Cu2O v + 2 H2O}$
\end{experience}

\begin{experience}[Test de Tollens (argent(I) complexé par ammoniaque, $\ce{[Ag(NH3)2]+}$)]
\textbf{Dispositif :} Solution de Tollens fraîche + aldéhyde $\rightarrow$ chauffage doux (bain-marie).
\textbf{Observation :} \textbf{Miroir d'argent} (dépôt d'Ag métallique brillant) ou précipité noir pour les \textbf{aldéhydes}. \textbf{Négatif} pour les cétones.
\textbf{Bilan :} $\ce{R-CHO + 2 [Ag(NH3)2]+ + 3 OH- -> R-COO- + 2 Ag v + 4 NH3 + 2 H2O}$
\end{experience}

\begin{aretenir}[Tableau récapitulatif des tests d'identification]
\begin{center}
\begin{tabularx}{\linewidth}{|l|X|X|}
\hline
\rowcolor{popblueL} \textbf{Réactif} & \textbf{Aldéhyde} & \textbf{Cétone} \\
\hline
2,4-DNPH & \textbf{+}\newline Précipité jaune-orangé & \textbf{+}\newline Précipité jaune-orangé \\
\hline
Schiff & \textbf{+}\newline Rose-violet & \textbf{--}\newline Incolore \\
\hline
Fehling & \textbf{+}\newline Précipité rouge brique ($\ce{Cu2O}$) & \textbf{--}\newline Bleu inchangé \\
\hline
Tollens & \textbf{+}\newline Miroir d'argent (Ag) & \textbf{--}\newline Inchangé \\
\hline
\end{tabularx}
\end{center}
\end{aretenir}

\courssub{6. Synthèse et applications au quotidien}

\paragraph{Préparation des aldéhydes et cétones : Oxydation des alcools.}
\begin{methode}[Oxydation ménagée (arrêt à l'aldehyde/cétone)]
\begin{enumerate}
    \item Mélanger l'alcool et l'oxydant ($\ce{K2Cr2O7/H+}$ ou $\ce{KMnO4/H+}$) dans un ballon à reflux.
    \item \textbf{Alcool primaire} $\xrightarrow{\text{distillation immédiate}}$ \textbf{Aldéhyde} (évite la sur-oxydation en acide).
    \item \textbf{Alcool secondaire} $\xrightarrow{\text{reflux}}$ \textbf{Cétone}.
    \item \textbf{Alcool tertiaire} : \textbf{Pas de réaction} (pas de H sur le C fonctionnel).
\end{enumerate}
\end{methode}

\begin{attention}[Sur-oxydation]
Si l'aldehyde n'est pas distillé au fur et à mesure, il s'oxyde en \textbf{acide carboxylique} : $\ce{R-CHO ->[Ox] R-COOH}$. C'est pour cela que le vin tourné (éthanol $\rightarrow$ acide acétique) devient du vinaigre.
\end{attention}

\paragraph{Distillation et bili-bili.} La fabrication traditionnelle du bili-bili (bière de mil/maïs) ou du vin de palme suivi de distillation artisanale (arak, odontol) repose sur la différence de $\theta_{\text{éb}}$ : éthanol ($\SI{78}{\celsius}$) vs eau ($\SI{100}{\celsius}$). \textbf{Attention :} Les premières fractions (têtes) contiennent du \textbf{méthanol} ($\SI{65}{\celsius}$), \textbf{très toxique} (cécité, décès).

\begin{savaistu}[CIMENCAM et SONARA]
À la SONARA (Limbe), le craquage et le reformage produisent des composés oxygénés (éthanol, éthers-oxydes comme MTBE/ETBE) ajoutés aux essences comme \textbf{oxygénants} (améliorent l'indice d'octane, réduisent \ce{CO}/\ce{NO_x}). CIMENCAM utilise des alcools (glycol) dans les adjuvants de ciment.
\end{savaistu}

% ============================================================
% EXERCICES
% ============================================================

\courssub{Je vérifie mes connaissances}

\begin{exercice}[QCM : les groupes fonctionnels]
Pour chaque question, une seule réponse est exacte. Recopie la lettre correspondante et la phrase complète.
\begin{enumerate}
\item Le groupe fonctionnel des alcools est :
\begin{enumerate}
\item $-\ce{CHO}$ \quad \item $-\ce{OH}$ \quad \item $-\ce{CO}-$ \quad \item $-\ce{COOH}$
\end{enumerate}
\item Un alcool est dit secondaire lorsque le carbone fonctionnel est lié à :
\begin{enumerate}
\item un seul atome de carbone \quad \item deux atomes de carbone \quad \item trois atomes de carbone \quad \item aucun atome de carbone
\end{enumerate}
\item La formule brute générale d'un alcool saturé non cyclique est :
\begin{enumerate}
\item $C_nH_{2n}O$ \quad \item $C_nH_{2n+2}O$ \quad \item $C_nH_{2n-2}O$ \quad \item $C_nH_{2n+1}OH$ avec $n=0$ possible
\end{enumerate}
\item Le réactif qui donne un précipité jaune-orangé avec les aldéhydes \emph{et} les cétones est :
\begin{enumerate}
\item la liqueur de Fehling \quad \item le réactif de Schiff \quad \item la 2,4-DNPH \quad \item le réactif de Tollens
\end{enumerate}
\item Le miroir d'argent est un test caractéristique :
\begin{enumerate}
\item des cétones \quad \item des alcools \quad \item des aldéhydes \quad \item des éthers-oxydes
\end{enumerate}
\end{enumerate}
\end{exercice}

\begin{exercice}[Vrai ou faux, en justifiant]
Réponds par VRAI ou FAUX, puis justifie ta réponse en une ou deux phrases.
\begin{enumerate}
\item Le propan-2-ol est un alcool primaire.
\item Les aldéhydes et les cétones possèdent tous deux le groupe carbonyle $\ce{C=O}$.
\item La liqueur de Fehling permet de distinguer un aldéhyde d'une cétone.
\item Le glycérol (propane-1,2,3-triol) est un polyalcool possédant trois groupes hydroxyle.
\item La 2,4-DNPH permet de distinguer le propanal de la propanone.
\item L'éthanol est miscible à l'eau en toutes proportions grâce aux liaisons hydrogène.
\end{enumerate}
\end{exercice}

\begin{exercice}[Texte à trous]
Recopie et complète le texte suivant avec les mots ou expressions de la liste : \emph{carbonyle ; hydroxyle ; primaire ; tertiaire ; rouge brique ; Schiff ; oxydation ménagée ; hydrogène}.

Les composés oxygénés renferment l'élément oxygène dans leur molécule. Les alcools possèdent le groupe \ldots\ldots\ldots $-\ce{OH}$, tandis que les aldéhydes et les cétones possèdent le groupe \ldots\ldots\ldots $\ce{C=O}$. Dans un alcool \ldots\ldots\ldots, le carbone fonctionnel n'est lié qu'à un seul atome de carbone ; dans un alcool \ldots\ldots\ldots, il est lié à trois atomes de carbone. La solubilité des petits alcools dans l'eau s'explique par l'existence de liaisons \ldots\ldots\ldots entre leurs molécules et celles de l'eau. Lors du test à la liqueur de Fehling, un aldéhyde donne un précipité \ldots\ldots\ldots d'oxyde de cuivre(I) \ce{Cu2O}. Le réactif de \ldots\ldots\ldots, incolore, devient rose-violet en présence d'un aldéhyde. Les aldéhydes s'obtiennent par \ldots\ldots\ldots des alcools primaires.
\end{exercice}

\begin{exercice}[Définitions]
Définis en une phrase claire et précise chacun des termes suivants :
\begin{enumerate}
\item composé oxygéné ;
\item alcool secondaire ;
\item aldéhyde ;
\item test d'identification ;
\item polyalcool.
\end{enumerate}
\end{exercice}

\courssub{Je m'entraîne}

\begin{exercice}[Nomenclature dans les deux sens]
\textbf{Partie A.} Donne le nom de chacun des composés suivants :
\begin{enumerate}
\item \chemfig{CH_3-CH_2-CH(-[2]OH)-CH_3}
\item \chemfig{CH_3-C(-[2]CH_3)(-[6]CH_3)-OH}
\item \chemfig{CH_3-CH_2-O-CH_2-CH_3}
\item \chemfig{CH_3-CH_2-C(=[1]O)-[7]H}
\item \chemfig{CH_3-C(=[1]O)-[7]CH_3}
\item \chemfig{CH_3-CH_2-CH_2-CH_2-C(=[1]O)-[7]H}
\end{enumerate}
\textbf{Partie B.} Écris la formule semi-développée de chacun des composés suivants :
\begin{enumerate}
\item le 3-méthylbutan-2-ol ;
\item le méthanal ;
\item la pentan-3-one ;
\item l'éther méthylique et propylique (méthoxypropane).
\end{enumerate}
\end{exercice}

\begin{exercice}[Classes des alcools]
On considère les quatre alcools suivants :
\begin{center}
A : \chemfig{CH_3-CH_2-CH_2-CH_2-OH} \qquad
B : \chemfig{CH_3-CH(-[2]OH)-CH_2-CH_3} \qquad
C : \chemfig{CH_3-C(-[2]CH_3)(-[6]OH)-CH_3} \qquad
D : \chemfig{CH_3-CH(-[2]CH_3)-CH_2-OH}
\end{center}
\begin{enumerate}
\item Indique, pour chaque alcool, sa classe (primaire, secondaire ou tertiaire) en justifiant.
\item Nomme chacun de ces alcools.
\item Quels sont, parmi ces composés, les isomères de constitution ? Précise leur formule brute.
\end{enumerate}
\end{exercice}

\begin{exercice}[Températures d'ébullition]
On donne les températures d'ébullition (sous pression normale) de quatre composés, dans le désordre : \SI{78}{\celsius}, \SI{118}{\celsius}, \SI{-42}{\celsius}, \SI{83}{\celsius}.
\begin{enumerate}
\item Attribue à chacun des composés suivants sa température d'ébullition : propane \ce{C3H8}, éthanol \ce{C2H5OH}, butan-1-ol \ce{C4H9OH}, 2-méthylpropan-2-ol \ce{C4H9OH}.
\item Justifie tes choix en faisant intervenir les liaisons hydrogène, la longueur de la chaîne carbonée et la ramification.
\item Lequel de ces composés est le plus soluble dans l'eau ? Justifie.
\end{enumerate}
\end{exercice}

\begin{exercice}[Tableau des tests]
Recopie et complète le tableau suivant en indiquant, pour chaque test, s'il est positif (+) ou négatif ($-$), et l'observation correspondante en cas de test positif.

\begin{center}
\begin{tabularx}{\linewidth}{|l|X|X|X|}
\hline
\rowcolor{popblueL} \textbf{Réactif} & \textbf{Propanal} & \textbf{Butanone} & \textbf{Éthanol} \\
\hline
2,4-DNPH & & & \\
\hline
Réactif de Schiff & & & \\
\hline
Liqueur de Fehling & & & \\
\hline
Réactif de Tollens & & & \\
\hline
\end{tabularx}
\end{center}

Deux flacons non étiquetés contiennent deux liquides inconnus X et Y. Le liquide X donne un précipité jaune avec la 2,4-DNPH mais ne réagit ni avec la liqueur de Fehling ni avec le réactif de Tollens. Le liquide Y donne un précipité jaune avec la 2,4-DNPH et un dépôt d'argent avec le réactif de Tollens. Détermine la famille chimique de X et celle de Y.
\end{exercice}

\courssub{Je me prépare au Probatoire}

\begin{exercice}[Le composé \ce{C4H10O} : isomères et identification]
Un laboratoire de lycée à Douala possède un flacon étiqueté simplement « \ce{C4H10O} ».

\textbf{Données :} masses molaires atomiques : $M(\ce{C}) = \SI{12}{g.mol^{-1}}$ ; $M(\ce{H}) = \SI{1}{g.mol^{-1}}$ ; $M(\ce{O}) = \SI{16}{g.mol^{-1}}$.
\begin{enumerate}
\item Calcule la masse molaire moléculaire du composé.
\item Écris les formules semi-développées de \textbf{tous} les isomères de formule brute \ce{C4H10O} appartenant à la famille des alcools. Nomme-les et précise la classe de chacun.
\item Écris les formules semi-développées de tous les isomères éthers-oxydes de formule \ce{C4H10O} et nomme-les.
\item Une oxydation ménagée du contenu du flacon par une solution acidifiée de dichromate de potassium fournit un composé B qui donne un précipité jaune avec la 2,4-DNPH et un précipité rouge brique avec la liqueur de Fehling.
\begin{enumerate}
\item Quelle est la famille du composé B ? Justifie.
\item En déduire la classe de l'alcool contenu dans le flacon, puis identifier les alcools possibles.
\item Écris la formule semi-développée et le nom de B sachant que sa chaîne carbonée est linéaire.
\end{enumerate}
\item Le réactif de Schiff, rosé par B, redevient-il incolore si l'on ajoute un excès d'acide sulfurique concentré ? Que peut-on en conclure sur la nature de B ?
\end{enumerate}
\end{exercice}

\begin{exercice}[Identification d'un composé \ce{C3H6O} et tests de Fehling]
Au cours d'une séance de travaux pratiques au Lycée de Nkongsamba, un professeur remet aux élèves un échantillon d'un composé organique B de formule brute \ce{C3H6O}. Les élèves réalisent trois tests :
\begin{itemize}
\item test 1 : B donne avec la 2,4-DNPH un précipité jaune-orangé ;
\item test 2 : B chauffé avec la liqueur de Fehling donne un précipité rouge brique ;
\item test 3 : B avec le réactif de Tollens donne un miroir d'argent.
\end{itemize}

\textbf{Données :} $M(\ce{C}) = \SI{12}{g.mol^{-1}}$ ; $M(\ce{H}) = \SI{1}{g.mol^{-1}}$ ; $M(\ce{O}) = \SI{16}{g.mol^{-1}}$ ; $M(\ce{Cu}) = \SI{63,5}{g.mol^{-1}}$.
\begin{enumerate}
\item Écris les formules semi-développées de tous les composés de formule brute \ce{C3H6O} possédant un groupe carbonyle. Nomme-les.
\item Interprète le résultat du test 1. Que peut-on en déduire ?
\item Interprète les tests 2 et 3. Identifie alors sans ambiguïté le composé B : donne sa formule semi-développée et son nom.
\item Écris l'équation-bilan de la réaction de B avec la liqueur de Fehling, sachant que l'ion cuivre(II) complexé \ce{Cu^2+} est réduit en oxyde de cuivre(I) \ce{Cu2O} et que B est oxydé en ion carboxylate \ce{C2H5COO-} en milieu basique.
\item Sachant que l'on a utilisé $n = \SI{0,010}{mol}$ de B, calcule la quantité de matière puis la masse du précipité de \ce{Cu2O} obtenu, la réaction étant totale.
\item Propose une expérience simple permettant de préparer B à partir de l'alcool correspondant. Précise le nom et la classe de cet alcool.
\end{enumerate}
\end{exercice}

\begin{exercice}[Alcool frelaté : le danger du méthanol]
Suite à plusieurs cas d'intoxication liés à la consommation d'alcool artisanal frelaté dans une localité de la région de l'Ouest, un laboratoire d'analyses reçoit un échantillon de boisson suspectée de contenir du méthanol \ce{CH3OH} mélangé à l'éthanol \ce{C2H5OH}.

\textbf{Données :} températures d'ébullition : méthanol $\theta_{\text{éb}} = \SI{65}{\celsius}$ ; éthanol $\theta_{\text{éb}} = \SI{78}{\celsius}$. Masses molaires : $M(\ce{CH3OH}) = \SI{32}{g.mol^{-1}}$ ; $M(\ce{C2H5OH}) = \SI{46}{g.mol^{-1}}$.
\begin{enumerate}
\item Écris les formules semi-développées du méthanol et de l'éthanol. À quelle classe d'alcool appartiennent-ils ?
\item On réalise l'oxydation ménagée de chacun de ces deux alcools par une solution acidifiée de permanganate de potassium.
\begin{enumerate}
\item Nomme les produits d'oxydation obtenus respectivement à partir du méthanol et de l'éthanol (produits de la première étape d'oxydation).
\item Explique pourquoi le méthanol et l'éthanol, tous deux alcools primaires, donnent les \emph{mêmes} résultats aux tests d'oxydation (2,4-DNPH, Schiff, Fehling) : ces tests permettent-ils de distinguer les deux alcools ?
\end{enumerate}
\item Les produits d'oxydation obtenus sont le méthanal et l'éthanal. Propose une démarche expérimentale, fondée sur une propriété physique ou sur une réaction ultérieure d'oxydation, permettant de distinguer ces deux produits d'oxydation.
\item En exploitant les températures d'ébullition données, nomme la technique de séparation qui permettrait de mettre en évidence la présence de méthanol dans la boisson. Décris brièvement le principe de cette technique et précise quel liquide s'évapore en premier.
\item Le méthanol est toxique : son oxydation dans l'organisme produit des composés qui attaquent le nerf optique. Rédige un court message de sensibilisation (5 lignes maximum) à l'attention des consommateurs, expliquant le danger des alcools frelatés.
\end{enumerate}
\end{exercice}

\newpage

\coursec{Corrigés des exercices}

\begin{corrige}[Exercice 1 — QCM : les groupes fonctionnels]
\begin{enumerate}
\item \textbf{b)} Le groupe fonctionnel des alcools est le groupe \cle{hydroxyle} $-\ce{OH}$, porté par un carbone \cle{tétragonal} (hybridation $sp^3$).
\item \textbf{b)} Un alcool est \cle{secondaire} lorsque le carbone fonctionnel (celui qui porte le $-\ce{OH}$) est lié à \textbf{deux} atomes de carbone (et à un seul atome d'hydrogène). Sa formule générale est $R-\ce{CH}(OH)-R'$.
\item \textbf{b)} La formule brute générale d'un alcool saturé non cyclique est $C_nH_{2n+2}O$ avec $n \ge 1$ (l'alcool le plus simple est le méthanol \ce{CH4O}, donc $n=0$ est impossible).
\item \textbf{c)} La \cle{2,4-DNPH} donne un précipité jaune-orangé avec \emph{tous} les composés carbonylés : aldéhydes \textbf{et} cétones. C'est un test de la \cle{famille carbonyle}, pas un test de distinction.
\item \textbf{c)} Le \cle{miroir d'argent} (réactif de Tollens) est caractéristique des \textbf{aldéhydes}, seuls réducteurs parmi les composés étudiés dans cette leçon.
\end{enumerate}
\textbf{Point de vigilance :} ne pas confondre « test de caractérisation du carbonyle » (2,4-DNPH) et « tests de distinction aldéhyde/cétone » (Schiff, Fehling, Tollens).
\end{corrige}

\begin{corrige}[Exercice 2 — Vrai ou faux, en justifiant]
\begin{enumerate}
\item \textbf{FAUX.} Dans le propan-2-ol \chemfig{CH_3-CH(-[2]OH)-CH_3}, le carbone fonctionnel est lié à \textbf{deux} atomes de carbone : c'est un alcool \textbf{secondaire}.
\item \textbf{VRAI.} Les aldéhydes ($R-\ce{CHO}$) et les cétones ($R-\ce{CO}-R'$) possèdent tous deux le groupe \cle{carbonyle} $\ce{C=O}$ ; ils diffèrent par la position de ce groupe (terminal pour l'aldéhyde, interne pour la cétone).
\item \textbf{VRAI.} Un aldéhyde, \cle{réducteur}, réduit les ions \ce{Cu^2+} de la liqueur de Fehling en précipité rouge brique d'oxyde de cuivre(I) \ce{Cu2O} ; une cétone ne réagit pas (la solution reste bleue).
\item \textbf{VRAI.} Le glycérol \ce{HO-CH2-CH(OH)-CH2-OH} porte \textbf{trois} groupes hydroxyle : c'est un \cle{trialcool}, donc un \cle{polyalcool}.
\item \textbf{FAUX.} Le propanal (aldéhyde) et la propanone (cétone) donnent \emph{tous deux} un précipité jaune-orangé avec la 2,4-DNPH : ce test ne les distingue pas. Il faut utiliser Schiff, Fehling ou Tollens.
\item \textbf{VRAI.} L'éthanol forme des \cle{liaisons hydrogène} avec les molécules d'eau grâce à son groupe $-\ce{OH}$ ; sa chaîne carbonée étant courte ($n=2$), il est \cle{miscible} à l'eau en toutes proportions.
\end{enumerate}
\end{corrige}

\begin{corrige}[Exercice 3 — Texte à trous]
Les \cle{composés oxygénés} renferment l'élément oxygène dans leur molécule. Les alcools possèdent le groupe \cle{hydroxyle} $-\ce{OH}$, tandis que les aldéhydes et les cétones possèdent le groupe \cle{carbonyle} $\ce{C=O}$. Dans un alcool \cle{primaire}, le carbone fonctionnel n'est lié qu'à un seul atome de carbone ; dans un alcool \cle{tertiaire}, il est lié à trois atomes de carbone. La \cle{solubilité} des petits alcools dans l'eau s'explique par l'existence de liaisons \cle{hydrogène} entre leurs molécules et celles de l'eau. Lors du test à la liqueur de Fehling, un aldéhyde donne un précipité \textbf{rouge brique} d'oxyde de cuivre(I) \ce{Cu2O}. Le réactif de \cle{Schiff}, incolore, devient rose-violet en présence d'un aldéhyde. Les aldéhydes s'obtiennent par \cle{oxydation ménagée} des alcools primaires.
\end{corrige}

\begin{corrige}[Exercice 4 — Définitions]
\begin{enumerate}
\item \cle{Composé oxygéné} : composé organique dont la molécule contient au moins un atome d'oxygène lié à la chaîne carbonée (dérivé d'hydrocarbure portant un groupe fonctionnel oxygéné).
\item \cle{Alcool secondaire} : alcool dont le carbone fonctionnel (porteur du groupe $-\ce{OH}$) est lié à \textbf{deux} atomes de carbone et à un atome d'hydrogène : $R-\ce{CH}(OH)-R'$.
\item \cle{Aldéhyde} : composé oxygéné possédant un groupe carbonyle $\ce{C=O}$ porté par un carbone \textbf{terminal} de la chaîne, lié à au moins un atome d'hydrogène : $R-\ce{CHO}$.
\item \cle{Test d'identification} : expérience simple dont le résultat observable (coloration, précipité, dépôt) permet de reconnaître de manière caractéristique une famille ou une espèce chimique.
\item \cle{Polyalcool} : molécule organique possédant \textbf{plusieurs} groupes hydroxyle $-\ce{OH}$ portés par des atomes de carbone différents (ex. glycol, glycérol).
\end{enumerate}
\end{corrige}

\begin{corrige}[Exercice 5 — Nomenclature dans les deux sens]
\textbf{Partie A.}
\begin{enumerate}
\item \chemfig{CH_3-CH_2-CH(-[2]OH)-CH_3} : chaîne de 4 C, $-\ce{OH}$ en position 2 $\Rightarrow$ \textbf{butan-2-ol} (alcool secondaire).
\item \chemfig{CH_3-C(-[2]CH_3)(-[6]CH_3)-OH} : chaîne principale de 3 C, deux méthyles sur le C2 $\Rightarrow$ \textbf{2-méthylpropan-2-ol} (alcool tertiaire).
\item \chemfig{CH_3-CH_2-O-CH_2-CH_3} : éther-oxyde symétrique $\Rightarrow$ \textbf{éthoxyéthane} (éther éthylique).
\item \chemfig{CH_3-CH_2-C(=[1]O)-[7]H} : aldéhyde à 3 C $\Rightarrow$ \textbf{propanal}.
\item \chemfig{CH_3-C(=[1]O)-[7]CH_3} : cétone à 3 C $\Rightarrow$ \textbf{propanone} (propan-2-one).
\item \chemfig{CH_3-CH_2-CH_2-CH_2-C(=[1]O)-[7]H} : aldéhyde à 5 C $\Rightarrow$ \textbf{pentanal}.
\end{enumerate}
\textbf{Partie B.}
\begin{enumerate}
\item 3-méthylbutan-2-ol : \chemfig{CH_3-CH(-[2]OH)-CH(-[6]CH_3)-CH_3}
\item Méthanal : \chemfig{H-C(=[1]O)-[7]H} \quad (\ce{HCHO})
\item Pentan-3-one : \chemfig{CH_3-CH_2-C(=[1]O)-[7]CH_2-CH_3}
\item Méthoxypropane : \chemfig{CH_3-O-CH_2-CH_2-CH_3}
\end{enumerate}
\textbf{Point de vigilance :} pour une cétone, numéroter la chaîne de façon à donner le plus petit locant au carbonyle ; pour un aldéhyde, le carbone du $-\ce{CHO}$ porte toujours le numéro 1 (locant omis).
\end{corrige}

\begin{corrige}[Exercice 6 — Classes des alcools]
\textbf{1. Classes.}
\begin{itemize}
\item \textbf{A} : \ce{CH3-CH2-CH2-CH2-OH} : le C fonctionnel est lié à \textbf{un} atome de carbone $\Rightarrow$ alcool \cle{primaire}.
\item \textbf{B} : \ce{CH3-CH(OH)-CH2-CH3} : le C fonctionnel est lié à \textbf{deux} atomes de carbone $\Rightarrow$ alcool \cle{secondaire}.
\item \textbf{C} : \ce{CH3-C(CH3)(OH)-CH3} : le C fonctionnel est lié à \textbf{trois} atomes de carbone $\Rightarrow$ alcool \cle{tertiaire}.
\item \textbf{D} : \ce{CH3-CH(CH3)-CH2-OH} : le C fonctionnel (celui du \ce{CH2-OH}) est lié à \textbf{un} atome de carbone $\Rightarrow$ alcool \cle{primaire}.
\end{itemize}
\textbf{2. Noms.} A : \textbf{butan-1-ol} ; B : \textbf{butan-2-ol} ; C : \textbf{2-méthylpropan-2-ol} ; D : \textbf{2-méthylpropan-1-ol}.

\textbf{3. Isomérie.} Les quatre composés ont la même formule brute \ce{C4H10O} : ce sont des \cle{isomères de constitution} (isomérie de position pour A et B ; isomérie de chaîne pour A, B et C, D).
\end{corrige}

\begin{corrige}[Exercice 7 — Températures d'ébullition et solubilité]
\textbf{1. Attribution.}
\begin{center}
\begin{tabularx}{\linewidth}{|l|X|c|}
\hline
\rowcolor{popblueL} \textbf{Composé} & \textbf{Justification} & $\theta_{\text{éb}}$ \\
\hline
Propane \ce{C3H8} & pas de liaison H (alcane) & \SI{-42}{\celsius} \\
\hline
Éthanol \ce{C2H5OH} & liaison H, chaîne courte ($n=2$) & \SI{78}{\celsius} \\
\hline
2-méthylpropan-2-ol \ce{C4H9OH} & liaison H, chaîne de 4 C \textbf{ramifiée} & \SI{83}{\celsius} \\
\hline
Butan-1-ol \ce{C4H9OH} & liaison H, chaîne de 4 C \textbf{linéaire} & \SI{118}{\celsius} \\
\hline
\end{tabularx}
\end{center}
\textbf{2. Justification.} Le propane ne forme pas de liaisons hydrogène : ses molécules ne sont retenues que par de faibles forces de Van der Waals, d'où une $\theta_{\text{éb}}$ très basse. Les alcools s'associent par liaisons hydrogène, ce qui élève fortement $\theta_{\text{éb}}$. À fonction identique, plus la chaîne carbonée est longue, plus $\theta_{\text{éb}}$ est élevée (butan-1-ol $>$ éthanol). À formule brute égale (\ce{C4H10O}), l'isomère \textbf{ramifié} (2-méthylpropan-2-ol) présente une surface de contact moindre entre molécules : ses forces intermoléculaires sont plus faibles et il bout plus bas (\SI{83}{\celsius} $<$ \SI{118}{\celsius}).

\textbf{3. Solubilité.} L'\textbf{éthanol} est le plus soluble dans l'eau (miscible en toutes proportions) : sa chaîne hydrophobe est la plus courte et son groupe $-\ce{OH}$ forme des liaisons hydrogène avec l'eau.
\end{corrige}

\begin{corrige}[Exercice 8 — Tableau des tests d'identification]
\begin{center}
\begin{tabularx}{\linewidth}{|l|X|X|X|}
\hline
\rowcolor{popblueL} \textbf{Réactif} & \textbf{Propanal (aldéhyde)} & \textbf{Butanone (cétone)} & \textbf{Éthanol (alcool)} \\
\hline
2,4-DNPH & \textbf{+} précipité jaune-orangé & \textbf{+} précipité jaune-orangé & \textbf{--} \\
\hline
Schiff & \textbf{+} rose-violet & \textbf{--} & \textbf{--} \\
\hline
Fehling & \textbf{+} précipité rouge brique \ce{Cu2O} & \textbf{--} & \textbf{--} \\
\hline
Tollens & \textbf{+} miroir d'argent & \textbf{--} & \textbf{--} \\
\hline
\end{tabularx}
\end{center}
\textbf{Identification de X et Y.}
\begin{itemize}
\item \textbf{X} : DNPH positive $\Rightarrow$ composé \cle{carbonylé} ; Fehling et Tollens négatifs $\Rightarrow$ ce n'est \textbf{pas} un aldéhyde. Donc \textbf{X est une cétone}.
\item \textbf{Y} : DNPH positive $\Rightarrow$ composé carbonylé ; Tollens positif (miroir d'argent) $\Rightarrow$ \cle{réducteur}. Donc \textbf{Y est un aldéhyde}.
\end{itemize}
\textbf{Point de vigilance :} l'éthanol ne réagit avec \emph{aucun} de ces quatre réactifs : les tests de cette leçon caractérisent le groupe carbonyle, pas le groupe hydroxyle.
\end{corrige}

\begin{corrige}[Exercice 9 — Le composé \ce{C4H10O} : isomères et identification]
\emph{Barème indicatif (sur 10) : 1. (1 pt) ; 2. (3 pts) ; 3. (1,5 pt) ; 4. (3 pts) ; 5. (1,5 pt).}

\textbf{1. Masse molaire.}
\[ M(\ce{C4H10O}) = 4\times 12 + 10\times 1 + 16 = 48 + 10 + 16 = \boxed{\SI{74}{g.mol^{-1}}} \]

\textbf{2. Isomères alcools de \ce{C4H10O}.}
\begin{center}
\begin{tabularx}{\linewidth}{|X|l|l|}
\hline
\rowcolor{popblueL} \textbf{Formule semi-développée} & \textbf{Nom} & \textbf{Classe} \\
\hline
\ce{CH3-CH2-CH2-CH2-OH} & butan-1-ol & primaire \\
\hline
\ce{CH3-CH2-CH(OH)-CH3} & butan-2-ol & secondaire \\
\hline
\ce{CH3-CH(CH3)-CH2-OH} & 2-méthylpropan-1-ol & primaire \\
\hline
\ce{CH3-C(CH3)(OH)-CH3} & 2-méthylpropan-2-ol & tertiaire \\
\hline
\end{tabularx}
\end{center}

\textbf{3. Isomères éthers-oxydes.}
\begin{itemize}
\item \ce{CH3-O-CH2-CH2-CH3} : méthoxypropane ;
\item \ce{CH3-O-CH(CH3)-CH3} : 2-méthoxypropane ;
\item \ce{CH3-CH2-O-CH2-CH3} : éthoxyéthane.
\end{itemize}

\textbf{4. Identification.}
\begin{enumerate}
\item B donne un précipité jaune avec la 2,4-DNPH : B est un \textbf{composé carbonylé}. De plus, B donne un précipité rouge brique avec la liqueur de Fehling : B est \textbf{réducteur}. Conclusion : \textbf{B est un aldéhyde}.
\item Seul un \textbf{alcool primaire} conduit par oxydation ménagée à un aldéhyde (un alcool secondaire donnerait une cétone, négative au test de Fehling ; un alcool tertiaire ne s'oxyde pas). L'alcool du flacon est donc le \textbf{butan-1-ol} ou le \textbf{2-méthylpropan-1-ol}.
\item B a une chaîne carbonée \textbf{linéaire} : B est le \textbf{butanal} : \chemfig{CH_3-CH_2-CH_2-C(=[1]O)-[7]H}. L'alcool du flacon est donc le butan-1-ol.
\end{enumerate}

\textbf{5. Réactif de Schiff.} Oui : la coloration rose du réactif de Schiff est \textbf{réversible} ; un excès d'acide sulfurique concentré (ou un chauffage) la fait disparaître. Cela confirme que B est bien un \textbf{aldéhyde} (seuls les aldéhydes donnent ce test réversible).

\textbf{Points de vigilance :} ne pas oublier l'isomère 2-méthylpropan-1-ol (primaire malgré la ramification) ; rappeler que la 2,4-DNPH seule ne suffit pas à conclure.
\end{corrige}

\begin{corrige}[Exercice 10 — Identification d'un composé \ce{C3H6O} et tests de Fehling]
\emph{Barème indicatif (sur 10) : 1. (1,5 pt) ; 2. (1 pt) ; 3. (1,5 pt) ; 4. (2 pts) ; 5. (2,5 pts) ; 6. (1,5 pt).}

\textbf{1. Composés carbonylés de formule \ce{C3H6O}.}
\begin{itemize}
\item \ce{CH3-CH2-CHO} : \textbf{propanal} (aldéhyde) ;
\item \ce{CH3-CO-CH3} : \textbf{propanone} (cétone).
\end{itemize}

\textbf{2. Test 1.} Le précipité jaune-orangé avec la 2,4-DNPH prouve que B possède un \textbf{groupe carbonyle} : B est un aldéhyde \emph{ou} une cétone. Ce test seul ne permet pas de trancher.

\textbf{3. Tests 2 et 3.} Le précipité rouge brique avec Fehling et le miroir d'argent avec Tollens montrent que B est \textbf{réducteur} : B est un \textbf{aldéhyde}. Conclusion : B est le \textbf{propanal}, \chemfig{CH_3-CH_2-C(=[1]O)-[7]H}.

\textbf{4. Équation-bilan avec la liqueur de Fehling.}
\begin{align*}
\text{Oxydation : } & \ce{CH3CH2CHO + 3 OH- -> CH3CH2COO- + 2 H2O + 2 e-} \\
\text{Réduction : } & \ce{2 Cu^2+ + 2 OH- + 2 e- -> Cu2O + H2O}
\end{align*}
Bilan :
\[ \ce{CH3CH2CHO + 2 Cu^2+ + 5 OH- -> CH3CH2COO- + Cu2O v + 3 H2O} \]

\textbf{5. Quantité et masse de \ce{Cu2O}.} Tableau d'avancement (réaction totale, B limitant) :
\begin{center}
\begin{tabularx}{\linewidth}{|l|X|c|c|}
\hline
\rowcolor{popblueL} \textbf{État} & \textbf{Avancement} & $n(\ce{CH3CH2CHO})$ & $n(\ce{Cu2O})$ \\
\hline
Initial & $x=0$ & \SI{0,010}{mol} & $0$ \\
\hline
Final & $x_{\max}$ & $0$ & $x_{\max}$ \\
\hline
\end{tabularx}
\end{center}
D'après la stœchiométrie : $n(\ce{Cu2O}) = n(\ce{B}) = x_{\max} = \SI{0,010}{mol}$.
\[ M(\ce{Cu2O}) = 2 \times 63{,}5 + 16 = \SI{143}{g.mol^{-1}} \]
\[ m(\ce{Cu2O}) = n \times M = 0{,}010 \times 143 = \boxed{\SI{1,43}{g}} \]

\textbf{6. Préparation de B.} On réalise l'\textbf{oxydation ménagée} du \textbf{propan-1-ol} \ce{CH3CH2CH2OH} (alcool \textbf{primaire}) par une solution acidifiée de dichromate de potassium (\ce{K2Cr2O7}/\ce{H+}), en \textbf{distillant le propanal au fur et à mesure} de sa formation pour éviter sa sur-oxydation en acide propanoïque.

\textbf{Points de vigilance :} vérifier l'équilibrage des demi-équations en milieu basique (avec \ce{OH-}) ; dans le calcul, ne pas oublier le facteur 2 sur le cuivre dans $M(\ce{Cu2O})$.
\end{corrige}

\begin{corrige}[Exercice 11 — Alcool frelaté : le danger du méthanol]
\emph{Barème indicatif (sur 10) : 1. (1,5 pt) ; 2. (3 pts) ; 3. (1,5 pt) ; 4. (2 pts) ; 5. (2 pts).}

\textbf{1. Formules et classes.} Méthanol : \chemfig{CH_3-OH} ; éthanol : \chemfig{CH_3-CH_2-OH}. Dans les deux cas, le carbone fonctionnel n'est lié qu'à \textbf{un seul} atome de carbone (zéro pour le méthanol, considéré par convention comme primaire) : ce sont des \textbf{alcools primaires}.

\textbf{2. Oxydation ménagée.}
\begin{enumerate}
\item Méthanol $\rightarrow$ \textbf{méthanal} \ce{HCHO} (formaldéhyde) ; éthanol $\rightarrow$ \textbf{éthanal} \ce{CH3CHO} (acétaldéhyde).
\item Les deux alcools étant \textbf{primaires}, ils donnent tous deux des \textbf{aldéhydes} : les produits obtenus réagissent positivement à la 2,4-DNPH (précipité jaune), au réactif de Schiff (rose-violet) et à la liqueur de Fehling (précipité rouge brique). Ces tests caractérisent la \textbf{famille} (aldéhyde) mais pas la \textbf{longueur de la chaîne} : ils ne permettent \textbf{pas} de distinguer le méthanol de l'éthanol.
\end{enumerate}

\textbf{3. Distinction du méthanal et de l'éthanal.} Deux démarches possibles :
\begin{itemize}
\item \textbf{Propriété physique :} mesurer la température d'ébullition : le méthanal est gazeux à température ambiante ($\theta_{\text{éb}} \approx \SI{-21}{\celsius}$) alors que l'éthanal bout vers \SI{21}{\celsius} ;
\item \textbf{Oxydation ultérieure :} oxyder chaque aldéhyde en acide carboxylique : le méthanal donne l'\textbf{acide méthanoïque} \ce{HCOOH}, qui possède encore un caractère réducteur (il réduit le réactif de Tollens), contrairement à l'acide éthanoïque \ce{CH3COOH} issu de l'éthanal.
\end{itemize}

\textbf{4. Technique de séparation.} On utilise la \textbf{distillation fractionnée}, qui sépare les liquides miscibles selon leurs températures d'ébullition. En chauffant le mélange, le liquide le plus \textbf{volatil} s'évapore en premier : c'est le \textbf{méthanol} ($\SI{65}{\celsius} < \SI{78}{\celsius}$). En recueillant la première fraction distillée et en mesurant la température en tête de colonne (palier vers \SI{65}{\celsius}), on met en évidence la présence de méthanol.

\textbf{5. Message de sensibilisation.}
\begin{quote}
\emph{« Attention aux alcools artisanaux frelatés ! Ils peuvent contenir du méthanol, un poison invisible : quelques gouttes suffisent à rendre aveugle, une petite quantité peut tuer. Ni la couleur, ni le goût, ni l'odeur ne permettent de le détecter. Ne consommez jamais d'alcool de source inconnue et signalez les cas suspects aux autorités sanitaires. Votre vie et votre vue n'ont pas de prix. »}
\end{quote}

\textbf{Points de vigilance :} le méthanol, bien que primaire, ne possède qu'un seul carbone ; les tests d'oxydation identifient une \emph{fonction}, jamais une \emph{molécule précise} : il faut toujours croiser avec une propriété physique ou une analyse complémentaire.
\end{corrige}

\newpage

\coursec{Fiche bilan}

\begin{popfigure}
\begin{center}
\begin{tikzpicture}[
  mindmap,
  grow cyclic,
  every node/.style={concept, text=white, font=\small\bfseries},
  root concept/.append style={concept color=popdark, font=\large\bfseries},
  level 1/.append style={level distance=3.6cm, sibling angle=72},
  level 2/.append style={level distance=2.1cm, sibling angle=45, font=\scriptsize\bfseries},
  concept connection/.append style={color=popmuted}
]
\node[root concept]{Composés oxygénés}
  child[concept color=popblue]{ node {Alcools $R-OH$}
    child { node[concept color=popblue!70] {3 classes : primaire, secondaire, tertiaire} }
    child { node[concept color=popblue!70] {Suffixe \emph{-ol}} }
  }
  child[concept color=poporange]{ node {Aldéhydes $R-CHO$}
    child { node[concept color=poporange!70] {Suffixe \emph{-al}} }
    child { node[concept color=poporange!70] {Schiff $+$, Fehling $+$, Tollens $+$} }
  }
  child[concept color=poppurple]{ node {Cétones $R-CO-R'$}
    child { node[concept color=poppurple!70] {Suffixe \emph{-one}} }
    child { node[concept color=poppurple!70] {DNPH $+$ seulement} }
  }
  child[concept color=popgreen]{ node {Propriétés physiques}
    child { node[concept color=popgreen!70] {Liaisons H : $\theta_{\text{éb}}$ élevée, solubilité} }
    child { node[concept color=popgreen!70] {Polyalcools : glycol, glycérol} }
  }
  child[concept color=poppink]{ node {Tests d'identification}
    child { node[concept color=poppink!70] {2,4-DNPH : précipité jaune} }
    child { node[concept color=poppink!70] {Fehling, Tollens, Schiff} }
  };
\end{tikzpicture}
\end{center}
\legende{Carte mentale de la leçon : les cinq idées-force à retenir sur les composés oxygénés.}
\end{popfigure}

\begin{bilan}[L'essentiel de la leçon]
\textbf{Définitions clés}
\begin{itemize}
  \item \cle{Composé oxygéné} : composé organique dont la molécule contient au moins un atome d'oxygène.
  \item \cle{Alcool} : groupe \cle{hydroxyle} $-\text{OH}$ porté par un carbone \emph{tétragonal} (hybridé $sp^3$) ; formule générale $\mathrm{R{-}OH}$, brute $C_nH_{2n+2}O$.
  \item \cle{Aldéhyde} : groupe \cle{carbonyle} \ce{C=O} en \emph{extrémité} de chaîne : $\mathrm{R{-}CHO}$, brute $C_nH_{2n}O$.
  \item \cle{Cétone} : groupe carbonyle \emph{à l'intérieur} de la chaîne : $\mathrm{R{-}CO{-}R'}$, brute $C_nH_{2n}O$.
  \item Aldéhydes et cétones de même formule brute sont des \cle{isomères de fonction}.
\end{itemize}

\textbf{Classes d'alcools (selon le nombre de groupes alkyle liés au carbone fonctionnel)}
\begin{center}
\begin{tabularx}{\linewidth}{|l|X|X|}
\hline
\rowcolor{popblueL} \textbf{Classe} & \textbf{Carbone fonctionnel lié à...} & \textbf{Exemple} \\
\hline
Primaire (I) & 1 groupe alkyle (ou 0) & \ce{CH3-CH2-OH} éthanol \\
\hline
Secondaire (II) & 2 groupes alkyles & \ce{CH3-CH(OH)-CH3} propan-2-ol \\
\hline
Tertiaire (III) & 3 groupes alkyles & \ce{(CH3)3C-OH} 2-méthylpropan-2-ol \\
\hline
\end{tabularx}
\end{center}

\textbf{Nomenclature express}
\begin{itemize}
  \item Alcool : chaîne la plus longue contenant $-\text{OH}$, indice le plus petit, suffixe \emph{-ol} : \ce{CH3-CH2-CH2-OH} = propan-1-ol.
  \item Aldéhyde : suffixe \emph{-al}, le carbone du $-\text{CHO}$ est toujours C1 : \ce{CH3-CHO} = éthanal.
  \item Cétone : suffixe \emph{-one} avec indice de position : \ce{CH3-CO-CH3} = propanone.
\end{itemize}

\textbf{Propriétés physiques}
\begin{itemize}
  \item Les liaisons hydrogène entre molécules d'alcool $\Rightarrow$ $\theta_{\text{éb}}$ des alcools \emph{supérieure} à celle des alcanes de masse voisine.
  \item Solubilité dans l'eau : grande pour $n \leq 4$ carbones, puis décroît quand la chaîne carbonée s'allonge.
  \item \cle{Polyalcools} : glycol \ce{HO-CH2-CH2-OH} (antigel), glycérol \ce{HO-CH2-CH(OH)-CH2-OH} (savonnerie, pharmacie).
\end{itemize}

\textbf{Tests d'identification (à connaître par c\oe ur)}
\begin{center}
\begin{tabularx}{\linewidth}{|l|X|X|X|}
\hline
\rowcolor{poppinkL} \textbf{Réactif} & \textbf{Aldéhyde} & \textbf{Cétone} & \textbf{Observation positive} \\
\hline
2,4-DNPH & $+$ & $+$ & précipité jaune-orangé \\
\hline
Réactif de Schiff & $+$ & $-$ & coloration rose-violette \\
\hline
Liqueur de Fehling & $+$ & $-$ & précipité rouge brique \ce{Cu2O} \\
\hline
Réactif de Tollens & $+$ & $-$ & miroir d'argent \\
\hline
\end{tabularx}
\end{center}

\textbf{Équations-types des tests}
\[ \ce{R-CHO + 2Cu^2+ + 5OH- -> R-COO- + Cu2O(s) + 3H2O} \quad \text{(Fehling)} \]
\[ \ce{R-CHO + 2[Ag(NH3)2]+ + 3OH- -> R-COO- + 2Ag(s) + 4NH3 + 2H2O} \quad \text{(Tollens)} \]

\textbf{Méthode express pour identifier un composé oxygéné}
\begin{enumerate}
  \item DNPH négative $\Rightarrow$ ce n'est ni un aldéhyde ni une cétone (penser à un alcool).
  \item DNPH positive puis Schiff (ou Fehling) positive $\Rightarrow$ \cle{aldéhyde}.
  \item DNPH positive mais Schiff et Fehling négatives $\Rightarrow$ \cle{cétone}.
\end{enumerate}
\end{bilan}

\begin{aretenir}[Checklist avant le Probatoire]
\begin{itemize}
  \item[$\square$] Je sais reconnaître les groupes hydroxyle $-\text{OH}$ et carbonyle $\ce{C=O}$ dans une formule développée.
  \item[$\square$] Je connais les formules brutes générales : $C_nH_{2n+2}O$ (alcool), $C_nH_{2n}O$ (aldéhyde et cétone).
  \item[$\square$] Je sais nommer un alcool (suffixe \emph{-ol}), un aldéhyde (\emph{-al}), une cétone (\emph{-one}) avec les bons indices.
  \item[$\square$] Je distingue sans hésiter alcool primaire, secondaire et tertiaire à partir de la formule développée.
  \item[$\square$] Je sais expliquer l'ébullition élevée et la solubilité des alcools par les liaisons hydrogène.
  \item[$\square$] Je connais les formules du glycol et du glycérol et une application de chacun.
  \item[$\square$] Je sais que la 2,4-DNPH caractérise le groupe carbonyle (aldéhydes \emph{et} cétones).
  \item[$\square$] Je sais que Schiff, Fehling et Tollens distinguent les aldéhydes (positifs) des cétones (négatifs).
  \item[$\square$] Je sais décrire les observations : précipité jaune (DNPH), rose (Schiff), rouge brique (Fehling), miroir d'argent (Tollens).
  \item[$\square$] Je sais écrire et équilibrer les équations-bilans des tests de Fehling et de Tollens.
  \item[$\square$] Je sais construire un raisonnement d'identification à partir des résultats de plusieurs tests.
  \item[$\square$] Je cite des exemples du quotidien camerounais : éthanol des boissons fermentées, glycérol des savonneries artisanales.
\end{itemize}
\end{aretenir}

\findecours

\end{document}
